% AMBAC 技術報告書（LuaLaTeX + jlreq）
% ビルド: report/ で latexmk -lualatex main.tex（技術報告書）
%        report/ で latexmk -lualatex main_detailed.tex（詳細版: \DETAILED を定義して本ファイルを読む）
\newif\ifdetailed
\ifdefined\DETAILED \detailedtrue \fi
\documentclass[paper=a4paper,fontsize=10.5pt,report,openany]{jlreq}
\usepackage{amsmath,amssymb,amsthm,bm,mathtools}
\usepackage{graphicx,xcolor}
\usepackage{longtable,booktabs,array,calc,multirow}
\usepackage{tikz}
\usetikzlibrary{arrows.meta,calc,decorations.pathmorphing,decorations.markings,patterns,angles,quotes,3d,positioning}
\usepackage{tikz-3dplot}
\usepackage{fancyvrb}
\usepackage{algorithm,algpseudocode}
\usepackage[most]{tcolorbox}
\usepackage[hidelinks,bookmarksnumbered]{hyperref}
\usepackage{url}
\graphicspath{{../figures/}}
\providecommand{\tightlist}{\setlength{\itemsep}{0pt}\setlength{\parskip}{0pt}}
\setcounter{tocdepth}{1}
\newcommand{\repo}{https://github.com/okamasaya/simulation_AMBAC}
\newcommand{\Sim}[1]{Sim#1}

% 図: \FIG{ファイル名} → 図環境（キャプションは \cap@ファイル名）
\makeatletter
\newcommand{\defcap}[2]{\expandafter\def\csname cap@#1\endcsname{#2}}
\newcommand{\FIG}[1]{%
  \begin{figure}[htbp]\centering
  \includegraphics[width=\linewidth,height=0.42\textheight,keepaspectratio]{#1}%
  \caption{\csname cap@#1\endcsname}\label{fig:#1}
  \end{figure}}
\makeatother
\defcap{sim01_timeseries.png}{腕 1 本の往復運動（下 → 上 → 下）における腕の角度と胴体の姿勢角（Sim01）。腕を戻すと姿勢も戻る。}
\defcap{sim01_sweep.png}{腕の質量比・長さ比と、180° の腕振りによる胴体の回転角（Sim01）。}
\defcap{sim02_cyclic.png}{左：曲率 $F(q)$ のマップと腕＋脚の長方形ループ。右：5 周分の胴体の姿勢角（Sim02）。}
\defcap{sim03_connection.png}{左：関節可動域内の曲率と最適円ループ。右：腕・脚・リアクションホイールの接続 $\Gamma$（Sim03）。}
\defcap{sim03_rest2rest.png}{90° の静止 → 静止旋回。左：リアクションホイール、右：AMBAC（Sim03）。}
\defcap{sim03_disturbance.png}{一定外乱トルクの下での姿勢保持と飽和（Sim03）。}
\defcap{sim04_human_cycle.png}{人間大モデルの順動力学（0.5 Hz の最適ループ）：姿勢、関節トルク、パワーと運動エネルギー（Sim04）。}
\defcap{sim04_scaling.png}{相似拡大による必要トルク・90° 旋回時間・角運動量容量の変化（Sim04）。}
\defcap{sim05a_learned_ambac.png}{学習した関節軌道と曲率マップ、姿勢の時系列、学習曲線（Sim05A）。}
\defcap{sim05b_pareto.png}{18 m 級の 90° 旋回における時間と推進剤のパレート曲線（左）と最適配分の例（右）（Sim05B）。}
\defcap{sim05b_robustness.png}{モデル誤差・外乱のあるプラントでの制御器の比較例（Sim05B）。}
\defcap{sim06_tipover.png}{荷物運び動作の ZMP（左上）、フルード則（右上）、軌道最適化（下）（Sim06）。}
\defcap{sim06_yaw_ratchet.png}{腕の水平振りによる台車のヨー（スティックスリップ）（Sim06）。}
\defcap{sim07_space.png}{宇宙での 3 次元 AMBAC：円錐運動による回転、回転軸、球面座標の特異点（Sim07）。}
\defcap{sim07_ground.png}{地上の 3 次元：定格荷重マップ、旋回中の ZMP、旋回中の腕の向き（Sim07）。}
\defcap{sim08_shoulder.png}{コッドマンのパラドックス（左 2 枚）と腕組みの姿勢・必要な内旋（右 2 枚）（Sim08）。}
\defcap{sim09_legs.png}{宇宙での脚の AMBAC 効果（左）、地上での戦車型と二本脚（中）、腰の戦略（右）（Sim09）。}
\defcap{sim10_steps.png}{段差の乗り越え：戦車型（左）、二本脚（中）、相似拡大（右）（Sim10）。}


\theoremstyle{definition}
\newtheorem{definition}{定義}[chapter]
\newtheorem{proposition}{命題}[chapter]
\newtheorem{example}{例題}[chapter]
\newtheorem{exercise}{考察課題}[chapter]
\renewcommand{\algorithmicrequire}{\textbf{入力:}}
\renewcommand{\algorithmicensure}{\textbf{出力:}}
\floatname{algorithm}{アルゴリズム}
% 詳細版だけに入れる内容
\newcommand{\Detail}[1]{\ifdetailed\input{detail/#1}\fi}
\newtcolorbox{keypoint}{colback=gray!6,colframe=gray!60,boxrule=0.5pt,arc=2pt,left=4pt,right=4pt,top=3pt,bottom=3pt}

\title{\textbf{AMBAC から考える\\二足歩行・形状・フィジカル AI}\\[0.5em]
\large ガンダムから生まれた空想科学は、どこまで真実に迫れるか\\
--- 第 1 段階 \ifdetailed 技術報告書（詳細版）\else 技術報告書\fi ---}
\author{simulation\_AMBAC プロジェクト}
\date{2026 年 11 月}

\begin{document}
\maketitle

\begin{abstract}
フィジカル AI（実世界で身体を持って動く AI）のハードウェアは高価で、試行錯誤のコストが大きい。
本報告書は、『機動戦士ガンダム』の SF 考証から生まれた AMBAC（Active Mass Balance Auto Control：手足の質量移動の反作用による姿勢制御）\cite{R1,R42} を出発点に、
安価なコンピュータシミュレーションで、ロボットの形状が姿勢制御・安定性・移動能力に与える意味をどこまで明らかにできるかを検討する。
中心となる問いは「二足歩行は本当にバランスが良いのか、なぜ人は二本脚なのか」である。

剛体力学・幾何力学・制御理論・機械学習を用いた 10 本のシミュレーション（Sim01〜Sim10）により、次のことを示した。
(1) AMBAC は角運動量保存の下で物理的に成立し、手足のループ運動による幾何学的位相として定量的に記述できる。リアクションホイールはその特殊な場合である。
(2) 二足歩行は静的にはバランスが悪いが、動かせる重い上体による能動的なバランス（地上版の AMBAC）と、人間の大きさでの段差への強さを持つ。
形状の優位は、宇宙・平地・段差、人間大・巨大機といった条件で入れ替わり、大型化では二乗三乗則が繰り返し効く。
(3) 人の肩のコッドマンのパラドックスは AMBAC と同じ幾何学的位相であり、腕組みのような人間らしい姿勢には、ひねりを含む 3 自由度の肩と肩甲骨の動きが必要になる。
(4) 学習は AMBAC を発見できるが、物理の構造を方策に組み込むことが成否を分ける。

以上から、用途や目的に応じて形状を選ぶ必要があること、フィジカル AI に必要なスペックとトレードオフを整理した。
最後に、組み合わせた形状について校正値を定義し、要求した動きや性能の実現性を検証する第 2 段階の構想を示す。

\medskip\noindent
\textbf{ソースコード・データ}：\url{\repo}

\end{abstract}

\tableofcontents
\Detail{notation}

\chapter{序章：二足歩行は本当にバランスが良いのか}\label{ch:intro}

\section{背景}

フィジカル AI は、ロボットのように実世界で身体を持って動く AI である。
その開発には高価なハードウェアが必要で、試行錯誤のたびに機体を壊すリスクや時間のコストがかかる。
そこで、安価なコンピュータシミュレーションで、どこまで実機の設計・制御・学習に使える知見が得られるかが重要になる。

本報告書が出発点とするのは、フィクションから生まれた概念である。
AMBAC（Active Mass Balance Auto Control、能動的質量移動による自動姿勢制御）は、
『機動戦士ガンダム』\cite{R1} の作品世界を支える SF 考証として、1981 年のムック本『ガンダムセンチュリー』で示されたとされる \cite{R42}。
人型の機動兵器が宇宙空間で手足を振り、その反作用で姿勢を変えるという考え方である。

これは単なる空想ではない。宇宙飛行士が手足の動きで体の向きを変えられることは古くから研究されており \cite{R8}、
落下する猫が空中で向きを変える現象も同じ原理で説明される \cite{R4,R7}。
実際に、宇宙飛行士の大西卓哉さんは国際宇宙ステーションで脚を使った姿勢制御に AMBAC の名を挙げ、ガンダムファンの間で話題になった \cite{R43,R44}。

\section{問い}

ロボットを題材にしたアニメやゲームでは、人型の機体が当たり前のように二本脚で立ち、手足を振って姿勢を整える。
一方で同じ作品世界には、脚のない機体について整備兵が言う「あんなの飾りです」という台詞もある \cite{R2}。
本報告書の中心となる問いは次のとおりである。

\begin{quote}
\textbf{二足歩行は本当にバランスが良いのか？　なぜ人は二本脚なのか？}
\end{quote}

\section{本報告書の答え（要約）}

\begin{itemize}
\item \textbf{静的には、二足歩行のバランスは悪い。}
  足裏の支持範囲が狭く重心が高いので、10 kg の荷物を腕の先で持つだけで、静止していても転倒する（第~\ref{ch:legs}章）。
  平地での安定性は戦車型の方が上である。
\item \textbf{それでも立てるのは、動かせる重い上体による能動的なバランスがあるから。}
  上体を前後に傾ける腰の戦略 \cite{R31} は、手足の代わりに上体を動かす地上版の AMBAC であり、これで同じ荷物を運べるようになる。
  段差を上るときも、上体を前に傾けることで、膝だけでは 18 cm だった限界が 60 cm まで広がる（第~\ref{ch:steps}章）。
\item \textbf{人間の大きさだからこそ、脚が活きる。}
  人間大の二本脚は戦車型の 2 倍以上の段差を越えるが、大型化すると二乗三乗則 \cite{R33} で逆転し、18 m 級の脚は 5 cm の段差しか越えられない。
\item \textbf{条件次第で優位は入れ替わる。}
  宇宙・平地・段差、人間大・巨大機で、どの形状が有利かが変わる。
  したがって、用途や目的に応じて形状を合わせる必要がある（パーツを組み合わせて機体を構成するゲームの「アセンブル」\cite{R3} のように）。
\end{itemize}

\section{本報告書の構成}

本報告書は、プロジェクトの第 1 段階「形状と、それが与える意味」の成果をまとめたものである。

\begin{description}
\item[第 I 部：AMBAC の物理（宇宙）] 角運動量保存から AMBAC の基本式を導き、幾何学的位相として理解する（第~\ref{ch:theory}章、第~\ref{ch:planar}章）。
  リアクションホイールとの関係（第~\ref{ch:rw}章）、トルクとスケーリング則（第~\ref{ch:dyn}章）、学習と最適制御（第~\ref{ch:ai}章）を扱う。
\item[第 II 部：下半身の形状（地上）] 腕の反作用による転倒・滑り（第~\ref{ch:tank}章）、「足は飾りか」（第~\ref{ch:legs}章）、段差（第~\ref{ch:steps}章）を扱う。
\item[第 III 部：上半身（肩と肩甲骨）] 3 次元の肩と回転の非可換性（第~\ref{ch:3d}章）、コッドマンのパラドックスと腕組み（第~\ref{ch:shoulder}章）を扱い、
  フィジカル AI に必要な可動域などのスペックとトレードオフを考察する。
\item[第 IV 部：結論] 問いへの答え、スペックとトレードオフの整理、第 2 段階の構想（第~\ref{ch:conclusion}章）。
\item[付録] 式の導出（付録~\ref{app:deriv}）、数値計算の方法と検証（付録~\ref{app:num}）、パラメータと再現手順（付録~\ref{app:param}）。
\end{description}

各章の見出しにある Sim01〜Sim10 は、リポジトリ \url{\repo} の \texttt{sims/} にあるシミュレーションスクリプトに対応する。
すべての数値と図は、付録~\ref{app:param} の手順で再現できる。

\section{前提とスコープ}

\begin{itemize}
\item \textbf{剛体の仮定}：機体は剛体リンクの集まりとして扱う。柔軟体・流体・ケーブル張力は扱わない。
\item \textbf{理想的なアクチュエータ}：多くの評価で、関節のトルク・速度を理想的に出せると仮定する。
  油圧・電動の動特性は第 2 段階の課題とする。
\item \textbf{環境}：宇宙（自由浮遊、角運動量が保存される）と地上（接触力・摩擦が働く）。重力は地球・火星・月を扱う。
\item \textbf{モデルの次元}：主に平面（2 次元）モデルで原理を確かめ、肩の運動など本質的に 3 次元の現象は 3 次元モデルで扱う。
\item \textbf{準静的な評価}：転倒・滑り・段差は、主にゆっくり動く場合の限界を求める。勢いを使う動的な動作は扱わない。
\end{itemize}

\Detail{intro_detail}

\ifdetailed
\part{基礎編}
\chapter{剛体力学の基礎}\label{ch:basics-rigid}

本章では、本報告書全体で使う剛体力学の道具を、必要な範囲に絞ってまとめる。
すでに力学に慣れている読者は飛ばしてよい。

\section{ベクトル・座標系・回転}

\subsection{慣性系と胴体座標系}

運動を記述する基準となる座標系として、空間に固定された\textbf{慣性系}と、機体の胴体に固定された\textbf{胴体座標系}の 2 つを使う。
同じベクトルでも、どちらの座標系で成分を書くかで数値が変わる。本報告書では、関節の位置や手足の重心位置は胴体座標系で与え、
胴体の向き（姿勢）は慣性系に対する回転として表す。

\subsection{平面の回転}

平面では、向きは 1 つの角度 $\theta$ で表せる。角度 $\theta$ の回転行列は
\begin{equation}
R(\theta) = \begin{pmatrix}\cos\theta & -\sin\theta\\ \sin\theta & \cos\theta\end{pmatrix}
\end{equation}
であり、次の性質を持つ。
\begin{itemize}
\item $R(\theta_1)R(\theta_2) = R(\theta_1+\theta_2)$：回転の合成は角度の足し算で、順序によらない（\textbf{可換}）。
\item $R(\theta)^\top = R(\theta)^{-1} = R(-\theta)$。
\item $\dfrac{d}{d\theta}R(\theta) = J R(\theta)$、ただし $J = R(\pi/2) = \begin{psmallmatrix}0&-1\\1&0\end{psmallmatrix}$。
\end{itemize}
平面のベクトル $\bm u,\bm v$ の外積はスカラー $\bm u\times\bm v = u_xv_y - u_yv_x$ とする。$\bm u\times(J\bm v) = \bm u\cdot\bm v$、
$(R\bm u)\times(R\bm v) = \bm u\times\bm v$（回転で不変）が成り立つ。

\subsection{3 次元の回転と SO(3)}

3 次元の向きは、$3\times3$ の回転行列 $R$（$R^\top R = E$、$\det R = 1$）で表す。このような行列全体を $\mathrm{SO}(3)$ と書く。
平面と違い、3 次元の回転の合成は\textbf{非可換}である：一般に $R_1R_2 \neq R_2R_1$。

\begin{example}[回転の順序]
本を机に置き、(a) 縦軸まわりに 90° 回してから横軸まわりに 90° 回す場合と、
(b) 横軸まわりに 90° 回してから縦軸まわりに 90° 回す場合を比べると、最終的な本の向きは異なる。
第~\ref{ch:3d}章では、腕の円錐運動による「ヨー → ピッチ」と「ピッチ → ヨー」で、最終姿勢が 29° 違うことを確かめた。
\end{example}

回転を数値で表す方法には次のものがある \cite{R18}。
\begin{longtable}{p{0.20\linewidth}p{0.36\linewidth}p{0.36\linewidth}}
\toprule
表現 & 長所 & 短所 \\ \midrule \endhead
回転行列（9 成分） & 合成が行列積で簡単、特異点がない & 冗長（拘束 6 個）、数値誤差で直交性が崩れる \\
オイラー角（3 成分） & 直観的、人が読みやすい & 特異点（ジンバルロック）がある \\
回転ベクトル（3 成分） & 「軸 × 角度」で直観的 & 180° 付近で不連続 \\
四元数（4 成分） & 特異点がない、補間が滑らか & 直観的でない、$q$ と $-q$ が同じ回転 \\
\bottomrule
\end{longtable}
本報告書の 3 次元計算では、姿勢の積分に回転ベクトル（指数写像）と回転行列を使い、オイラー角は使わない。

\subsection{角速度}

姿勢 $R(t)$ が時間とともに変わるとき、$R^\top\dot R$ は歪対称行列になり、ベクトル $\bm\omega$ を使って
\begin{equation}
R^\top\dot R = [\bm\omega]_\times = \begin{pmatrix}0&-\omega_3&\omega_2\\ \omega_3&0&-\omega_1\\ -\omega_2&\omega_1&0\end{pmatrix}
\end{equation}
と書ける。この $\bm\omega$ を\textbf{胴体座標で表した角速度}という（$\dot RR^\top$ から得られる角速度は慣性系の成分で、両者は $\bm\omega_{\text{慣性}} = R\bm\omega$ の関係にある）。
本報告書の 3 次元モデルでは、胴体座標の角速度 $\bm\omega$ を求めて $\dot R = R[\bm\omega]_\times$ を積分する。
短い時間 $\Delta t$ の間 $\bm\omega$ が一定なら、$R(t+\Delta t) = R(t)\exp([\bm\omega\Delta t]_\times)$ である（指数写像、ロドリゲスの公式で計算できる）。

\section{質点系の運動量と角運動量}

\subsection{重心と運動量}

質量 $m_i$ の質点が位置 $\bm r_i$ にあるとき、系の重心と全運動量は
\begin{equation}
\bm c = \frac{\sum_i m_i\bm r_i}{M},\quad M=\sum_i m_i,\qquad \bm P = \sum_i m_i\dot{\bm r}_i = M\dot{\bm c}.
\end{equation}
外力の合力を $\bm F$ とすると $\dot{\bm P} = \bm F$。内力（質点どうしが及ぼし合う力）は作用・反作用の法則で打ち消し合うので、
\textbf{内力では重心の運動を変えられない}。宇宙空間で手足をどう動かしても、重心は等速直線運動を続ける。

\subsection{角運動量と重心まわりの分解}

点 $\bm o$ まわりの角運動量は $\bm L_o = \sum_i (\bm r_i-\bm o)\times m_i\dot{\bm r}_i$ である。
重心を基準にとると、次のように分解できる（ケーニッヒの定理）：
\begin{equation}
\bm L_o = (\bm c-\bm o)\times M\dot{\bm c} + \underbrace{\sum_i (\bm r_i-\bm c)\times m_i(\dot{\bm r}_i-\dot{\bm c})}_{\text{重心まわりの角運動量 }\bm L}.
\end{equation}
外力のモーメントを $\bm N$ とすると $\dot{\bm L} = \bm N$（重心まわり）。内力のモーメントは（中心力なら）打ち消し合うので、
\begin{keypoint}
\textbf{外力のモーメントがなければ、重心まわりの角運動量 $\bm L$ は保存される。内部の運動（手足の動き）では $\bm L$ を変えられない。}
\end{keypoint}
これが AMBAC の出発点である。AMBAC は $\bm L$ を変えるのではなく、$\bm L$ を一定に保ったまま、胴体と手足の間で回転を「やりとり」する。

\subsection{剛体の角運動量と慣性}

剛体の重心まわりの角運動量は、平面では $L = I\omega$、3 次元では $\bm L = I\bm\omega$（$I$ は慣性テンソル）である。
平面での代表的な慣性モーメントを挙げる。
\begin{longtable}{p{0.40\linewidth}p{0.50\linewidth}}
\toprule
物体 & 重心まわりの慣性モーメント \\ \midrule \endhead
長さ $\ell$ の一様な棒（棒に垂直な軸） & $m\ell^2/12$ \\
幅 $a$、高さ $b$ の一様な長方形（面に垂直な軸） & $m(a^2+b^2)/12$ \\
半径 $r$ の一様な円盤（面に垂直な軸） & $mr^2/2$ \\
\bottomrule
\end{longtable}
回転軸が重心から距離 $d$ だけ離れている場合は、\textbf{平行軸の定理} $I_{\text{軸}} = I_{\text{重心}} + md^2$ を使う。
3 次元の慣性テンソルについても同様に $I_{\text{点}} = I_{\text{重心}} + m(|\bm d|^2E - \bm d\bm d^\top)$ が成り立つ。

\begin{example}[本報告書の人間大モデルの胴体]
胴体 60 kg を幅 0.35 m・高さ 1.0 m の一様な長方形とみなすと、$I_b = 60(0.35^2+1.0^2)/12 = 5.61$ kg\,m$^2$。
腕（4 kg、0.7 m の棒）の重心まわりの慣性は $4\times0.7^2/12 = 0.163$ kg\,m$^2$ で、胴体の 3\% にすぎない。
ところが腕の重心は肩から 0.35 m、肩は胴体重心から 0.35 m 離れているので、平行軸の項が加わり、腕が胴体の回転に与える影響はずっと大きくなる。
\end{example}

\subsection{2 体の重心まわりの角運動量}

質量 $m_1, m_2$ の 2 つの剛体が相対位置 $\bm r = \bm c_2-\bm c_1$ にあるとき、系全体の重心まわりの角運動量は
\begin{equation}
L = I_1\omega_1 + I_2\omega_2 + \mu\,\bm r\times\dot{\bm r},\qquad \mu = \frac{m_1m_2}{m_1+m_2}
\end{equation}
と書ける（$\mu$ は換算質量）。付録~\ref{app:twobody} では、これを使って胴体＋手足 1 本の $A(q)$, $B(q)$ を閉じた形で求めた。

\section{運動方程式}

\subsection{ニュートン・オイラー方程式}

各剛体について、重心の並進運動 $m\ddot{\bm c} = \bm F$ と、重心まわりの回転運動
\begin{equation}
I\dot{\bm\omega} + \bm\omega\times I\bm\omega = \bm N\quad(\text{胴体座標、オイラーの運動方程式})
\end{equation}
を連立させる方法をニュートン・オイラー法という。ZMP の式（付録~\ref{app:zmp}）は、地面に置かれた系全体にこれを適用したものである。

\subsection{ラグランジュ方程式}

一般化座標 $\bm x$（ここでは $(\theta, q)$）と運動エネルギー $T(\bm x,\dot{\bm x})$ を使うと、外力のない系の運動方程式は
\begin{equation}
\frac{d}{dt}\frac{\partial T}{\partial\dot x_k} - \frac{\partial T}{\partial x_k} = Q_k
\end{equation}
で与えられる（$Q_k$ は一般化力、関節トルクなど）。本報告書の系では、$T = \tfrac12\bm u^\top M(q)\bm u$（$\bm u = (\dot\theta,\dot q)$）であり、
$M$ が $\theta$ に依存しない。このような座標を\textbf{循環座標}という。

\subsection{循環座標とネーターの定理}

$\partial T/\partial\theta = 0$ かつ $\theta$ に対応する一般化力がゼロなら、ラグランジュ方程式から
\begin{equation}
\frac{d}{dt}\frac{\partial T}{\partial\dot\theta} = 0\quad\Longrightarrow\quad p_\theta = \frac{\partial T}{\partial\dot\theta} = A(q)\dot\theta + B(q)\cdot\dot q = L = \text{一定}.
\end{equation}
「系全体を回転させても運動エネルギーが変わらない」という対称性から、角運動量の保存が導かれる。これがネーターの定理の一例である。
保存量 $L$ を使って $\dot\theta$ を消去し、形状 $q$ だけの運動方程式に縮約する手続きを\textbf{ラウスの簡約}という \cite{R5}。

\section{例題：回転椅子の上で腕を振る}

\begin{example}
回転椅子に座った人（胴体と椅子をまとめて慣性モーメント $I_b = 3$ kg\,m$^2$、縦軸まわり）が、
水平に伸ばした腕（質量 4 kg、長さ 0.7 m、肩は回転軸から 0.2 m）を前から横へ 90° 振る。椅子の摩擦は無視する。

縦軸まわりの角運動量は保存されるので、腕の角度 $q$ について付録~\ref{app:twobody} と同じ形の式が成り立つ。
回転軸を胴体重心とみなし（椅子は固定軸なので、換算質量の代わりに腕の質量そのものを使う）、
$\bm a$ を肩の位置、$\ell_c=0.35$ m とすると
\[
A(q) = I_b + I_l + m_l\left(|\bm a|^2+\ell_c^2+2\ell_c|\bm a|\cos(q-\varphi_a)\right),\quad
B(q) = I_l + m_l\left(\ell_c^2+\ell_c|\bm a|\cos(q-\varphi_a)\right).
\]
$|\bm a| = 0.2$、肩が真横（$\varphi_a=90^\circ$）にあり、腕を前方（$q=0$）から横（$q=90^\circ$）へ振るとすると、
$\Gamma = B/A$ はおよそ 0.18〜0.21 の範囲で変わり、胴体は逆向きにおよそ $-\int_0^{\pi/2}\Gamma\,dq \approx -17^\circ$ 回る。
実際に回転椅子で試すと、腕を振った向きと逆に体が回るのを感じられる。
\end{example}

\section*{考察課題}
\begin{exercise}
胴体の慣性モーメントを 2 倍にすると、腕を 180° 振ったときの胴体の回転角はおよそ何倍になるか。式~\eqref{eq:AB-closed} を使って見積もれ。
\end{exercise}
\begin{exercise}
3 次元で、縦軸まわりに 90°、横軸まわりに 90° の回転を順序を変えて合成し、結果の回転行列が異なることを確かめよ。
\end{exercise}

\chapter{幾何力学入門：接続・曲率・ホロノミー}\label{ch:basics-geom}

AMBAC、落下する猫、宇宙飛行士の姿勢変更、コッドマンのパラドックスは、すべて同じ数学で説明できる。
本章では、その数学（幾何力学）の考え方を、式の厳密さよりも直観を優先して紹介する。詳しくは \cite{R5,R4,R6} を参照。

\section{形状と位置・姿勢を分ける}

自由に浮かぶ機体の状態は、次の 2 つに分けて考えると見通しがよい。
\begin{itemize}
\item \textbf{形状}（shape）：関節角 $q$。機体の「形」で、胴体から見た手足の配置。
\item \textbf{位置・姿勢}（group）：重心の位置と胴体の向き $\theta$（3 次元なら $R$）。形を保ったまま機体全体を動かしても変わらない量。
\end{itemize}
数学的には、全体の配位空間は「形状空間」の各点に「位置・姿勢の空間」が付いた構造（主ファイバー束）を持つ。
物理法則（運動エネルギー）が位置・姿勢の取り方によらない（対称性がある）ことが、以下の議論の前提になる。

\section{接続：形を変えたとき向きがどう変わるか}

角運動量がゼロのとき、形を少し変える（$dq$）と、向きは決まった量だけ変わる：
\begin{equation}
d\theta = -\Gamma(q)\cdot dq .
\end{equation}
この「形の変化から向きの変化を決める規則」$\Gamma$ を\textbf{接続}（connection）という。力学から決まる接続なので、機械的接続と呼ぶ。

\begin{keypoint}
接続は、形状空間の各点に「どちらに形を変えると、向きがどれだけ変わるか」を与える。時間や速さは含まれない。
したがって、形の変え方（経路）だけで向きの変化が決まる。
\end{keypoint}

\paragraph{例：手足 1 本の平面モデル}
形状空間は 1 次元（$q$ だけ）で、$\Gamma(q) = B(q)/A(q)$。付録~\ref{app:twobody} の閉じた形を使えば、腕の角度ごとの $\Gamma$ を手計算でも求められる。

\section{曲率とホロノミー}

\subsection{閉じた経路で何が起きるか}

形状を変えて元に戻す（形状空間の閉曲線 $C$ を 1 周する）と、向きの変化は
\begin{equation}
\Delta\theta = -\oint_C\Gamma\cdot dq .
\end{equation}
形状空間が 1 次元なら、閉曲線は「行って戻る」だけなので $\Delta\theta=0$ である。2 次元以上なら、閉曲線は面積を囲めるので、一般に $\Delta\theta\neq0$ となる。
この「1 周したときに残るずれ」を\textbf{ホロノミー}（holonomy）、あるいは\textbf{幾何学的位相}（geometric phase）という。

\subsection{曲率}

ストークスの定理により、閉曲線に沿った積分は、囲まれた領域での「曲率」の面積分に置き換えられる：
\begin{equation}
\Delta\theta = -\iint_S F\,dq_1\,dq_2,\qquad F = \frac{\partial\Gamma_2}{\partial q_1}-\frac{\partial\Gamma_1}{\partial q_2}.
\end{equation}
曲率 $F$ は「形状空間の各点で、小さなループを回したときにどれだけ向きが進むか」の密度である。
曲率マップ（第~\ref{ch:planar}章、第~\ref{ch:rw}章）は、この密度を色で表したもので、AMBAC の「効きやすい形の領域」を示す地図になる。

\begin{example}[小さな正方形ループ]
形状空間で一辺 $\varepsilon$ の小さな正方形を反時計回りに回すと、
4 辺の積分は、向かい合う辺どうしで $\Gamma$ の差を拾うので
\[
\Delta\theta \approx -\left[\Gamma_1(q_2)-\Gamma_1(q_2+\varepsilon)\right]\varepsilon - \left[\Gamma_2(q_1+\varepsilon)-\Gamma_2(q_1)\right]\varepsilon
\approx -F\,\varepsilon^2 .
\]
一辺を 2 倍にすると回転は 4 倍になる（2 次の効果）。Sim02 の曲率マップで $|F|\approx0.04$ の領域に一辺 0.5 rad（約 29°）の正方形を置くと、
1 周あたり $0.04\times0.25 = 0.01$ rad（約 0.6°）の回転が得られる。
\end{example}

\subsection{可換と非可換}

平面では向き $\theta$ が 1 つの数なので、ループを続けて回すと回転は足し算になり、順序によらない（可換）。
3 次元では向きが回転行列なので、ループごとの回転を合成すると順序によって結果が変わる（非可換）。
このとき曲率には、$\Gamma$ の微分に加えて、$\Gamma$ どうしの交換子の項が加わる。
非可換性は、第~\ref{ch:3d}章の「ヨー → ピッチ」と「ピッチ → ヨー」の違いとして現れた。

\section{非ホロノミック系とリー括弧}

\subsection{ホロノミックと非ホロノミック}

速度に関する拘束（ここでは $L=0$、すなわち $\dot\theta = -\Gamma\cdot\dot q$）が、位置に関する拘束（$\theta = f(q)$）に積分できるとき、ホロノミックという。
積分できないとき、非ホロノミックという。手足 1 本の AMBAC はホロノミック（$\theta$ は $q$ の関数）、手足 2 本以上で曲率がゼロでない AMBAC は非ホロノミックである。

非ホロノミック系の身近な例は、車の縦列駐車である。車は真横には動けない（速度の拘束）が、前後進とハンドル操作を組み合わせると、横にずれることができる。
AMBAC で「形を元に戻しながら向きだけを変える」のは、これと同じ構造である。

\subsection{リー括弧}

2 つの動き（ベクトル場）$X$, $Y$ を、$X$ を $\varepsilon$ → $Y$ を $\varepsilon$ → $X$ を $-\varepsilon$ → $Y$ を $-\varepsilon$ の順に行うと、
元の点には戻らず、$\varepsilon^2[X,Y]$ だけずれる。この $[X,Y]$ を\textbf{リー括弧}という。
リー括弧の方向にも動けることが、非ホロノミック系の可制御性の鍵である（Chow--Rashevskii の定理）。
AMBAC の曲率は、このリー括弧の向き成分にあたる。正弦波の入力で系を操る方法 \cite{R10} は、この考え方に基づく。

\subsection{Brockett の定理}

非ホロノミック系は、可制御であっても、滑らかで時間に依存しないフィードバック則では原点に漸近安定化できない \cite{R9}。
AMBAC で姿勢を目標に収束させるには、周期的な入力（ループ運動）、時間に依存する制御、あるいは切り替えが必要になる。
第~\ref{ch:ai}章で、構造のない学習方策が失敗し、振動子（周期的な構造）を持つ方策が成功したことは、この定理と整合している。

\section{同じ数学が現れる他の例}

\begin{longtable}{p{0.24\linewidth}p{0.30\linewidth}p{0.38\linewidth}}
\toprule
現象 & 形状空間 & 1 周で残るずれ \\ \midrule \endhead
AMBAC（本報告書） & 関節角 & 胴体の向き \\
落下する猫 \cite{R4,R7} & 上半身と下半身の曲げ・ひねり & 体の向き（足から着地できる） \\
宇宙飛行士の姿勢変更 \cite{R8} & 腕・脚の関節角 & 体の向き \\
コッドマンのパラドックス（第~\ref{ch:shoulder}章） & 腕の向き（単位球面） & 上腕のひねり（＝立体角） \\
低レイノルズ数の泳ぎ \cite{R6} & 微生物の体の形 & 位置（前進） \\
車の縦列駐車 & 前後進とハンドル角 & 横方向の位置 \\
\bottomrule
\end{longtable}

\section*{考察課題}
\begin{exercise}
形状空間で、曲率 $F$ が正の領域と負の領域にまたがる大きなループを回すと、1 周あたりの回転は小さくなる。その理由を説明せよ。
\end{exercise}
\begin{exercise}
腕と脚を同じ位相で同時に往復させた場合（形状空間で直線上を往復する場合）、正味の回転がゼロになることを示せ。
\end{exercise}

\chapter{最適化・制御・学習の基礎}\label{ch:basics-opt}

本報告書では、軌道の最適化（第~\ref{ch:tank}章、第~\ref{ch:3d}章、第~\ref{ch:legs}章）、
推進剤の配分（第~\ref{ch:ai}章）、方策の学習（第~\ref{ch:ai}章）に、最適化・制御・学習の手法を使った。本章ではその基礎をまとめる。

\section{最適化問題}

最適化問題は、変数 $\bm x$、目的関数 $f(\bm x)$、制約 $g_i(\bm x)\ge0$、$h_j(\bm x)=0$ から成る：
\begin{equation}
\min_{\bm x}\ f(\bm x)\quad\text{s.t.}\quad g_i(\bm x)\ge0,\quad h_j(\bm x)=0 .
\end{equation}

\subsection{線形計画（LP）}

目的関数と制約がすべて 1 次式のとき、線形計画という：
\begin{equation}
\min_{\bm x}\ \bm c^\top\bm x\quad\text{s.t.}\quad A_{ub}\bm x\le\bm b_{ub},\quad A_{eq}\bm x=\bm b_{eq},\quad \bm l\le\bm x\le\bm u .
\end{equation}
線形計画は凸最適化の一種で、局所解が大域解になり、大規模な問題でも単体法や内点法で確実に解ける。
本報告書では、推進剤最小の配分（付録~\ref{app:lp}）と、段差を這い上がるのに必要な摩擦係数（付録~\ref{app:steps}）を線形計画で解いた（ソルバは HiGHS \cite{R25}）。

\paragraph{絶対値の扱い}
推進剤は $|\tau|$ に比例するが、絶対値は 1 次式ではない。$\tau = \tau^+ - \tau^-$、$\tau^\pm\ge0$ と分けて $\tau^++\tau^-$ を最小化すると、
最適解では $\tau^+$ と $\tau^-$ の少なくとも一方が 0 になり、$\tau^++\tau^- = |\tau|$ となる。

\subsection{非線形計画と SQP}

目的関数や制約が非線形のとき、逐次二次計画法（SQP）がよく使われる。各反復で、目的関数を二次式、制約を 1 次式で近似した部分問題を解き、解を更新する。
本報告書の軌道最適化では、SQP の一種である SLSQP \cite{R26} を使い、動作時間 $T$ と軌道の補正係数を同時に最適化した。
非線形計画は局所解に陥りうるので、複数の初期値から解き、制約を満たす最良解を採用した。

\paragraph{最適性の条件（KKT 条件）}
制約付き最適化の解では、目的関数の勾配が、有効な（等号で効いている）制約の勾配の 1 次結合で書ける。
直観的には「これ以上どの方向に動いても、目的関数を減らすか制約を破るかのどちらか」という状態である。
第~\ref{ch:tank}章の最適軌道で ZMP が許容範囲の端に張り付いたのは、転倒の制約が有効になっていることを示す。

\section{最適制御とモデル予測制御}

\subsection{直接法}

時間とともに変わる入力 $u(t)$ を決める問題（最適制御）は、時間を区切って各区間の入力を変数とすれば、有限次元の最適化問題になる（直接法）。
第~\ref{ch:ai}章の推進剤配分は、平均化モデルが線形なので、直接法で線形計画になった。

\subsection{モデル予測制御（MPC）}

最適制御で一度だけ計画して実行する（開ループ）と、モデルの誤差や外乱で結果がずれる。
\textbf{モデル予測制御}は、一定時間ごとに現在の状態を測り、そこから残り時間（または一定の先読み時間）について最適制御問題を解き直し、
最初の部分だけを実行する。計画の最適性と、フィードバックによる誤差の修正を両立できる。

\begin{algorithm}[htbp]
\caption{モデル予測制御（第~\ref{ch:ai}章の推進剤配分）}
\begin{algorithmic}[1]
\Require 目標姿勢、旋回時間 $T$、再計画の間隔 $\Delta T$
\For{$t = 0, \Delta T, 2\Delta T, \dots < T$}
  \State 現在の状態（周回平均の姿勢 $\bar\theta$、角運動量 $L$、位相 $\varphi$）を測る
  \State 残り時間 $T-t$ について、推進剤最小の線形計画を解く
  \State 計画の最初の $\Delta T$ 秒分のスラスタトルクと周回周波数を実行する
\EndFor
\end{algorithmic}
\end{algorithm}

\section{強化学習と進化戦略}

\subsection{強化学習の枠組み}

強化学習では、エージェントが状態（観測）$\bm o$ を見て行動 $\bm a = \pi_{\bm w}(\bm o)$ を選び、環境から報酬を受け取る。
方策 $\pi_{\bm w}$ のパラメータ $\bm w$（ニューラルネットワークの重み）を、報酬の期待値 $J(\bm w)$ が大きくなるように調整する。
物理の式を明示的には使わず、試行錯誤から方策を得る点が、最適制御（モデルを使う）との違いである。

\subsection{進化戦略（ES）}

\begin{figure}[htbp]\centering
\resizebox{0.85\linewidth}{!}{% 進化戦略の 1 反復
\begin{tikzpicture}[>=Stealth,font=\footnotesize,box/.style={draw,rounded corners,fill=gray!8,align=center,minimum width=2.6cm,minimum height=1.0cm}]
\node[box] (a) at (0,0) {方策のパラメータ $\bm w$};
\node[box] (b) at (3.6,0) {ノイズを加えた\\$\bm w\pm\sigma\bm\varepsilon_k$};
\node[box] (c) at (7.2,0) {シミュレーションで\\一括評価（報酬）};
\node[box] (d) at (7.2,-1.9) {報酬を順位に変換};
\node[box] (e) at (3.6,-1.9) {勾配を推定\\$\frac{1}{n\sigma}\sum(r^+-r^-)\bm\varepsilon$};
\node[box] (f) at (0,-1.9) {Adam で更新};
\draw[->,thick] (a)--(b); \draw[->,thick] (b)--(c); \draw[->,thick] (c)--(d); \draw[->,thick] (d)--(e); \draw[->,thick] (e)--(f); \draw[->,thick] (f)--(a);
\end{tikzpicture}
}
\caption{進化戦略（OpenAI-ES）の 1 反復の流れ。}\label{fig:tikz-esloop}
\end{figure}

進化戦略は、方策のパラメータそのものにノイズを加えて試し、報酬の良かった方向へパラメータを動かす方法である \cite{R19}。
パラメータ $\bm w$ のまわりで、正規分布のノイズ $\bm\varepsilon_k$ を加えた $\bm w+\sigma\bm\varepsilon_k$ を $n$ 通り試すと、
報酬の勾配は
\begin{equation}
\nabla_{\bm w}J \approx \frac{1}{n\sigma}\sum_{k=1}^n F(\bm w+\sigma\bm\varepsilon_k)\,\bm\varepsilon_k
\end{equation}
で推定できる。本報告書では、次の工夫を加えた OpenAI-ES を使った。
\begin{itemize}
\item \textbf{対称サンプリング}：$+\bm\varepsilon_k$ と $-\bm\varepsilon_k$ を対で試し、推定のばらつきを減らす。
\item \textbf{順位変換}：報酬の値そのものではなく順位を使い、外れ値の影響を抑える。
\item \textbf{Adam} \cite{R20}：勾配の移動平均と大きさで歩幅を調整する。
\end{itemize}

\begin{algorithm}[htbp]
\caption{OpenAI-ES（本報告書の実装）}
\begin{algorithmic}[1]
\Require 初期パラメータ $\bm w$、ノイズの大きさ $\sigma$、対の数 $n$、学習率 $\eta$
\For{各反復}
  \State $\bm\varepsilon_1,\dots,\bm\varepsilon_n\sim\mathcal N(0,E)$ を生成する
  \State $2n$ 個の方策 $\bm w\pm\sigma\bm\varepsilon_k$ を、まとめてシミュレーションで評価し、報酬 $F_k^\pm$ を得る
  \State 全 $2n$ 個の報酬を順位に変換し、$[-0.5, 0.5]$ に正規化する（$r_k^\pm$）
  \State 勾配の推定 $\bm g = \frac{1}{n\sigma}\sum_k (r_k^+-r_k^-)\bm\varepsilon_k$
  \State Adam で $\bm w$ を更新する
\EndFor
\end{algorithmic}
\end{algorithm}

勾配を使わない（シミュレーションを微分しない）ので、報酬が最後にしか決まらない問題や、長い試行（数百ステップ）に強い。
一方、パラメータが多いと必要な試行数が増える。本報告書の方策は数百〜数千パラメータの小さなネットワークで、ES に向いた規模である。

\subsection{帰納バイアス：方策にどんな構造を持たせるか}

学習の成否は、方策にどんな構造（帰納バイアス）を持たせるかで大きく変わる。
\begin{itemize}
\item \textbf{中枢パターン生成器（CPG）} \cite{R22}：動物の歩行のリズムを生む神経回路にならい、周期的な信号を方策に与える。
\item \textbf{軌道生成器の調整（PMTG）} \cite{R21}：あらかじめ周期軌道を生成する仕組みを持たせ、学習ではその形や速さを調整する。
\end{itemize}
第~\ref{ch:ai}章では、周期信号を観測に加えるだけ（CPG 的）では不十分で、ループを描く振動子を持たせる（PMTG 的）と学習が成功した。
幾何学的位相が振幅の 2 次の効果であるため（第~\ref{ch:basics-geom}章）、最初から有限の大きさのループを描ける構造が必要だったと解釈できる。

\subsection{シミュレーションから実機へ}

シミュレーションで学習した方策を実機で使うと、モデルの誤差で性能が落ちる（sim-to-real の問題）。
対策の 1 つが\textbf{ドメインランダム化} \cite{R23} で、学習時にシミュレーションのパラメータ（質量、摩擦、外乱など）を乱数で変え、どれでもうまくいく方策を学ぶ。
第~\ref{ch:ai}章の学習方策は、ループの効き方 $\kappa$、慣性 $\alpha$、外乱 $\tau_d$ をランダムにして学習し、MPC より誤差・外乱に強くなった。
第 2 段階では、校正値（実機の測定値）でランダム化の範囲を決めることを想定している。

\section*{考察課題}
\begin{exercise}
推進剤最小の配分で、AMBAC の周回周波数 $f$ に小さな正の費用（$10^{-6}$）を付けた理由を、線形計画の解が一意でない場合を考えて説明せよ（第~\ref{ch:ai}章の考察 C-2 も参照）。
\end{exercise}
\begin{exercise}
進化戦略で、ノイズの大きさ $\sigma$ を大きくしすぎた場合と小さくしすぎた場合に、それぞれ何が起きるか考えよ。
\end{exercise}

\fi

\part{AMBAC の物理（宇宙）}
\chapter{理論的準備：角運動量保存と AMBAC}\label{ch:theory}

\section{AMBAC とは何か（物理的に言い直す）}

AMBAC（Active Mass Balance Auto Control）は、『機動戦士ガンダム』\cite{R1} の SF 考証として示された概念で \cite{R42}、 「四肢を動かした反作用で姿勢を制御する」と説明される。物理学の言葉にすると次の 1 文に尽きる。

\begin{quote}
\textbf{外トルクのない系で、内部の質量配置（形状）を変えることにより、全角運動量 L を保存したまま 胴体の向きを変える制御。}
\end{quote}

推進剤を消費しないこと、反力を外部に依存しないことが本質的な利点である。 一方で、\textbf{全角運動量そのものは 1 ミリも変えられない}（回転を止める・始めることは単独ではできない） という原理的制約がある。この「できること / できないこと」の境界を明確にするのが本リポジトリの第一目標。

\section{平面自由浮遊マルチボディ系}

胴体（姿勢角 θ）に n 本のリンク（関節角 q = (q\textsubscript{1},\ldots,q\textsubscript{n})）が付いた系を考える。 外力ゼロなので重心は等速直線運動、重心系で見れば静止とみなせる。 系重心まわりの全角運動量は

\[
L = \sum_i \left[ I_i \dot\phi_i + m_i (\mathbf r_i - \mathbf c) \times (\dot{\mathbf r}_i - \dot{\mathbf c}) \right]
\]

（φ\textsubscript{i}: 各剛体の絶対角、r\textsubscript{i}: 各剛体重心、c: 系重心）。 回転対称性から L は θ に依存せず、速度について線形なので

\[
L = A(q)\,\dot\theta + B(q)\cdot\dot q
\]

と書ける。A(q) \textgreater{} 0 は「形状 q で固まった系全体の慣性モーメント」、B(q) は関節運動と姿勢の結合項。

\section{AMBAC の基本式（L = 0 の場合）}

手足が 1 本の場合には $A(q)$, $B(q)$ を閉じた形で書ける（付録~\ref{app:twobody}）。基本式の導出と幾何学的な意味は付録~\ref{app:connection} にまとめた。


静止状態から始めれば L ≡ 0 なので

\[
\boxed{\ \dot\theta = -\,\Gamma(q)\cdot\dot q,\qquad \Gamma(q) = B(q)/A(q)\ }
\]

Γ は幾何力学で \textbf{機械的接続（mechanical connection）} と呼ばれる量。重要な性質:

\begin{enumerate}
\def\labelenumi{\arabic{enumi}.}
\tightlist
\item
  \textbf{時間に依存しない。} θ の変化量は関節空間の「経路」だけで決まり、速さには依存しない。 ゆっくり動かしても速く動かしても回転角は同じ（必要トルクとエネルギーが変わるだけ）。
\item
  \textbf{スケール不変。} Γ は質量比・長さ比のみの関数。人間大でも 18 m 級でも、 幾何学的に相似で質量分布が相似なら回転角は同じ（Sim01 で数値確認済み）。 大型化で効くのはトルク要求・構造強度・時間スケールの方。
\end{enumerate}

\section{可逆的 AMBAC と非ホロノミック AMBAC}

図~\ref{fig:tikz-phase} に幾何学的位相の考え方を示す。

\begin{figure}[htbp]\centering
\resizebox{0.85\linewidth}{!}{% 幾何学的位相：形状空間のループと正味の回転
\begin{tikzpicture}[>=Stealth,font=\small]
\begin{scope}
  \draw[->] (-0.2,0)--(3.6,0) node[right]{$q_1$（腕）};
  \draw[->] (0,-0.2)--(0,3.0) node[above]{$q_2$（脚）};
  \shade[inner color=blue!35,outer color=white] (1.8,1.4) ellipse (1.5 and 1.15);
  \draw[thick,red,postaction={decorate},decoration={markings,mark=at position 0.25 with {\arrow{Stealth}},mark=at position 0.75 with {\arrow{Stealth}}}] (1.8,1.4) ellipse (1.0 and 0.75);
  \node[red] at (2.95,2.15) {$C$};
  \node at (1.8,1.4) {$S$};
  \node[font=\footnotesize,align=center] at (1.8,-0.6) {形状空間：色は曲率 $F$ の大きさ};
\end{scope}
\draw[->,very thick,gray] (4.0,1.4)--(4.9,1.4);
\begin{scope}[shift={(6.6,1.4)}]
  \foreach \a/\c in {0/black!25,25/black!45,50/black!75}{
    \draw[thick,\c,rotate=\a,fill=white,fill opacity=0.5] (-0.18,-0.7) rectangle (0.18,0.7);}
  \draw[->,red,thick] (0,1.0) arc (90:150:1.0);
  \node[red,align=left,font=\footnotesize] at (0.8,1.25) {1 周ごとに\\$\Delta\theta=-\iint_S F$};
  \node[font=\footnotesize,align=center] at (0,-2.0) {形状が戻っても\\向きだけが進む};
\end{scope}
\end{tikzpicture}
}
\caption{幾何学的位相：形状空間（関節空間）の閉曲線 $C$ を 1 周すると、形状は元に戻るが、胴体の向きは $C$ が囲む曲率の量だけ進む。}\label{fig:tikz-phase}
\end{figure}

\textbf{関節 1 つ（n = 1）} の場合、θ = -∫Γ(q)dq は q だけの関数になり、 腕を戻すと姿勢も必ず元に戻る（\textbf{可逆的・ホロノミック}）。 「腕を振った分だけ一時的に向きを変える」ことしかできず、回転可能範囲も上限がある。

\textbf{関節 2 つ以上（n ≥ 2）} で、関節空間の閉ループ C を 1 周すると、形状は元に戻るのに

\[
\Delta\theta = -\oint_C \Gamma\cdot dq = -\iint_S F\,dq_1 dq_2,\qquad
F = \frac{\partial \Gamma_2}{\partial q_1} - \frac{\partial \Gamma_1}{\partial q_2}
\]

だけ姿勢が進む（ストークスの定理）。F は接続の \textbf{曲率}。これは \textbf{幾何学的位相（geometric phase）} の一例で、落下する猫が空中で向きを変えるのと同じ原理である。 ループを繰り返せば \textbf{任意の角度まで回転できる}。これが「AMBAC で姿勢を自在に変える」ことの数学的な正体。

設計上の含意: - 1 周あたりの回転量は「ループが囲む曲率の量」で決まる → \textbf{曲率の大きい形状領域でループを回すべき}。 - 曲率マップは機体設計（四肢の質量・長さ・取付位置）で決まる → \textbf{機体設計そのものが AMBAC 性能を決める}。

\section{L ≠ 0 の場合（回転を止められるか？）}

L ≠ 0 なら θ̇ = (L − B·q̇)/A(q)。内部運動では L は変えられないが、 \textbf{A(q) を大きくする（四肢を伸ばす）ことで回転速度を下げる}ことはできる（フィギュアスケートの逆）。 完全停止には外部トルク（スラスタ、リアクションホイール、環境力）が必要。 → 今後のシミュレーション課題。

\section{フィジカル AI との接点（予告）}

\begin{itemize}
\tightlist
\item
  姿勢制御が \textbf{非ホロノミック拘束下の経路計画問題} になる（Brockett 条件により、 滑らかな時不変フィードバックでは原点を漸近安定化できない → 周期入力・時変制御・学習が必要）。
\item
  動作が「経路のみ」に依存するため、\textbf{形状空間での計画}と\textbf{トルク・時間の配分}を分離できる。
\item
  曲率マップは学習した方策の解釈手段（なぜその動きを選んだか）としても使える。
\end{itemize}

\Detail{ch01}
\chapter{シミュレーション結果：平面モデル（Sim01, Sim02）}\label{ch:planar}

モデルは、人間を矢状面（体を左右に分ける面）に投影した粗い平面モデルである（胴体 60 kg, 0.35×1.0 m / 腕 4 kg, 0.7 m / 脚 12 kg, 0.9 m）。図~\ref{fig:tikz-planar} に構成を示す。

\begin{figure}[htbp]\centering
\resizebox{0.75\linewidth}{!}{% 平面モデル：胴体＋腕＋脚
\begin{tikzpicture}[>=Stealth,scale=1.15,font=\small]
\begin{scope}[rotate=12]
  \draw[thick,fill=gray!15] (-0.35,-1) rectangle (0.35,1);
  \fill (0,0) circle (1.5pt) node[left=6pt,font=\footnotesize]{胴体重心};
  \draw[dashed,gray] (0,-1.4)--(0,1.5);
  \coordinate (S) at (0,0.7); \coordinate (H) at (0,-0.3);
  \draw[line width=3pt,blue!60] (S) -- ++(35:1.4) coordinate (A);
  \draw[line width=4pt,red!60] (H) -- ++(-60:1.8) coordinate (L);
  \fill (S) circle (2pt); \fill (H) circle (2pt);
  \draw[->] (S)++(0.7,0) arc (0:35:0.7); \node at ($(S)+(10:0.95)$) {$q_1$};
  \draw[gray] (S)--++(1,0); \draw[gray] (H)--++(1,0);
  \draw[->] (H)++(0.8,0) arc (0:-60:0.8); \node at ($(H)+(-25:1.05)$) {$q_2$};
  \node[blue!70!black,above] at (A) {腕 $m_1,\ \ell_1$};
  \node[red!70!black,right] at (L) {脚 $m_2,\ \ell_2$};
  \node[right,font=\footnotesize] at (S) {肩};
  \node[left=4pt,font=\footnotesize,yshift=-6pt] at (H) {股関節};
\end{scope}
\draw[gray] (0,0)--(0,1.8);
\draw[->] (0,1.5) arc (90:102:1.5); \node[above] at (-0.15,1.55) {$\theta$};
\node[align=left,font=\footnotesize] at (4.2,1.3) {$L = A(q)\dot\theta + B(q)\cdot\dot q$\\$L=0 \Rightarrow \dot\theta = -\Gamma(q)\cdot\dot q$};
\end{tikzpicture}
}
\caption{平面モデル：胴体＋腕＋脚。胴体の姿勢角 $\theta$、関節角 $q_1$（腕）、$q_2$（脚）。}\label{fig:tikz-planar}
\end{figure}

\section{\texorpdfstring{Sim01: 腕 1 本の往復（\texttt{sims/sim01\_single\_arm.py}）}{Sim01: 腕 1 本の往復（sims/sim01\_single\_arm.py）}}

\FIG{sim01_timeseries.png}

\begin{longtable}[]{@{}
  >{\raggedright\arraybackslash}p{(\columnwidth - 2\tabcolsep) * \real{0.5000}}
  >{\raggedright\arraybackslash}p{(\columnwidth - 2\tabcolsep) * \real{0.5000}}@{}}
\toprule\noalign{}
\begin{minipage}[b]{\linewidth}\raggedright
項目
\end{minipage} & \begin{minipage}[b]{\linewidth}\raggedright
結果
\end{minipage} \\
\midrule\noalign{}
\endhead
\bottomrule\noalign{}
\endlastfoot
腕を下→上へ 180° 振ったときの胴体回転 & \textbf{−16.0°}（逆向き） \\
腕を戻した後の残留姿勢 & 約 1e-9°（可逆的） \\
角運動量の残差 & 約 1e-5 kg·m²/s（数値微分による誤差の範囲） \\
18 m 級（長さ ×10、質量 ×1000）での同条件 & \textbf{−16.0432°（人間大と同一）} \\
\end{longtable}

\FIG{sim01_sweep.png}

\textbf{考察} - 腕 1 本では「一時的な首振り」しかできない。姿勢を保ったまま腕を戻す手段がない。 - 回転量は腕の質量比 μ と長さ比にほぼ比例して増える。μ≈0.06（人間の腕）で約 16°。 AMBAC を主機能とするなら、四肢を \textbf{重く・長く} する設計圧力がかかる（推進剤タンクを四肢に置く等）。 - スケール不変性により「巨大ロボだから AMBAC が効く/効かない」ということは \textbf{回転角については無い}。

\section{\texorpdfstring{Sim02: 腕＋脚の周期運動（\texttt{sims/sim02\_cyclic\_gait.py}）}{Sim02: 腕＋脚の周期運動（sims/sim02\_cyclic\_gait.py）}}

\FIG{sim02_cyclic.png}

\begin{longtable}[]{@{}
  >{\raggedright\arraybackslash}p{(\columnwidth - 2\tabcolsep) * \real{0.5000}}
  >{\raggedright\arraybackslash}p{(\columnwidth - 2\tabcolsep) * \real{0.5000}}@{}}
\toprule\noalign{}
\begin{minipage}[b]{\linewidth}\raggedright
項目
\end{minipage} & \begin{minipage}[b]{\linewidth}\raggedright
結果
\end{minipage} \\
\midrule\noalign{}
\endhead
\bottomrule\noalign{}
\endlastfoot
1 周あたりの正味回転（数値積分） & \textbf{+2.9778°} \\
1 周あたりの正味回転（曲率の面積分） & \textbf{+2.9778°}（一致） \\
360° 回転に必要な周回数 & 約 121 周 \\
\end{longtable}

\textbf{考察} - 形状を周期的に戻しても姿勢が累積する ＝ 非ホロノミック AMBAC が成立。 - 今回のループは曲率の弱い領域に置かれている。曲率マップ中央の負のピーク （腕 ≈ −20°、脚 ≈ +50°）を囲めば 1 周あたりの回転量は増やせる → ループ最適化が次の課題。 - 1 周あたり約 3° は人間の宇宙飛行士の経験（ゆっくりとしか向きを変えられない）と定性的に整合。

\Detail{ch02}
\chapter{リアクションホイールと AMBAC の統一的理解（Sim03）}\label{ch:rw}

\section{同じ式、違う関節}

導出は付録~\ref{app:rw}（ホイールは曲率ゼロの関節）と付録~\ref{app:saturation}（角運動量の容量）を参照。

\begin{figure}[htbp]\centering
\resizebox{0.8\linewidth}{!}{% リアクションホイールと手足
\begin{tikzpicture}[>=Stealth,font=\small]
\begin{scope}
  \draw[thick,fill=gray!15] (-0.5,-1) rectangle (0.5,1);
  \draw[thick,fill=orange!30] (0,0) circle (0.35);
  \draw[->,thick] (0.5,0) arc (0:300:0.5);
  \fill (0,0) circle (1.5pt);
  \node[align=center,font=\footnotesize,below] at (0,-1.15) {リアクションホイール\\重心が回転軸上：形が変わらない\\$\Gamma=$ 一定、曲率 $F=0$\\何回転でも回せる};
\end{scope}
\begin{scope}[shift={(5,0)}]
  \draw[thick,fill=gray!15] (-0.5,-1) rectangle (0.5,1);
  \fill (0,0.6) circle (2pt);
  \draw[line width=3pt,blue!60] (0,0.6)--++(20:1.3);
  \draw[line width=3pt,blue!25] (0,0.6)--++(80:1.1);
  \draw[->,thick] (0,0.6)++(20:0.9) arc (20:80:0.9);
  \node[align=center,font=\footnotesize,below] at (0,-1.15) {手足（AMBAC）\\重心が動く：形が変わる\\$\Gamma(q)$ が形に依存、$F\neq0$\\可動域がある};
\end{scope}
\node[font=\footnotesize] at (2.5,0.4) {$\dot\theta=-\Gamma\,\dot q$};
\end{tikzpicture}
}
\caption{リアクションホイールと手足（AMBAC）。どちらも $\dot\theta=-\Gamma\dot q$ で書けるが、関節の性質が異なる。}\label{fig:tikz-rw}
\end{figure}


人工衛星のリアクションホイール（RW）と AMBAC は、どちらも \textbf{運動量交換} による姿勢制御で、同じ式で書ける。

\[
\dot\theta = \frac{L - B(q)\cdot\dot q}{A(q)} = \frac{L}{A(q)} - \Gamma(q)\cdot\dot q
\]

\begin{longtable}[]{@{}
  >{\raggedright\arraybackslash}p{(\columnwidth - 4\tabcolsep) * \real{0.3333}}
  >{\raggedright\arraybackslash}p{(\columnwidth - 4\tabcolsep) * \real{0.3333}}
  >{\raggedright\arraybackslash}p{(\columnwidth - 4\tabcolsep) * \real{0.3333}}@{}}
\toprule\noalign{}
\begin{minipage}[b]{\linewidth}\raggedright
\end{minipage} & \begin{minipage}[b]{\linewidth}\raggedright
リアクションホイール
\end{minipage} & \begin{minipage}[b]{\linewidth}\raggedright
AMBAC（四肢）
\end{minipage} \\
\midrule\noalign{}
\endhead
\bottomrule\noalign{}
\endlastfoot
関節 & 連続回転できる & 可動域がある \\
回したときの質量配置 & 変わらない（重心が回転軸上） & 変わる \\
接続 Γ & \textbf{一定}（I\_w / A） & \textbf{形状 q に依存} \\
曲率 F & 0 & ≠ 0 \\
任意角の回転 & 回し続ければ可能（ホロノミック） & ループ運動で可能（非ホロノミック・幾何学的位相） \\
たとえ & 低ギア・無限ストローク & 高ギア・短ストローク \\
\end{longtable}

\textbf{RW は「AMBAC の特殊な場合」}、すなわち曲率ゼロ・可動域無限の関節とみなせる。 逆に AMBAC は「可動域を補うために、形状空間のループで回転を稼ぐ RW」と言える。

\section{Part A: 接続とループ最適化}

\FIG{sim03_connection.png}

関節可動域は腕 −150°〜150°、脚 −120°〜30°。

\begin{longtable}[]{@{}
  >{\raggedright\arraybackslash}p{(\columnwidth - 2\tabcolsep) * \real{0.5000}}
  >{\raggedright\arraybackslash}p{(\columnwidth - 2\tabcolsep) * \real{0.5000}}@{}}
\toprule\noalign{}
\begin{minipage}[b]{\linewidth}\raggedright
項目
\end{minipage} & \begin{minipage}[b]{\linewidth}\raggedright
値
\end{minipage} \\
\midrule\noalign{}
\endhead
\bottomrule\noalign{}
\endlastfoot
RW の Γ\_w & 0.0014（ホイール 1 回転で胴体 0.51°） \\
脚の Γ\textsubscript{2} & 約 0.33（脚を 1° 振ると胴体 0.33°） \\
Sim02 の長方形ループ & +2.98°/周 \\
\textbf{最適円ループ}（中心 腕 −20°, 脚 −40°、半径 70°） & \textbf{+6.70°/周} \\
可動域内の理論上界（同符号曲率の全領域を囲む） & +9.57°/周 \\
\end{longtable}

\begin{itemize}
\tightlist
\item
  四肢の Γ は RW の 200 倍以上ある。四肢は「1 回の動きで大きく向きを変える」装置。
\item
  ループ最適化で 1 周あたりの回転は 2.25 倍になった。理論上界の 70\% に達している。 曲率マップの破線（F=0）が「囲むべき領域」の境界になっている。
\item
  上界は関節可動域で決まる。\textbf{AMBAC 性能は可動域設計の問題}でもある。
\end{itemize}

\section{Part B: 90° の姿勢変更（静止→静止）}

\FIG{sim03_rest2rest.png}

\begin{longtable}[]{@{}
  >{\raggedright\arraybackslash}p{(\columnwidth - 4\tabcolsep) * \real{0.3333}}
  >{\raggedright\arraybackslash}p{(\columnwidth - 4\tabcolsep) * \real{0.3333}}
  >{\raggedright\arraybackslash}p{(\columnwidth - 4\tabcolsep) * \real{0.3333}}@{}}
\toprule\noalign{}
\begin{minipage}[b]{\linewidth}\raggedright
\end{minipage} & \begin{minipage}[b]{\linewidth}\raggedright
RW
\end{minipage} & \begin{minipage}[b]{\linewidth}\raggedright
AMBAC
\end{minipage} \\
\midrule\noalign{}
\endhead
\bottomrule\noalign{}
\endlastfoot
所要時間 & 10 s（設定値） & 26 s（13 周 × 0.5 Hz） \\
実行量 & ホイール 178 回転、最大 2003 rpm & 四肢ループ 13 周 \\
最終姿勢 & 90.00° & 87.16°（整数周回の端数） \\
途中の姿勢 & 単調・滑らか & \textbf{最大 113°（周回ごとに ±20〜25° 揺れる）} \\
\end{longtable}

\textbf{考察} - AMBAC では、正味の回転の上に、形状変化による大きな \textbf{可逆的な揺れ} が重なる。 観測・照準・作業を伴う姿勢制御では、この揺れは実用上の大問題になる。 対策としては、ループ形状を「揺れが小さく曲率が大きい」ものに最適化すること、 他の肢で揺れを打ち消すこと（冗長肢の利用）が考えられる。 - 端数（87° vs 90°）は、最後の部分ループや単関節の可逆運動で微調整できる （可逆 AMBAC＝微調整、非ホロノミック AMBAC＝大角度、という役割分担）。

\section{Part C: 外乱トルク下の姿勢保持と「飽和」}

\FIG{sim03_disturbance.png}

一定の外乱トルク τ\_d = 0.01 N·m を胴体に加える。全角運動量は L(t) = τ\_d t で増え続け、 \textbf{内部運動では L を消せない}。姿勢を保つには、増えた L をどこかに「溜める」しかない。

\begin{itemize}
\tightlist
\item
  \textbf{RW}: ホイールの回転に溜める。ホイール回転数が時間に比例して増え、ω\_max に達すると飽和する。
\item
  \textbf{AMBAC}: \textbf{四肢の周期運動に溜める}。ループ周回周波数 f が時間に比例して増え、f\_max に達すると飽和する。
\end{itemize}

AMBAC の平均化モデル（1 周平均で dθ̄/dt ≈ f·Δθ\_c + L·$\langle$1/A$\rangle$）から、角運動量容量は次のとおり。

\[
h_\mathrm{RW} = I_w\,\omega_{\max}, \qquad
h_\mathrm{AMBAC} = \frac{f_{\max}\,|\Delta\theta_c|}{\langle 1/A\rangle}
\]

\begin{longtable}[]{@{}
  >{\raggedright\arraybackslash}p{(\columnwidth - 4\tabcolsep) * \real{0.3333}}
  >{\raggedright\arraybackslash}p{(\columnwidth - 4\tabcolsep) * \real{0.3333}}
  >{\raggedright\arraybackslash}p{(\columnwidth - 4\tabcolsep) * \real{0.3333}}@{}}
\toprule\noalign{}
\begin{minipage}[b]{\linewidth}\raggedright
\end{minipage} & \begin{minipage}[b]{\linewidth}\raggedright
角運動量容量
\end{minipage} & \begin{minipage}[b]{\linewidth}\raggedright
0.01 N·m での保持限界
\end{minipage} \\
\midrule\noalign{}
\endhead
\bottomrule\noalign{}
\endlastfoot
RW（2 kg、5700 rpm） & 8.64 kg·m²/s & 864 s \\
AMBAC（腕 4 kg + 脚 12 kg、1 Hz） & 1.18 kg·m²/s & 118 s \\
\end{longtable}

保持中の AMBAC は、周回平均姿勢 \textbar θ̄\textbar{} ≤ 0.30° に保てたが、瞬時姿勢は ±23° 揺れる。

\textbf{考察} - 「AMBAC で回転を止められない」ことと「RW の運動量飽和」は \textbf{同一の物理現象} である。 RW が角運動量をホイールの回転に溜めるのと同じく、AMBAC は四肢の周期運動に溜める （連続的なループ運動は、実効的に「回り続けるホイール」になっている）。 - この人間大モデルでは、2 kg の RW の方が 16 kg の四肢より 7 倍以上の容量を持つ。 AMBAC の強みは容量ではなく、\textbf{推進剤も専用機器も使わず、既存の四肢で姿勢を変えられる}ことにある。 - どちらも飽和したら、外部トルク（スラスタ、磁気トルカ、接触）による \textbf{アンローディング} が必要になる。 劇中で AMBAC とスラスタが併用されていることは、物理的に自然な設計と言える。 - 周回内の揺れの補償には、ループ形状から計算した周期成分 θ\_osc(φ) を使った。 \textbf{幾何構造（接続・曲率）を知っていると、制御設計が格段に簡単になる}。 これはフィジカル AI（学習で方策を得る場合）が「何を内部表現として獲得すべきか」という問いに直結する。

\section{次への論点}

\begin{enumerate}
\def\labelenumi{\arabic{enumi}.}
\tightlist
\item
  \textbf{トルクとエネルギー（Sim04）}: f\_max は本来アクチュエータのトルク・出力で決まる。 順動力学で、関節トルク・消費エネルギー・f\_max を見積もる。
\item
  \textbf{スケーリング（18 m 級）}: Δθ\_c はスケール不変だが、A ∝ M ℓ² であり、f\_max もスケールに依存する。 巨大ロボで h\_AMBAC がどう伸びるかは、アクチュエータのスケーリング則次第（Sim04 で検証予定）。
\item
  \textbf{揺れの小さいループ}: 「正味回転は大きく、周回内の揺れは小さい」ループ形状の多目的最適化。
\item
  \textbf{フィジカル AI}: RW＋AMBAC＋スラスタの統合制御（運動量の配分）を、最適制御と学習でそれぞれ解いて比較する。
\end{enumerate}

\Detail{ch03}
\chapter{順動力学・トルク・スケーリング則（Sim04）}\label{ch:dyn}

\section{運動方程式}

導出とスケーリング則の詳細は付録~\ref{app:eom} を参照。


一般化速度 u = (θ̇, q̇) とする。外力なし・全運動量ゼロなら系重心は静止し、 運動エネルギーは T = ½ u\textsuperscript{T}M(q)u（θ に依存しない）。ラグランジュ方程式は次のとおり。

\[
M(q)\dot u + \dot M(q) u - \tfrac12\left[0;\ u^\top \tfrac{\partial M}{\partial q_k} u\right] = [0;\ \tau]
\]

\begin{itemize}
\tightlist
\item
  M の第 1 行は (A, B) そのもので、第 1 式は dL/dt = 0 を表す。
\item
  関節トルク τ は胴体とリンクの間の \textbf{内力} なので、胴体には直接トルクがかからない。
\end{itemize}

\textbf{時間スケーリングの性質（厳密）}: 同じ経路を f 倍速で回すと、τ = M q̈ + (速度の 2 次項) なので \textbf{τ ∝ f²、パワー ∝ f³}。一方、回転角は経路のみで決まり、f に依存しない（第~\ref{ch:theory}章）。

\section{Part A: 検証}

関節トルク τ(t) だけを入力し、開ループで 2 周分（0.5 Hz）を数値積分した。

\begin{longtable}[]{@{}
  >{\raggedright\arraybackslash}p{(\columnwidth - 2\tabcolsep) * \real{0.5000}}
  >{\raggedright\arraybackslash}p{(\columnwidth - 2\tabcolsep) * \real{0.5000}}@{}}
\toprule\noalign{}
\begin{minipage}[b]{\linewidth}\raggedright
検証項目
\end{minipage} & \begin{minipage}[b]{\linewidth}\raggedright
結果
\end{minipage} \\
\midrule\noalign{}
\endhead
\bottomrule\noalign{}
\endlastfoot
関節軌道の追従誤差 & 2.6e-8° \\
2 周後の正味回転 & +13.4086°（運動学モデル Sim03 と一致） \\
角運動量 & max 2e-9 kg·m²/s（保存） \\
エネルギー収支 ΔT − ∫P dt & 9e-4 J（台形積分誤差の範囲） \\
\end{longtable}

→ 運動学モデル（Sim01〜03）と動力学モデルは整合している。

\section{Part B: 人間大（76 kg）}

\FIG{sim04_human_cycle.png}

\begin{longtable}[]{@{}
  >{\raggedright\arraybackslash}p{(\columnwidth - 2\tabcolsep) * \real{0.5000}}
  >{\raggedright\arraybackslash}p{(\columnwidth - 2\tabcolsep) * \real{0.5000}}@{}}
\toprule\noalign{}
\begin{minipage}[b]{\linewidth}\raggedright
項目
\end{minipage} & \begin{minipage}[b]{\linewidth}\raggedright
0.5 Hz で最適ループを回す場合
\end{minipage} \\
\midrule\noalign{}
\endhead
\bottomrule\noalign{}
\endlastfoot
必要トルク peak & 肩 8.2 N·m、股 28.8 N·m \\
パワー peak & 肩 21 W、股 47 W \\
平均正仕事率（回生なし） & 18.5 W \\
持続トルク目安（肩 40、股 120 N·m）での f\_max & \textbf{1.02 Hz} \\
90° 旋回時間（f\_max） & \textbf{13.2 s} \\
\end{longtable}

\begin{itemize}
\tightlist
\item
  人間の筋力で十分実行できる負荷である。90° に十数秒という数値は、宇宙飛行士の 「腕や脚を回して、ゆっくりなら向きを変えられる」という経験と整合する。
\item
  パワーは周期的に正負を繰り返し、1 周あたりの正味仕事はゼロ（内力のみ・摩擦なし）。 \textbf{回生・弾性要素（ばね）を使えば、消費エネルギーは原理的にほぼゼロにできる}。 動物が腱の弾性を使うのと同じ発想で、フィジカル AI の「身体性の活用」の論点になる。
\end{itemize}

\section{Part C: スケーリング則 --- 二乗三乗則が AMBAC を直撃する}

\FIG{sim04_scaling.png}

人間大モデルを長さ s 倍・同密度で相似拡大した（s=10 で約 18 m・76 t）。 アクチュエータは「同じ技術で相似拡大」と仮定し、最大トルク ∝ s³（断面積 s² × レバー s）とした。

\begin{longtable}[]{@{}
  >{\raggedright\arraybackslash}p{(\columnwidth - 4\tabcolsep) * \real{0.3333}}
  >{\raggedright\arraybackslash}p{(\columnwidth - 4\tabcolsep) * \real{0.3333}}
  >{\raggedright\arraybackslash}p{(\columnwidth - 4\tabcolsep) * \real{0.3333}}@{}}
\toprule\noalign{}
\begin{minipage}[b]{\linewidth}\raggedright
量
\end{minipage} & \begin{minipage}[b]{\linewidth}\raggedright
スケーリング（理論）
\end{minipage} & \begin{minipage}[b]{\linewidth}\raggedright
数値で得た指数
\end{minipage} \\
\midrule\noalign{}
\endhead
\bottomrule\noalign{}
\endlastfoot
1 周あたりの回転角 Δθ\_c & s\textsuperscript{0}（不変） & --- \\
同じ周波数で回すのに必要なトルク & s\textsuperscript{5} & 5.00 \\
達成可能な周回周波数 f\_max & s\textsuperscript{-1} & −1.00 \\
90° 旋回時間 & s\textsuperscript{1} & --- \\
AMBAC の角運動量容量 h & s\textsuperscript{4} & 4.00 \\
\end{longtable}

\begin{longtable}[]{@{}
  >{\raggedright\arraybackslash}p{(\columnwidth - 8\tabcolsep) * \real{0.2000}}
  >{\raggedright\arraybackslash}p{(\columnwidth - 8\tabcolsep) * \real{0.2000}}
  >{\raggedright\arraybackslash}p{(\columnwidth - 8\tabcolsep) * \real{0.2000}}
  >{\raggedright\arraybackslash}p{(\columnwidth - 8\tabcolsep) * \real{0.2000}}
  >{\raggedright\arraybackslash}p{(\columnwidth - 8\tabcolsep) * \real{0.2000}}@{}}
\toprule\noalign{}
\begin{minipage}[b]{\linewidth}\raggedright
s
\end{minipage} & \begin{minipage}[b]{\linewidth}\raggedright
全長
\end{minipage} & \begin{minipage}[b]{\linewidth}\raggedright
質量
\end{minipage} & \begin{minipage}[b]{\linewidth}\raggedright
f\_max
\end{minipage} & \begin{minipage}[b]{\linewidth}\raggedright
90° 旋回
\end{minipage} \\
\midrule\noalign{}
\endhead
\bottomrule\noalign{}
\endlastfoot
0.5 & 約 0.9 m & 9.5 kg & 2.0 Hz & 6.6 s \\
1 & 約 1.8 m & 76 kg & 1.0 Hz & 13 s \\
10 & 約 18 m & 76 t & 0.10 Hz & \textbf{132 s} \\
20 & 約 36 m & 608 t & 0.05 Hz & 263 s \\
\end{longtable}

\textbf{18 m 級について} - 人間と同じ機敏さ（0.5 Hz）で回すには、股関節に 2.9 MN·m が必要になる（人間の 10 万倍）。 - 90° を 10 s で旋回するには 20.8 MN·m が必要で、\textbf{相似則で拡大したアクチュエータの 173 倍} にあたる。 - 相似則のままなら、90° の旋回に 2 分以上かかる。

\textbf{考察} - 「回転角はスケール不変」（Sim01）でも、\textbf{速さは 1/s で落ちる}。巨大ロボの AMBAC は遅い。 - 逆に、小さい生物ほど AMBAC が速い。猫が空中で素早く向きを変えられるのは、この則の小さい側にいるため。 - 劇中のように機敏な AMBAC を実現するには、人間の筋肉や既存のモータに比べて \textbf{トルク密度が 2 桁以上高いアクチュエータ} が必要になる。これは物理的な要求として定量化できた。

\section{Part D: スラスタとの比較（18 m 級）}

慣性モーメント約 1.0×10\textsuperscript{6} kg·m²、スラスタのモーメントアーム 5 m、Isp 300 s（化学推進相当）とし、 90° 旋回（静止→静止、加速→減速）に必要な推進剤を求めた。推進剤は 1/T に比例する。

\begin{longtable}[]{@{}
  >{\raggedright\arraybackslash}p{(\columnwidth - 4\tabcolsep) * \real{0.3333}}
  >{\raggedright\arraybackslash}p{(\columnwidth - 4\tabcolsep) * \real{0.3333}}
  >{\raggedright\arraybackslash}p{(\columnwidth - 4\tabcolsep) * \real{0.3333}}@{}}
\toprule\noalign{}
\begin{minipage}[b]{\linewidth}\raggedright
旋回時間 T
\end{minipage} & \begin{minipage}[b]{\linewidth}\raggedright
スラスタ推進剤 / 回
\end{minipage} & \begin{minipage}[b]{\linewidth}\raggedright
AMBAC
\end{minipage} \\
\midrule\noalign{}
\endhead
\bottomrule\noalign{}
\endlastfoot
5 s & 88 kg & 不可能（相似則アクチュエータ） \\
10 s & 44 kg & 不可能（173 倍のトルクが必要） \\
30 s & 15 kg & 不可能 \\
132 s & 3.4 kg & \textbf{推進剤 0 kg}、関節の正仕事 約 2.1 MJ（回生なしの場合） \\
\end{longtable}

\textbf{考察} - ゆっくり回るならスラスタの推進剤も少なくて済む。AMBAC の節約効果が大きくなるのは素早い旋回のときだが、 巨大機では素早い旋回こそ AMBAC が苦手とする領域である。\textbf{トレードオフが逆向きに噛み合っている}。 - ただし、推進剤ゼロという性質は、長期間の運用・推進剤が尽きた状況・排気を出せない状況 （汚染、被探知）では決定的な利点になる。 - 合理的な設計は \textbf{ハイブリッド} になる。素早い旋回はスラスタ（あるいは大トルクの CMG）で行い、 低速の姿勢調整・外乱の吸収・推進剤節約は AMBAC で行う。 「どの時間スケールの運動をどのアクチュエータに配分するか」は、そのまま最適制御・学習の問題になる。

\section{次への論点}

\begin{enumerate}
\def\labelenumi{\arabic{enumi}.}
\tightlist
\item
  \textbf{統合制御（フィジカル AI 編）}: スラスタ・RW/CMG・AMBAC の配分を、推進剤・時間・揺れの多目的で 最適制御と強化学習によって解き、比較する。
\item
  \textbf{弾性・回生}: 直列弾性アクチュエータ（ばね）で周回運動のエネルギーをほぼゼロにできるか。 共振周波数とループ周波数の整合。
\item
  \textbf{3 次元化}: 3 次元ではロール・ピッチ・ヨーの結合（非可換性）により、平面では見えない回転が生じる。
\end{enumerate}

\Detail{ch04}
\chapter{フィジカル AI 編：学習は AMBAC を発見できるか、どう配分すべきか（Sim05）}\label{ch:ai}

\section{問題設定と問い}

ここまでで分かったことは次の 3 点である。

\begin{itemize}
\tightlist
\item
  AMBAC は \textbf{非ホロノミック}（形状ループの幾何学的位相で回る）。
\item
  AMBAC は \textbf{周回内で大きく揺れる}。
\item
  巨大機では \textbf{遅い}。
\end{itemize}

この性質を持つ系を、フィジカル AI はどう扱うべきか。本章では次の 2 つの問いを立てる。

\begin{enumerate}
\def\labelenumi{\arabic{enumi}.}
\tightlist
\item
  \textbf{発見}: 物理の知識を与えずに「胴体を目標角まで回せ」とだけ学習させたとき、エージェントは AMBAC（形状空間のループ）を自力で発見できるか。発見に必要な前提知識（帰納バイアス）は何か。
\item
  \textbf{配分}: スラスタ・リアクションホイール・AMBAC を併用するとき、最適制御（モデルベース）と学習（モデルフリー）はそれぞれどう配分し、モデル誤差・外乱にどう反応するか。
\end{enumerate}

学習手法には進化戦略（ES）を使った。勾配を使わない方策探索で、報酬が終端でしか決まらない問題や、 長いロールアウトに強い。強化学習の一種（ブラックボックス方策最適化）として扱う。

\begin{center}\rule{0.5\linewidth}{0.5pt}\end{center}

\section{Part A: 学習エージェントは AMBAC を発見できるか（人間大）}

\FIG{sim05a_learned_ambac.png}

\textbf{設定}: L=0 の運動学モデル（腕＋脚）。行動は関節角速度（上限は Sim04 の f\_max 相当）。 報酬は姿勢誤差の時間平均と終端誤差。接続・曲率・ループの知識は一切与えない。

\begin{longtable}[]{@{}
  >{\raggedright\arraybackslash}p{(\columnwidth - 2\tabcolsep) * \real{0.5000}}
  >{\raggedright\arraybackslash}p{(\columnwidth - 2\tabcolsep) * \real{0.5000}}@{}}
\toprule\noalign{}
\begin{minipage}[b]{\linewidth}\raggedright
方策の構造
\end{minipage} & \begin{minipage}[b]{\linewidth}\raggedright
中身
\end{minipage} \\
\midrule\noalign{}
\endhead
\bottomrule\noalign{}
\endlastfoot
\textbf{plain} & 観測（姿勢誤差, 関節角）→ MLP → 関節角速度 \\
\textbf{cpg} & plain の観測に周期信号 sin/cos(2πft) を追加 \\
\textbf{tg} & 振動子（内部位相 ψ）が目標関節角 q* = c + A$\odot$(cos ψ, sin(ψ+δ)) を生成する。学習するのはループの中心 c・振幅 A・形 δ と、姿勢誤差 → 位相速度のフィードバック則（MLP）。PMTG \cite{R21} と同系統の「軌道生成器を学習で調整する」構造 \\
expert（比較） & Sim03 の最適円ループを、揺れ補償付きフィードバックで回すモデルベース制御 \\
\end{longtable}

\textbf{結果（25 s 後の姿勢誤差）}

\begin{longtable}[]{@{}
  >{\raggedright\arraybackslash}p{(\columnwidth - 8\tabcolsep) * \real{0.2000}}
  >{\raggedright\arraybackslash}p{(\columnwidth - 8\tabcolsep) * \real{0.2000}}
  >{\raggedright\arraybackslash}p{(\columnwidth - 8\tabcolsep) * \real{0.2000}}
  >{\raggedright\arraybackslash}p{(\columnwidth - 8\tabcolsep) * \real{0.2000}}
  >{\raggedright\arraybackslash}p{(\columnwidth - 8\tabcolsep) * \real{0.2000}}@{}}
\toprule\noalign{}
\begin{minipage}[b]{\linewidth}\raggedright
目標
\end{minipage} & \begin{minipage}[b]{\linewidth}\raggedright
plain
\end{minipage} & \begin{minipage}[b]{\linewidth}\raggedright
cpg
\end{minipage} & \begin{minipage}[b]{\linewidth}\raggedright
\textbf{tg}
\end{minipage} & \begin{minipage}[b]{\linewidth}\raggedright
expert
\end{minipage} \\
\midrule\noalign{}
\endhead
\bottomrule\noalign{}
\endlastfoot
+90° & 61.8° & 61.8° & \textbf{0.2°} & 8.7° \\
−90° & 49.7° & 49.7° & \textbf{0.1°} & 7.9° \\
+45° & 16.8° & 16.8° & \textbf{0.1°} & 6.1° \\
+150° & 121.8° & 121.8° & \textbf{0.4°} & 21.3° \\
報酬 & −3.00 & −3.00 & \textbf{−0.46} & −1.18 \\
\end{longtable}

\textbf{tg が学習したループ}: 中心（腕 −11°, 脚 −10°）、振幅（120°, 128°）。可動域で切り取られ、 \textbf{曲率 F\textless0 の領域をほぼ埋め尽くす形} になった。1 周あたり \textbf{+9.34°} で、 Sim03 の最適円（+6.70°）を上回り、可動域内の理論上界（+9.57°）の \textbf{98\%} に達する。 回す向きも、目標 +90° では反時計回り（掃引面積 +33.5 rad²）、−90° では時計回り（−33.3 rad²）で、 曲率の符号と整合している。

\subsection{考察 A-1: なぜ構造なしでは失敗するのか ---「幾何学的位相は 2 次の効果」}

plain と cpg は、手足を可動域の端まで振って止まる（×印）という局所解に収束した。 これは単なるハイパーパラメータの問題ではなく、物理的な理由がある。

\begin{itemize}
\tightlist
\item
  1 回の動きの効果（可逆的な首振り, −Γ·Δq）は振幅の \textbf{1 次}。
\item
  ループの正味回転（曲率 × 囲んだ面積）は振幅の \textbf{2 次}（∝ A\textsubscript{1}A\textsubscript{2}）。
\end{itemize}

静止に近い方策から探索を始めると、ループを大きくする方向の勾配はほぼゼロになり、 1 次の効果である「端まで振る」が先に最適化されてしまう。 制御理論の言葉では、正味回転は入力ベクトル場の \textbf{リー括弧}（2 次の項）から生じる。 Brockett の定理で、滑らかな時不変フィードバックでは安定化できないのと同じ構造である。

\subsection{考察 A-2: 正しい帰納バイアスを与えれば、学習は専門家を超える}

周期信号を観測に足すだけ（cpg）では不十分だった。\textbf{「ループを描く振動子」という構造} を与え、 その位置・大きさ・回す向き・速さを学習に任せたところ（tg）、数十イテレーションで解を見つけた。

\begin{itemize}
\tightlist
\item
  学習したループは、人間が設計した「最適円」より大きい（円という形の制約がないため）。
\item
  学習方策は最後に \textbf{揺れのない位相で止まる} ことを覚え、終端誤差が専門家より 1 桁以上小さい。 専門家は周回平均の姿勢 θ̄ を制御するので、止まった位相の揺れがそのまま誤差に残る。
\end{itemize}

\textbf{フィジカル AI への示唆}: 身体のダイナミクスが持つ幾何構造（ここでは「周期運動 × 曲率」）を 方策の構造として埋め込むことが、学習の成否を分ける。歩行ロボットで CPG や 軌道生成器が有効なのと同じ原理を、AMBAC という別の非ホロノミック系で確認できた。

\begin{center}\rule{0.5\linewidth}{0.5pt}\end{center}

\section{Part B: 最適制御による配分（18 m 級）}

線形計画の定式化は付録~\ref{app:lp} を参照。


\FIG{sim05b_pareto.png}

\textbf{モデル}: 平均化モデル θ̄' = (L−h)/Ā + fΔθc、L' = τ\_th、h' = τ\_rw。 90° 静止→静止旋回の推進剤最小化を \textbf{線形計画（LP）} として厳密に解いた。

\begin{longtable}[]{@{}
  >{\raggedright\arraybackslash}p{(\columnwidth - 2\tabcolsep) * \real{0.5000}}
  >{\raggedright\arraybackslash}p{(\columnwidth - 2\tabcolsep) * \real{0.5000}}@{}}
\toprule\noalign{}
\begin{minipage}[b]{\linewidth}\raggedright
機器
\end{minipage} & \begin{minipage}[b]{\linewidth}\raggedright
主な値
\end{minipage} \\
\midrule\noalign{}
\endhead
\bottomrule\noalign{}
\endlastfoot
スラスタ & τ\_max 2.5×10\textsuperscript{5} N·m、アーム 5 m、Isp 300 s \\
AMBAC & Δθc 6.70°/周、f\_max 0.102 Hz（Sim04） \\
RW & 2 kg・半径 0.12 m の RW を相似拡大（h\_max 8.6×10\textsuperscript{4} kg·m²/s、τ\_max 200 N·m） \\
\end{longtable}

\begin{longtable}[]{@{}
  >{\raggedright\arraybackslash}p{(\columnwidth - 8\tabcolsep) * \real{0.2000}}
  >{\raggedright\arraybackslash}p{(\columnwidth - 8\tabcolsep) * \real{0.2000}}
  >{\raggedright\arraybackslash}p{(\columnwidth - 8\tabcolsep) * \real{0.2000}}
  >{\raggedright\arraybackslash}p{(\columnwidth - 8\tabcolsep) * \real{0.2000}}
  >{\raggedright\arraybackslash}p{(\columnwidth - 8\tabcolsep) * \real{0.2000}}@{}}
\toprule\noalign{}
\begin{minipage}[b]{\linewidth}\raggedright
旋回時間 T
\end{minipage} & \begin{minipage}[b]{\linewidth}\raggedright
スラスタのみ
\end{minipage} & \begin{minipage}[b]{\linewidth}\raggedright
+RW
\end{minipage} & \begin{minipage}[b]{\linewidth}\raggedright
+AMBAC
\end{minipage} & \begin{minipage}[b]{\linewidth}\raggedright
+AMBAC+RW
\end{minipage} \\
\midrule\noalign{}
\endhead
\bottomrule\noalign{}
\endlastfoot
8 s & 30.1 kg & 30.0 & 28.1 & 28.1 \\
30 s & 7.3 & 7.1 & 5.6 & 5.4 \\
73 s & 3.0 & 2.5 & 1.3 & 0.85 \\
113 s & 1.9 & 1.2 & 0.28 & \textbf{0} \\
推進剤ゼロの最短時間 & --- & 183 s & 137 s & \textbf{97 s} \\
\end{longtable}

（推進剤ゼロの最短時間は格子上の値。AMBAC 単独の理論値は 132 s）

\textbf{考察} - 最適な配分は「AMBAC と RW を常に最大で使い、足りない分だけスラスタで補う」になる。 T=60 s の例（右図）では、AMBAC は全区間フル稼働、RW は前半で加速・後半で減速し、 スラスタは最初と最後に少しだけ噴く。 - 内部機器（AMBAC・RW）は \textbf{ゆっくりした旋回ほど効く}。素早い旋回（10 s 以下）ではほぼスラスタ頼みになる。 これは Sim04 で見た「トレードオフが逆向きに噛み合う」構造の定量版である。 - AMBAC と RW は相補的で、併用すると推進剤ゼロ領域が約 100 s まで広がる。 四肢（AMBAC）は「大きく、遅い」、RW は「小さく、トルクが弱い」が、角運動量の受け皿として足し合わせられる。

\begin{center}\rule{0.5\linewidth}{0.5pt}\end{center}

\section{Part C: モデル誤差・外乱がある場合 --- 最適制御 vs 学習}

\FIG{sim05b_robustness.png}

\textbf{設定}: 18 m 級、スラスタ＋AMBAC、90° 旋回を T=60 s で行う。

\begin{itemize}
\tightlist
\item
  \textbf{評価用プラント}: 周回内の揺れを含む \textbf{位相モデル}（制御器は平均化モデルしか知らない）。
\item
  \textbf{モデル誤差}: Δθc の倍率 κ ∈ {[}0.7, 1.3{]}（ループ形状の誤差）、慣性の倍率 α ∈ {[}0.9, 1.1{]}。
\item
  \textbf{外乱}: 一定トルク τ\_d ∈ {[}−500, 500{]} N·m。
\item
  \textbf{指標}: 終端姿勢誤差（実際の θ）、旋回後に残る回転（残留角速度）、推進剤。
\end{itemize}

\begin{longtable}[]{@{}
  >{\raggedright\arraybackslash}p{(\columnwidth - 2\tabcolsep) * \real{0.5000}}
  >{\raggedright\arraybackslash}p{(\columnwidth - 2\tabcolsep) * \real{0.5000}}@{}}
\toprule\noalign{}
\begin{minipage}[b]{\linewidth}\raggedright
制御器
\end{minipage} & \begin{minipage}[b]{\linewidth}\raggedright
中身
\end{minipage} \\
\midrule\noalign{}
\endhead
\bottomrule\noalign{}
\endlastfoot
開ループ最適制御 & 公称モデルで一度だけ LP を解き、そのまま実行 \\
MPC & 2 s ごとに現在の状態（θ̄, L）から LP を再計画 \\
MPC＋揺れ補償 & MPC に加えて、計画から終端位相 φ\_end を予測し、目標を「90° − 揺れ(φ\_end)」に置換。終盤 10 s は AMBAC を止め、スラスタで仕上げる \\
学習方策（ES） & 観測（θ̄ 誤差, L, 残り時間, 位相 φ）→ MLP → (τ\_th, f)。上記のランダムなプラントで 800 イテレーション学習（ドメインランダム化） \\
\end{longtable}

参考（公称・平均化モデルでの最適値）: スラスタのみ 3.62 kg、スラスタ＋AMBAC 1.97 kg

\begin{longtable}[]{@{}
  >{\raggedright\arraybackslash}p{(\columnwidth - 10\tabcolsep) * \real{0.1667}}
  >{\raggedright\arraybackslash}p{(\columnwidth - 10\tabcolsep) * \real{0.1667}}
  >{\raggedright\arraybackslash}p{(\columnwidth - 10\tabcolsep) * \real{0.1667}}
  >{\raggedright\arraybackslash}p{(\columnwidth - 10\tabcolsep) * \real{0.1667}}
  >{\raggedright\arraybackslash}p{(\columnwidth - 10\tabcolsep) * \real{0.1667}}
  >{\raggedright\arraybackslash}p{(\columnwidth - 10\tabcolsep) * \real{0.1667}}@{}}
\toprule\noalign{}
\begin{minipage}[b]{\linewidth}\raggedright
条件
\end{minipage} & \begin{minipage}[b]{\linewidth}\raggedright
制御器
\end{minipage} & \begin{minipage}[b]{\linewidth}\raggedright
終端誤差 平均（95\%）
\end{minipage} & \begin{minipage}[b]{\linewidth}\raggedright
残留角速度
\end{minipage} & \begin{minipage}[b]{\linewidth}\raggedright
推進剤
\end{minipage} & \begin{minipage}[b]{\linewidth}\raggedright
AMBAC 周回
\end{minipage} \\
\midrule\noalign{}
\endhead
\bottomrule\noalign{}
\endlastfoot
公称 & 開ループ最適制御 & 17.4° & 0 & \textbf{1.97 kg} & 6.1 \\
公称 & MPC & 10.5° & 0 & 3.65 kg & 5.4 \\
公称 & MPC＋揺れ補償 & \textbf{0.02°} & 0 & 3.05 kg & 0.6 \\
公称 & 学習方策 & 0.09° & 0.002°/s & 4.66 kg & 0.3 \\
ランダム 100 例 & 開ループ最適制御 & 30.5°（68.8°） & 0.90°/s & 1.97 kg & 6.1 \\
ランダム 100 例 & MPC & 14.6°（23.3°） & 0.030°/s & 5.41 kg & 2.9 \\
ランダム 100 例 & MPC＋揺れ補償 & 0.19°（0.69°） & 0.031°/s & 5.56 kg & 0.8 \\
ランダム 100 例 & 学習方策 & \textbf{0.09°（0.10°）} & \textbf{0.002°/s} & 4.92 kg & 0.3 \\
\end{longtable}

\subsection{考察 C-1: 平均化モデルは正確だが「揺れ」が最後に効く}

開ループで実行すると、60 s 後の周回平均姿勢 θ̄ は 89.6° で、平均化モデルの予測（90.0°）とよく一致した。 それでも実際の姿勢 θ は 17° ずれる。止まった位相の揺れが誤差としてそのまま残るためである。 \textbf{AMBAC の精密な制御には、平均（ゆっくりした回転）と揺れ（速い形状変化）の 2 つの時間スケールを両方扱う必要がある}。 MPC＋揺れ補償はこれを人手で設計して解決した。学習方策は観測に位相 φ を含めただけで、同じことを自力で獲得した。

\subsection{考察 C-2: 角運動量を一度入れると、AMBAC は推進剤を節約しなくなる}

最適計画では AMBAC を 6 周使って 1.97 kg まで下げているが、MPC 系も学習方策も、実際には AMBAC をあまり使わない。 再計画の過程を追跡すると、理由は次のとおりだった。

\begin{itemize}
\tightlist
\item
  スラスタで角運動量 L を入れると、止めるときに \textbar L\textbar{} ぶんの逆噴射が必ず必要になる（減速の時期を変えても費用は同じ）。
\item
  そのため、\textbf{一度回転を始めた後は、AMBAC を使っても推進剤は 1 g も減らない}。 AMBAC が推進剤を節約できるのは、\textbf{最初にどれだけ噴くかを決める段階} だけである。
\item
  再計画の時点では「AMBAC を使う／使わない」の費用が同じになり、LP は使わない解を選ぶ。 揺れ補償 MPC は、終盤で AMBAC を止める制約の分だけ初期の噴射が増えた。
\end{itemize}

学習方策の推進剤（約 4.7〜4.9 kg）は、スラスタのみの最適値（3.62 kg）より多い。 同じ構造により、学習方策は「まずスラスタで回し始め、精度を出す」解に落ち着き、AMBAC で得をする余地を自ら消している。 \textbf{AMBAC の価値は「最初の判断」に集中しており、報酬が終端でしか与えられない学習では、その判断を学ぶのが難しい}。

\subsection{考察 C-3: 最適制御と学習の得意・不得意}

\begin{longtable}[]{@{}
  >{\raggedright\arraybackslash}p{(\columnwidth - 4\tabcolsep) * \real{0.3333}}
  >{\raggedright\arraybackslash}p{(\columnwidth - 4\tabcolsep) * \real{0.3333}}
  >{\raggedright\arraybackslash}p{(\columnwidth - 4\tabcolsep) * \real{0.3333}}@{}}
\toprule\noalign{}
\begin{minipage}[b]{\linewidth}\raggedright
\end{minipage} & \begin{minipage}[b]{\linewidth}\raggedright
最適制御（LP / MPC）
\end{minipage} & \begin{minipage}[b]{\linewidth}\raggedright
学習方策（ES）
\end{minipage} \\
\midrule\noalign{}
\endhead
\bottomrule\noalign{}
\endlastfoot
公称条件での推進剤 & ◎ 最適（1.97 kg、ただし揺れで精度が出ない） & △ スラスタのみの最適値より悪い \\
モデル誤差・外乱への強さ & △ 再計画すれば精度は出るが、外乱を知らないため回転が残る & ◎ 誤差 0.1°、残留回転は MPC の 1/10 \\
揺れ（2 つの時間スケール）の扱い & 人手の設計が必要 & 観測に位相を入れれば自力で獲得 \\
実行時の計算 & 再計画のたびに LP（ここでは 100 例で 15〜40 s） & 行列積だけ（100 例で 0.1 s） \\
解釈性・保証 & 制約・最適性が明示的 & ブラックボックス \\
\end{longtable}

\textbf{フィジカル AI への示唆} - どちらか一方では足りない。\textbf{計画（いつ・どれだけ噴くか、AMBAC をどれだけ回すか）は最適制御が担い、 実行時の誤差・外乱・揺れの補正は学習方策が担う} という役割分担（残差学習、学習済み MPC コスト）が自然な次の一手になる。 - Part A と合わせると、学習が成功するかどうかは \textbf{物理構造をどこまで方策に埋め込むか} で決まる。 - ループという「形」は構造として与える（Part A の tg）。 - 角運動量という「保存量」の扱いは、最適制御のように明示的に計画する。

物理の知識と学習の分担を設計することが、フィジカル AI の中心的な課題であることを、 AMBAC という具体例で示せた。

\begin{center}\rule{0.5\linewidth}{0.5pt}\end{center}

\section{まとめと次への論点}

\begin{enumerate}
\def\labelenumi{\arabic{enumi}.}
\tightlist
\item
  学習で AMBAC を \textbf{発見できる}。ただし「振動子でループを描く」構造が必要で、構造がないと 2 次効果の壁に阻まれる。 構造を与えた学習は、可動域内の理論上界の 98\% に達するループを見つけ、専門家制御を上回った。
\item
  18 m 級では、AMBAC＋RW により \textbf{約 100 s 以上かける旋回なら推進剤ゼロ} にできる。素早い旋回はスラスタ頼みになる。
\item
  AMBAC が推進剤を節約するのは \textbf{最初の噴射量を決める計画段階だけ}。一方、揺れと外乱への対処は学習が得意である。 → \textbf{計画は最適制御、実行は学習} というハイブリッドが有望。
\end{enumerate}

\textbf{次の候補} - 最適制御の計画に学習した残差方策を重ねるハイブリッド制御の実装と評価 - 3 次元化（ロール・ピッチ・ヨーの非可換性）で同じ問いを立て直す - 弾性要素（ばね）による周回エネルギーの回生と、学習によるリズムの獲得

\Detail{ch05}

\part{下半身の形状（地上）}
\chapter{地上を動く機体：戦車型台車＋腕（Sim06）}\label{ch:tank}

\section{宇宙から地上へ：何が変わるか}

ZMP の式の導出とフルード則は付録~\ref{app:zmp}、スティックスリップのモデルは付録~\ref{app:stickslip} を参照。

\begin{figure}[htbp]\centering
\resizebox{0.85\linewidth}{!}{% 戦車型台車と ZMP
\begin{tikzpicture}[>=Stealth,font=\small,scale=1.2]
\fill[pattern=north east lines,pattern color=gray] (-2.6,-0.15) rectangle (2.8,0);
\draw[thick] (-2.6,0)--(2.8,0);
\draw[thick,fill=gray!20,rounded corners=4pt] (-1.2,0) rectangle (1.2,0.55);
\draw[thick,fill=gray!10] (-0.35,0.55) rectangle (0.35,1.6);
\fill (0,0.9) circle (1.5pt) node[left,font=\footnotesize]{重心};
\coordinate (S) at (0,1.6);
\draw[line width=3pt,blue!60] (S)--++(30:1.0) coordinate (E) --++(-20:1.0) coordinate (T);
\fill[orange] (T) circle (4pt) node[right=3pt,font=\footnotesize]{荷物 10 kg};
\draw[->,thick,orange!80!black] (T)--++(-0.9,0) node[left,font=\footnotesize]{$-m\ddot x$};
\draw[<->] (-1.2,-0.45)--(1.2,-0.45) node[midway,below,font=\footnotesize,xshift=-18pt]{接地範囲 $[-a,a]$};
\draw[->,very thick,red] (0.7,-0.0)++(0,-0.9)--(0.7,0) node[pos=0,left,font=\footnotesize]{$N$};
\draw[->,very thick,red] (0.7,0.05)--++(-0.6,0) node[above,font=\footnotesize]{$F_x$};
\fill[red] (0.7,0) circle (2pt);
\node[red,font=\footnotesize] at (2.3,-1.0) {ZMP $p$};
\draw[red,->] (2.0,-0.95)--(0.78,-0.05);
\node[align=left,font=\footnotesize] at (4.6,2.6) {転倒しない：$|p|\le a$\\滑らない：$|F_x|\le\mu N$};
\end{tikzpicture}
}
\caption{戦車型台車の腕の反作用。地面の反力の合力が通る点（ZMP）が接地範囲を出ると転倒、水平反力が摩擦を超えると滑る。}\label{fig:tikz-zmp}
\end{figure}


宇宙では外力がないので、角運動量が保存された（dL/dt = 0）。 地上では地面からの反力が働くので、角運動量の変化は接触反力のモーメントに等しい（dL/dt = 接触モーメント）。 ただし、接触反力には次の限界がある。

\begin{longtable}[]{@{}
  >{\raggedright\arraybackslash}p{(\columnwidth - 4\tabcolsep) * \real{0.3333}}
  >{\raggedright\arraybackslash}p{(\columnwidth - 4\tabcolsep) * \real{0.3333}}
  >{\raggedright\arraybackslash}p{(\columnwidth - 4\tabcolsep) * \real{0.3333}}@{}}
\toprule\noalign{}
\begin{minipage}[b]{\linewidth}\raggedright
限界
\end{minipage} & \begin{minipage}[b]{\linewidth}\raggedright
条件
\end{minipage} & \begin{minipage}[b]{\linewidth}\raggedright
超えると
\end{minipage} \\
\midrule\noalign{}
\endhead
\bottomrule\noalign{}
\endlastfoot
ZMP（反力の合力の作用点）は接地範囲の中にしか来られない & −a ≤ p ≤ a & \textbf{転倒} \\
水平反力は摩擦の範囲までしか出せない & \textbar Fx\textbar{} ≤ μN & \textbf{滑り} \\
旋回（ヨー）の反力も摩擦の範囲まで & \textbar τ\textbar{} ≤ μ N r̄ & \textbf{その場で回る} \\
\end{longtable}

\[
p = \frac{\sum m_i\left(x_i(\ddot z_i+g) - z_i\ddot x_i\right) + \sum I_i\ddot\phi_i}{\sum m_i(\ddot z_i+g)}
\]

腕を動かしたときの反力のうち、動的な部分は「角運動量の時間変化」そのものである。 \textbf{地上版 AMBAC とは、腕の運動量の変化が接触の限界を超えるかどうかの問題} と言える。

\textbf{モデル}: 戦車型台車（台車＋胴体 80 kg、履帯接地長 0.8 m）に 2 関節の腕（上腕 4 kg・前腕 3 kg、各 0.5 m）、 先端に荷物 10 kg（総重量 97 kg）。脚の問題を切り離すため、脚は戦車型の台車として単純化した。剛体を仮定する。

\section{Part A: 腕の反作用で転倒するか（矢状面）}

\FIG{sim06_tipover.png}

\textbf{動作}: 前方の地面にある荷物を拾い、頭上を通して後部デッキに置く（最小躍度の単純な関節軌道）。

\begin{longtable}[]{@{}
  >{\raggedright\arraybackslash}p{(\columnwidth - 4\tabcolsep) * \real{0.3333}}
  >{\raggedright\arraybackslash}p{(\columnwidth - 4\tabcolsep) * \real{0.3333}}
  >{\raggedright\arraybackslash}p{(\columnwidth - 4\tabcolsep) * \real{0.3333}}@{}}
\toprule\noalign{}
\begin{minipage}[b]{\linewidth}\raggedright
動作時間 T
\end{minipage} & \begin{minipage}[b]{\linewidth}\raggedright
ZMP の範囲（接地範囲 ±0.4 m）
\end{minipage} & \begin{minipage}[b]{\linewidth}\raggedright
必要な摩擦係数
\end{minipage} \\
\midrule\noalign{}
\endhead
\bottomrule\noalign{}
\endlastfoot
4.0 s & −0.12 〜 +0.17 m & 0.04 \\
2.0 s & −0.21 〜 +0.27 m & 0.15 \\
\textbf{1.62 s} & −0.28 〜 +0.35 m（安全余裕 5 cm の限界） & 0.23 \\
1.2 s & −0.49 〜 \textbf{+0.56 m（転倒）} & 0.46 \\
\end{longtable}

\begin{itemize}
\tightlist
\item
  腕を持ち上げて加速する前半は、ZMP が前方に出る（前のめり）。後方で減速する後半は、後方に出る（のけぞり）。
\item
  摩擦係数 0.6 の路面では、\textbf{滑るより先に転倒する}（滑らない限界 1.09 s ＜ 転倒しない限界 1.62 s）。
\end{itemize}

\subsection{スケーリング：フルード則 \cite{R29}}

\begin{longtable}[]{@{}
  >{\raggedright\arraybackslash}p{(\columnwidth - 6\tabcolsep) * \real{0.2500}}
  >{\raggedright\arraybackslash}p{(\columnwidth - 6\tabcolsep) * \real{0.2500}}
  >{\raggedright\arraybackslash}p{(\columnwidth - 6\tabcolsep) * \real{0.2500}}
  >{\raggedright\arraybackslash}p{(\columnwidth - 6\tabcolsep) * \real{0.2500}}@{}}
\toprule\noalign{}
\begin{minipage}[b]{\linewidth}\raggedright
長さ倍率 s
\end{minipage} & \begin{minipage}[b]{\linewidth}\raggedright
地球
\end{minipage} & \begin{minipage}[b]{\linewidth}\raggedright
火星
\end{minipage} & \begin{minipage}[b]{\linewidth}\raggedright
月
\end{minipage} \\
\midrule\noalign{}
\endhead
\bottomrule\noalign{}
\endlastfoot
1（人間程度） & 1.62 s & 2.62 s & 3.97 s \\
2 & 2.29 s & 3.71 s & 5.61 s \\
10（18 m 級） & 5.11 s & 8.30 s & 12.55 s \\
\end{longtable}

安全に動ける最短時間は \textbf{T ∝ √(s/g)}（フルード則）に \textbf{厳密に} 従った（T/√(s·g\_地球/g) = 1.616 で一定）。 これは時間スケーリングの性質から導ける。動的な反力は 1/T² に比例し、重力は一定なので、 限界は「加速度/重力」の比（フルード数）だけで決まる。

\textbf{考察} - \textbf{大型化}: Sim04（宇宙）では、アクチュエータの制約で速さが 1/s で落ちた。地上では、アクチュエータが十分強くても、 \textbf{転倒の制約だけで速さが 1/√s で落ちる}。18 m 級の機体は、人間程度の機体の約 3 倍ゆっくり腕を動かさないと転倒する。 - \textbf{低重力}: 月面では重力が 1/6 なので、同じ動作に地球の約 2.5 倍の時間がかかる。 \textbf{低重力天体では、腕の動きの反作用が相対的に大きく効く}。月面ローバーにマニピュレータを載せる場合に直接関係する。

\section{Part B: 腕の動かし方で転倒を防ぐ（地上版の能動的質量バランス）}

関節軌道を最小躍度＋補正（関節あたり 4 つの基底関数）で表し、 転倒・滑り・関節速度（6 rad/s）・関節加速度（60 rad/s²）の制約の下で、動作時間を最小化した（SLSQP）。

\begin{longtable}[]{@{}
  >{\raggedright\arraybackslash}p{(\columnwidth - 4\tabcolsep) * \real{0.3333}}
  >{\raggedright\arraybackslash}p{(\columnwidth - 4\tabcolsep) * \real{0.3333}}
  >{\raggedright\arraybackslash}p{(\columnwidth - 4\tabcolsep) * \real{0.3333}}@{}}
\toprule\noalign{}
\begin{minipage}[b]{\linewidth}\raggedright
\end{minipage} & \begin{minipage}[b]{\linewidth}\raggedright
単純軌道（最小躍度）
\end{minipage} & \begin{minipage}[b]{\linewidth}\raggedright
最適化軌道
\end{minipage} \\
\midrule\noalign{}
\endhead
\bottomrule\noalign{}
\endlastfoot
地球 & 1.62 s & \textbf{1.03 s（36\% 短縮）} \\
月 & 3.97 s & \textbf{2.01 s（49\% 短縮）} \\
1.03 s で実行したときの ZMP & −0.72 〜 +0.78 m（転倒） & −0.35 〜 +0.35 m（安全余裕ぎりぎりを使い切る） \\
\end{longtable}

\textbf{最適化で見つかった戦略}（左下図） - \textbf{荷物を体に引き寄せて、低く小さな弧で運ぶ。} 回転半径を小さくして、腕と荷物の慣性（＝角運動量の変化）を減らしている。 - ZMP が許容範囲の上限・下限に \textbf{張り付く} 形になる（右下図）。最短時間制御で典型的な「制約をぎりぎりまで使う」解である。 - これは、人間が重い物を素早く運ぶときに体に引き寄せるのと同じ戦略である。 物理モデルの最適化が、人間らしい動きを \textbf{再発見} した例と言える（第~\ref{ch:conclusion}章の「人体運動を基準値にするか」の論点への一つの答え： 人間の動きの合理性は、物理と制約から説明できる場合がある）。

\section{Part C: 腕を振ると台車がその場で回るか（上から見たヨー）}

\FIG{sim06_yaw_ratchet.png}

\textbf{設定}: 台車中心の旋回関節（ターレット）で、腕（7 kg・1 m）＋荷物（10 kg）を水平に 90° 振る。 台車の並進は止まっているとし、ヨー回転だけを考える。履帯の摩擦は、クーロン摩擦（静止・動摩擦）でモデル化した。

\begin{longtable}[]{@{}
  >{\raggedright\arraybackslash}p{(\columnwidth - 4\tabcolsep) * \real{0.3333}}
  >{\raggedright\arraybackslash}p{(\columnwidth - 4\tabcolsep) * \real{0.3333}}
  >{\raggedright\arraybackslash}p{(\columnwidth - 4\tabcolsep) * \real{0.3333}}@{}}
\toprule\noalign{}
\begin{minipage}[b]{\linewidth}\raggedright
重力
\end{minipage} & \begin{minipage}[b]{\linewidth}\raggedright
摩擦トルクの上限 μNr̄
\end{minipage} & \begin{minipage}[b]{\linewidth}\raggedright
台車が滑り始める振りの速さ
\end{minipage} \\
\midrule\noalign{}
\endhead
\bottomrule\noalign{}
\endlastfoot
地球 & 155 N·m & 0.85 s より速い \\
火星 & 59 N·m & 1.38 s より速い \\
月 & 26 N·m & 2.09 s より速い \\
\end{longtable}

\textbf{「速く振って、ゆっくり（3 s）戻す」往復を繰り返すと、関節が 1 つでも台車が回る。}

\begin{longtable}[]{@{}
  >{\raggedright\arraybackslash}p{(\columnwidth - 6\tabcolsep) * \real{0.2500}}
  >{\raggedright\arraybackslash}p{(\columnwidth - 6\tabcolsep) * \real{0.2500}}
  >{\raggedright\arraybackslash}p{(\columnwidth - 6\tabcolsep) * \real{0.2500}}
  >{\raggedright\arraybackslash}p{(\columnwidth - 6\tabcolsep) * \real{0.2500}}@{}}
\toprule\noalign{}
\begin{minipage}[b]{\linewidth}\raggedright
振りの時間
\end{minipage} & \begin{minipage}[b]{\linewidth}\raggedright
地球
\end{minipage} & \begin{minipage}[b]{\linewidth}\raggedright
火星
\end{minipage} & \begin{minipage}[b]{\linewidth}\raggedright
月
\end{minipage} \\
\midrule\noalign{}
\endhead
\bottomrule\noalign{}
\endlastfoot
0.25 s & −37°/往復 & −49° & −54° \\
0.47 s & −8° & −31° & −44° \\
0.87 s & 0° & −2° & −19° \\
摩擦ゼロ（宇宙と同じ） & 0°（1.9×10\textsuperscript{-4}°） & & \\
\end{longtable}

0.6 s で振る往復を 5 回繰り返すと、地球では −2.8°、\textbf{月では −177°} 回った。

\textbf{考察: 宇宙と地上で「回る仕組み」が違う}

\begin{longtable}[]{@{}
  >{\raggedright\arraybackslash}p{(\columnwidth - 4\tabcolsep) * \real{0.3333}}
  >{\raggedright\arraybackslash}p{(\columnwidth - 4\tabcolsep) * \real{0.3333}}
  >{\raggedright\arraybackslash}p{(\columnwidth - 4\tabcolsep) * \real{0.3333}}@{}}
\toprule\noalign{}
\begin{minipage}[b]{\linewidth}\raggedright
\end{minipage} & \begin{minipage}[b]{\linewidth}\raggedright
宇宙（Sim01〜05）
\end{minipage} & \begin{minipage}[b]{\linewidth}\raggedright
地上（Sim06）
\end{minipage} \\
\midrule\noalign{}
\endhead
\bottomrule\noalign{}
\endlastfoot
関節 1 つの往復 & 回らない（可逆） & \textbf{回る}（速い振りで滑り、遅い戻りで止まる） \\
正味回転の源 & 形状ループの曲率（幾何学的位相） & 摩擦による \textbf{時間対称性の破れ}（スティックスリップ） \\
速さの影響 & 関係ない（経路だけで決まる） & \textbf{決定的}（速さが摩擦の限界を超えるかどうか） \\
\end{longtable}

\begin{itemize}
\tightlist
\item
  宇宙では「速さは関係ない」が鉄則だった（第~\ref{ch:theory}章）。地上では摩擦が速さに対して非対称なので、\textbf{速さそのものが制御入力} になる。
\item
  \textbf{低重力天体では、腕の作業がそのまま機体の姿勢を乱す。} 月面で荷物を素早く振れば、ローバーが意図せず回る。 逆に、これを利用して車輪や履帯を使わずに向きを変えることもできる（地上版 AMBAC）。
\item
  地球の通常の作業速度（1 s 以上）では、ほぼ問題にならない。
\end{itemize}

\section{フィジカル AI への示唆}

\begin{enumerate}
\def\labelenumi{\arabic{enumi}.}
\tightlist
\item
  \textbf{安価なシミュレーションで、高価な実機の「危ない領域」を事前に地図にできる。} 転倒・滑りの限界は ZMP と摩擦の式から求まり、時間・大きさ・重力についてのスケーリング則（フルード則）まで厳密に得られた。 実機での試行錯誤（＝転倒させて壊すリスク）を大きく減らせる。
\item
  \textbf{学習・最適化に入れるべき制約が明確になる。} ZMP と摩擦の限界は、強化学習の報酬や安全制約として、 そのまま方策に組み込める形をしている。
\item
  \textbf{重力や摩擦が違う環境（月・火星・氷上）への移行は、パラメータの変更だけで評価できる。} これはシミュレーションの大きな強みで、本章の結果は「地球で学習した方策を月にそのまま持っていくと危ない」ことを定量的に示している。
\end{enumerate}

\section{限界と次の課題}

\begin{itemize}
\tightlist
\item
  台車は静止しているとして、転倒・滑りの \textbf{開始条件} だけを調べた。転倒が始まった後の運動（台車が傾く）は扱っていない。
\item
  台車の並進（前後に滑る）と旋回の連成、走行中の加減速は扱っていない。
\item
  地面は剛体と仮定した（不整地・軟弱地盤・レゴリスの沈み込みは扱っていない）。
\item
  アクチュエータは理想的（油圧・電動の動特性は今後の課題、第~\ref{ch:conclusion}章参照）。
\end{itemize}

\textbf{次の候補} - 3 次元化（腕の空間的な動きで、前後の転倒・左右の転倒・旋回が同時に起きる） - 走行中の腕の作業（台車の加減速と腕の反作用の連成） - 学習方策に ZMP 制約を入れ、Part B の最適化結果と比較する

\begin{quote}
\textbf{修正履歴}: 当初 Part B の最適化で、肘の関節の速度・加速度の制約が課されていなかった（制約を間引くときの配列の取り方の誤り）。 修正後、地球での最短時間は 1.00 s（38\% 短縮）から 1.03 s（36\% 短縮）に変わった。月面の結果（2.01 s）は変わらない。
\end{quote}

\Detail{ch06}
\chapter{「足なんて飾り」は本当か（Sim09）}\label{ch:legs}

\section{問い}

「足なんて飾りです」という有名なセリフがある。元は『機動戦士ガンダム』第 42 話で、整備兵がシャアに脚のない機体ジオングについて言う「あんなの飾りです。偉い人にはそれがわからんのですよ」で、それが広まる過程でこの形になった \cite{R2}。 コスト面だけを見れば、脚の代わりに戦車型の台車にした方が安い。では物理的に、脚は本当に飾りなのか。 宇宙（AMBAC）と地上（安定性）の 2 つの場面で検証した。

\section{Part A: 宇宙 --- 脚は AMBAC の道具になるか}

\subsection{A1: 平面モデル、総質量 76 kg を揃えて比較}

1 周あたりに稼げる回転の上限（関節の可動域の中で、同じ符号の曲率の領域をすべて囲んだ場合）を比べた。

\begin{longtable}[]{@{}
  >{\raggedright\arraybackslash}p{(\columnwidth - 4\tabcolsep) * \real{0.3333}}
  >{\raggedright\arraybackslash}p{(\columnwidth - 4\tabcolsep) * \real{0.3333}}
  >{\raggedright\arraybackslash}p{(\columnwidth - 4\tabcolsep) * \real{0.3333}}@{}}
\toprule\noalign{}
\begin{minipage}[b]{\linewidth}\raggedright
構成
\end{minipage} & \begin{minipage}[b]{\linewidth}\raggedright
1 周あたりの上限
\end{minipage} & \begin{minipage}[b]{\linewidth}\raggedright
90° に必要な周回
\end{minipage} \\
\midrule\noalign{}
\endhead
\bottomrule\noalign{}
\endlastfoot
腕＋脚（人間型） & 9.57° & 10 周 \\
腕＋腕（脚を外し、脚の質量は胴体へ） & 5.71° & 16 周 \\
腕だけ（脚を外す） & \textbf{0°}（関節 1 つでは往復で元に戻る） & 不可能 \\
腕＋重い脚（脚 2 倍：推進剤タンクや装甲を脚に置く） & \textbf{12.87°} & 7 周 \\
\end{longtable}

\subsection{A2: 3 次元、股関節（球関節）で脚に円錐を描かせる}

\begin{longtable}[]{@{}
  >{\raggedright\arraybackslash}p{(\columnwidth - 2\tabcolsep) * \real{0.5000}}
  >{\raggedright\arraybackslash}p{(\columnwidth - 2\tabcolsep) * \real{0.5000}}@{}}
\toprule\noalign{}
\begin{minipage}[b]{\linewidth}\raggedright
動作
\end{minipage} & \begin{minipage}[b]{\linewidth}\raggedright
1 周あたりの胴体の回転
\end{minipage} \\
\midrule\noalign{}
\endhead
\bottomrule\noalign{}
\endlastfoot
腕（肩）: 前下方 30° の軸、半頂角 30° & 24.8° \\
\textbf{脚（股関節）: 同じ軸、同じ半頂角} & \textbf{48.2°（腕の 1.9 倍）} \\
脚（股関節）: 同じ軸、半頂角 45° & 96.8° \\
\end{longtable}

股関節の可動域は肩より狭いが、脚は腕の 3 倍重く長いので、同じ動きでもはるかに大きく胴体を回せる。

\textbf{結論 A: 宇宙では、脚は飾りではない。} 脚は AMBAC で最も効く「重い肢」であり、 脚を外すと姿勢変更の能力は半分近くまで落ちる（腕が 1 本なら平面では不可能になる）。 脚に質量（推進剤タンク、装甲）を集めるほど、AMBAC は強くなる。 皮肉なことに、セリフが語られた「宇宙」こそ、脚が最も役に立つ場面だった。 実際に、宇宙飛行士の大西卓哉さんが国際宇宙ステーションで脚を使って姿勢を制御し、AMBAC に言及している \cite{R43,R44}。

\section{Part B: 地上 --- 二本脚で立つ機体 vs 戦車型台車}

\begin{figure}[htbp]\centering
\resizebox{0.9\linewidth}{!}{% 戦車型と二本脚の支持範囲、腰の戦略
\begin{tikzpicture}[>=Stealth,font=\small]
\begin{scope}
  \fill[pattern=north east lines,pattern color=gray] (-1.6,-0.12) rectangle (1.6,0);
  \draw[thick,fill=gray!20,rounded corners=3pt] (-1.1,0) rectangle (1.1,0.45);
  \draw[thick,fill=gray!10] (-0.3,0.45) rectangle (0.3,1.4);
  \fill (0,0.75) circle (1.5pt);
  \draw[<->,red,thick] (-1.1,-0.35)--(1.1,-0.35) node[midway,below,font=\footnotesize]{支持 $\pm0.40$ m};
  \node at (0,2.0) {戦車型};
\end{scope}
\begin{scope}[shift={(4.2,0)}]
  \fill[pattern=north east lines,pattern color=gray] (-1.3,-0.12) rectangle (1.3,0);
  \draw[thick] (-0.15,0)--(0.25,0);
  \draw[line width=3pt,gray] (0,0.05)--(0,1.2);
  \draw[thick,fill=gray!10] (-0.25,1.2) rectangle (0.25,2.1);
  \fill (0,1.3) circle (1.5pt);
  \draw[<->,red,thick] (-0.15,-0.35)--(0.25,-0.35) node[midway,below,font=\footnotesize]{足裏 $\pm0.125$ m};
  \node at (0,2.6) {二本脚（直立）};
\end{scope}
\begin{scope}[shift={(8.2,0)}]
  \fill[pattern=north east lines,pattern color=gray] (-1.3,-0.12) rectangle (1.6,0);
  \draw[thick] (-0.15,0)--(0.25,0);
  \draw[line width=3pt,gray] (0,0.05)--(0,1.2);
  \begin{scope}[shift={(0,1.2)},rotate=25]
    \draw[thick,fill=gray!10] (-0.25,0) rectangle (0.25,0.9);
  \end{scope}
  \draw[line width=2.5pt,blue!60] (0.25,1.95)--(1.2,1.55);
  \fill[orange] (1.2,1.55) circle (3.5pt);
  \draw[->,thick] (-0.25,2.35) arc (100:130:0.6);
  \node[font=\footnotesize,align=center] at (-1.2,1.5) {上体を\\後ろに反らす};
  \node at (0.3,3.4) {二本脚＋腰の戦略};
\end{scope}
\end{tikzpicture}
}
\caption{支持範囲の比較と腰の戦略。二本脚は足裏の支持範囲が狭く、荷物を持つと上体を反らせて ZMP を足裏に戻す必要がある。}\label{fig:tikz-legs}
\end{figure}

Sim06 と同じ荷物運び（10 kg を前方の低い位置から後方の高い位置へ）を、同じ腕・同じ総質量（97 kg）で比べた。 すべての方策に、腕の関節速度 6 rad/s・加速度 60 rad/s²、上体の傾きの速度 3 rad/s・加速度 15 rad/s²、傾き ±40° 以内、 床から浮かないこと（垂直反力 ≥ 自重の 30\%）の制約を課した。

\begin{longtable}[]{@{}
  >{\raggedright\arraybackslash}p{(\columnwidth - 6\tabcolsep) * \real{0.2500}}
  >{\raggedright\arraybackslash}p{(\columnwidth - 6\tabcolsep) * \real{0.2500}}
  >{\raggedright\arraybackslash}p{(\columnwidth - 6\tabcolsep) * \real{0.2500}}
  >{\raggedright\arraybackslash}p{(\columnwidth - 6\tabcolsep) * \real{0.2500}}@{}}
\toprule\noalign{}
\begin{minipage}[b]{\linewidth}\raggedright
機体
\end{minipage} & \begin{minipage}[b]{\linewidth}\raggedright
支持範囲
\end{minipage} & \begin{minipage}[b]{\linewidth}\raggedright
腕だけ最適化
\end{minipage} & \begin{minipage}[b]{\linewidth}\raggedright
腕＋上体の前後傾
\end{minipage} \\
\midrule\noalign{}
\endhead
\bottomrule\noalign{}
\endlastfoot
戦車型（接地長 0.8 m、重心が低い） & ±0.40 m & 0.98 s & \textbf{0.64 s} \\
二本脚（足裏 0.25 m、重心が高い） & ±0.125 m & \textbf{不可能（静止しても転倒）} & 0.69 s \\
\end{longtable}

\FIG{sim09_legs.png}

\textbf{考察} - \textbf{二本脚で立つ機体は、10 kg の荷物を腕の先で持つだけで、静止していても転倒する。} 足裏の支持範囲（±0.125 m）が狭いため、荷物を前に持つと ZMP が足のつま先を越える。 - \textbf{上体を前後に傾ければ（腰の戦略）、二本脚でも運べる。} 荷物を前に持つときは後ろに反り（−40°）、 後ろに置くときは前に傾く（+28°）。重い上体（56 kg）を動かして腕の反作用を打ち消す、\textbf{地上版の AMBAC} である（人の姿勢制御でいう腰の戦略 \cite{R31}）。 人が重い物を持つときに自然に体を反らせるのと同じ動きを、最適化が再発見した。 - \textbf{ただし、上体を傾けられるなら戦車型の方がさらに速い}（0.64 s 対 0.69 s）。 どちらも最終的には関節の速度・加速度の上限で決まっており、広い支持範囲は余裕として効いている。 - \textbf{地上で安定性を決めているのは「脚」ではなく、「動かせる重い上体（腰）」と「広い支持範囲」である。} 平地での作業に限れば、脚は安定性を損なうだけで、戦車型＋腰の関節が最も合理的になる。

\section{コスト面}

\begin{longtable}[]{@{}
  >{\raggedright\arraybackslash}p{(\columnwidth - 4\tabcolsep) * \real{0.3333}}
  >{\raggedright\arraybackslash}p{(\columnwidth - 4\tabcolsep) * \real{0.3333}}
  >{\raggedright\arraybackslash}p{(\columnwidth - 4\tabcolsep) * \real{0.3333}}@{}}
\toprule\noalign{}
\begin{minipage}[b]{\linewidth}\raggedright
\end{minipage} & \begin{minipage}[b]{\linewidth}\raggedright
戦車型＋腰
\end{minipage} & \begin{minipage}[b]{\linewidth}\raggedright
二本脚（人間型）
\end{minipage} \\
\midrule\noalign{}
\endhead
\bottomrule\noalign{}
\endlastfoot
移動・姿勢に使う関節・モータ & 履帯 2 ＋ 腰 1〜2 & 片脚 6 自由度 × 2 ＝ 12 ＋ 腰 \\
立っているだけで必要な制御 & 不要（静的に安定） & 常にバランス制御が必要 \\
荷物を持つ作業（平地） & 腰を使わなくても可能 & 腰の戦略が必須 \\
段差・階段・不整地 & 苦手（本章では評価していない） & 得意（本章では評価していない） \\
\end{longtable}

\section{判定}

\begin{longtable}[]{@{}
  >{\raggedright\arraybackslash}p{(\columnwidth - 4\tabcolsep) * \real{0.3333}}
  >{\raggedright\arraybackslash}p{(\columnwidth - 4\tabcolsep) * \real{0.3333}}
  >{\raggedright\arraybackslash}p{(\columnwidth - 4\tabcolsep) * \real{0.3333}}@{}}
\toprule\noalign{}
\begin{minipage}[b]{\linewidth}\raggedright
場面
\end{minipage} & \begin{minipage}[b]{\linewidth}\raggedright
「足は飾り」か
\end{minipage} & \begin{minipage}[b]{\linewidth}\raggedright
理由
\end{minipage} \\
\midrule\noalign{}
\endhead
\bottomrule\noalign{}
\endlastfoot
\textbf{宇宙} & \textbf{飾りではない} & 脚は AMBAC で最も効く重い肢。脚を外すと姿勢変更能力が半減する \\
\textbf{地上・平地の作業} & \textbf{ほぼ飾り（むしろ不利）} & 支持範囲が狭く、静止して荷物を持つことすら難しい。戦車型＋腰の方が速く、安く、安定 \\
\textbf{地上・不整地や段差} & 未検証 & 脚の本来の強みはここにある（次の課題） \\
\end{longtable}

人体ベースのフィジカル AI として見ると、\textbf{人の脚は「歩くため」と「重い肢として姿勢を操るため」の 2 つの役割を持つ}。 平地での作業だけなら戦車型で十分だが、人間らしい動き（腰の戦略、宇宙での姿勢変更）には、 重い下半身を能動的に動かせることが本質的に効いている。

\section{注釈：判定には「用途」と「前提」が付く}

\subsection{注釈 1: 推進剤を惜しまなければ、宇宙でも脚は飾りになる}

Sim04・Sim05 の 18 m 級の数値で比べると、次のようになる。

\begin{longtable}[]{@{}
  >{\raggedright\arraybackslash}p{(\columnwidth - 4\tabcolsep) * \real{0.3333}}
  >{\raggedright\arraybackslash}p{(\columnwidth - 4\tabcolsep) * \real{0.3333}}
  >{\raggedright\arraybackslash}p{(\columnwidth - 4\tabcolsep) * \real{0.3333}}@{}}
\toprule\noalign{}
\begin{minipage}[b]{\linewidth}\raggedright
90° 旋回
\end{minipage} & \begin{minipage}[b]{\linewidth}\raggedright
時間
\end{minipage} & \begin{minipage}[b]{\linewidth}\raggedright
推進剤
\end{minipage} \\
\midrule\noalign{}
\endhead
\bottomrule\noalign{}
\endlastfoot
スラスタ & 5 s & 88 kg \\
スラスタ & 10 s & 44 kg \\
AMBAC（脚と腕） & 約 130 s & 0 kg \\
\end{longtable}

採算を度外視して推進剤を使えば、機動力はスラスタが圧倒的に上である。 宇宙で脚が「飾りではない」のは、\textbf{推進剤に限りがある・長期間運用する・推進剤が尽きた・排気を出せない} といった前提の下での話である。 短時間の高機動戦闘のように推進剤を使い切ってよい用途では、脚は飾りという判定にも理がある。

\subsection{注釈 2: 平地の戦車型には、方向転換の課題が残る}

履帯（キャタピラ）の機体は、左右の履帯の速度差で向きを変える（スキッドステア）。 このとき履帯は地面の上を横に滑る必要があり、Sim06 Part C で求めた旋回の摩擦抵抗（97 kg の機体で約 155 N·m）に打ち勝たなければならない。 そのため、旋回には大きなエネルギーが必要で、路面も傷める。旋回の多い用途では、車輪（全方向移動）や脚の方が有利になりうる。

\subsection{注釈 3: 段差・不整地では、汎用性に勝る脚が有利}

段差・階段・瓦礫では、接地点を選べる脚が有利になることは容易に想像できる（本章では未評価）。

\subsection{用途別の整理}

\begin{longtable}[]{@{}
  >{\raggedright\arraybackslash}p{(\columnwidth - 4\tabcolsep) * \real{0.3333}}
  >{\raggedright\arraybackslash}p{(\columnwidth - 4\tabcolsep) * \real{0.3333}}
  >{\raggedright\arraybackslash}p{(\columnwidth - 4\tabcolsep) * \real{0.3333}}@{}}
\toprule\noalign{}
\begin{minipage}[b]{\linewidth}\raggedright
用途
\end{minipage} & \begin{minipage}[b]{\linewidth}\raggedright
有利な下半身
\end{minipage} & \begin{minipage}[b]{\linewidth}\raggedright
根拠
\end{minipage} \\
\midrule\noalign{}
\endhead
\bottomrule\noalign{}
\endlastfoot
宇宙・長期運用・推進剤に制約あり & \textbf{重い脚}（AMBAC 用） & Sim03〜05, 09 \\
宇宙・短時間の高機動（推進剤を使い切ってよい） & \textbf{推進剤}（脚は飾り） & Sim04, 05 \\
地上・平地での作業 & \textbf{戦車型＋腰} & Sim06, 09 \\
地上・平地で旋回が多い & 車輪（全方向移動）など & 注釈 2 \\
地上・段差・不整地 & \textbf{脚} & 注釈 3（未評価） \\
\end{longtable}

\textbf{「足は飾りか」の答えは用途で決まる。} 用途に応じて下半身（脚部）の形式を選ぶという設計思想は、 ロボットを題材にしたゲームなどで脚部の種類を選べる仕組みとも通じる。 人体ベースのフィジカル AI でも、すべてを人の形にするのではなく、用途に応じて「人の形をどこまで残すか」を決めることになる。

\section{限界と次の課題}

\begin{itemize}
\tightlist
\item
  二本脚は足をそろえた静止立位のみ。足を前後に開く・踏み出す（ステップ戦略）は扱っていない。人は踏み出すことで支持範囲を広げる。
\item
  上体の傾きは関節の速度・加速度の制限で代用し、股関節のトルクは直接制限していない。
\item
  段差・階段・不整地での移動能力（脚の本来の強み）は評価していない。
\item
  宇宙での脚の重さの増加は、関節トルクの増加（Sim04）とのトレードオフを含めていない。
\end{itemize}

\Detail{ch09}
\chapter{段差の乗り越え：戦車型 vs 二本脚（Sim10）}\label{ch:steps}

\section{問い}

Sim09 で「段差や不整地では脚が有利だろう」という注釈を付けた。これを確かめる。 ゆっくり動く（準静的な）モデルで、越えられる段差の高さの上限を求めた。走り込みや蹴り出しの勢いは使わない。

\section{Part A: 戦車型（Sim06 の機体、履帯の接地長 0.8 m）}

静力学の式は付録~\ref{app:steps} を参照。

\begin{figure}[htbp]\centering
\resizebox{0.95\linewidth}{!}{% 段差の乗り越え
\begin{tikzpicture}[>=Stealth,font=\small]
\begin{scope}
  \fill[pattern=north east lines,pattern color=gray] (-0.5,-0.12) rectangle (2.4,0);
  \fill[pattern=north east lines,pattern color=gray] (2.4,-0.12) rectangle (4.0,0.8);
  \draw[thick] (-0.5,0)--(2.4,0)--(2.4,0.8)--(4.0,0.8);
  \begin{scope}[rotate around={18.4:(0,0)}]
    \draw[thick,fill=gray!20,rounded corners=3pt] (0,0) rectangle (2.6,0.45);
    \draw[thick,fill=gray!10] (0.95,0.45) rectangle (1.45,1.2);
    \fill (1.2,0.7) circle (1.5pt);
  \end{scope}
  \fill[red] (0,0) circle (2pt) node[below left,font=\footnotesize]{R};
  \fill[red] (2.4,0.8) circle (2pt) node[above left,font=\footnotesize]{E};
  \draw[<->] (4.2,0)--(4.2,0.8) node[midway,right]{$h$};
  \draw[->,red,thick] (0,0)--(0,0.7) node[left,font=\footnotesize]{$N_1$};
  \draw[->,red,thick] (2.4,0.8)--++(108.4:0.7) node[left,font=\footnotesize]{$N_2$};
  \node[align=center,font=\footnotesize] at (1.8,-0.75) {戦車型：後ろの接地点 R と段の角 E で支え、\\摩擦の範囲内で這い上がる};
\end{scope}
\begin{scope}[shift={(7.2,0)}]
  \fill[pattern=north east lines,pattern color=gray] (-1.2,-0.12) rectangle (0.2,0);
  \fill[pattern=north east lines,pattern color=gray] (0.2,-0.12) rectangle (1.6,0.8);
  \draw[thick] (-1.2,0)--(0.2,0)--(0.2,0.8)--(1.6,0.8);
  \coordinate (Ank) at (0.8,0.8); \coordinate (Hip) at (0.8,2.1); \coordinate (Kn) at (1.35,1.45);
  \draw[line width=3pt,gray] (Ank)--(Kn)--(Hip);
  \draw[dashed] (Ank)--(Hip);
  \draw[<->] (0.8,1.45)--(Kn) node[midway,below,font=\footnotesize]{$d$};
  \draw[thick,fill=gray!10,rotate around={-20:(Hip)}] ($(Hip)+(-0.22,0)$) rectangle ($(Hip)+(0.22,0.9)$);
  \draw[line width=3pt,gray!50] (Hip)--(-0.3,1.1)--(-0.4,0.05);
  \draw[<->] (1.85,0)--(1.85,0.8) node[midway,right]{$h$};
  \node[align=center,font=\footnotesize] at (0.3,-0.75) {二本脚：膝のはみ出し $d$ が\\膝トルク $mg(d-x_c)$ を決める};
\end{scope}
\end{tikzpicture}
}
\caption{段差の乗り越えの静力学。左：戦車型（2 点支持と摩擦）。右：二本脚（膝のはみ出し $d$ と膝トルク）。}\label{fig:tikz-steps}
\end{figure}


越えられる段差は、次の 2 つの限界のうち小さい方で決まる。

\begin{enumerate}
\def\labelenumi{\arabic{enumi}.}
\tightlist
\item
  \textbf{形の限界}: 履帯の前端（持ち上がった部分）の高さ（0.25 m とした）より高い段差には、前端が引っかかる。
\item
  \textbf{力の限界}: 前端を段の角に乗せて車体が傾いた状態で、後ろの接地点と段の角の 2 点で支えながら、 履帯の摩擦（μ = 0.6）の範囲内で這い上がれるか。重心が後ろの接地点より後ろに来ると、後ろに転倒する。
\end{enumerate}

\begin{longtable}[]{@{}
  >{\raggedright\arraybackslash}p{(\columnwidth - 6\tabcolsep) * \real{0.2500}}
  >{\raggedright\arraybackslash}p{(\columnwidth - 6\tabcolsep) * \real{0.2500}}
  >{\raggedright\arraybackslash}p{(\columnwidth - 6\tabcolsep) * \real{0.2500}}
  >{\raggedright\arraybackslash}p{(\columnwidth - 6\tabcolsep) * \real{0.2500}}@{}}
\toprule\noalign{}
\begin{minipage}[b]{\linewidth}\raggedright
腕の姿勢（重心の位置）
\end{minipage} & \begin{minipage}[b]{\linewidth}\raggedright
力の限界
\end{minipage} & \begin{minipage}[b]{\linewidth}\raggedright
越えられる段差
\end{minipage} & \begin{minipage}[b]{\linewidth}\raggedright
20 cm の段差に必要な摩擦係数
\end{minipage} \\
\midrule\noalign{}
\endhead
\bottomrule\noalign{}
\endlastfoot
腕を後ろ上方にたたむ（重心が後ろ） & 0.26 m & 0.25 m & 0.16 \\
腕を下ろす（中立） & 0.30 m & 0.25 m & 0.23 \\
腕を前に伸ばす（重心が前） & \textbf{0.33 m} & 0.25 m & 0.24 \\
\end{longtable}

\begin{itemize}
\tightlist
\item
  今回の形（前端の高さ 0.25 m）では、\textbf{形の限界が先に効く}。どの腕の姿勢でも越えられる段差は 0.25 m である。
\item
  \textbf{腕を前に伸ばして重心を前に移すと、力の限界は 0.26 m → 0.33 m に上がる}（後ろへの転倒が起きにくくなる）。 前端を高くする（フリッパー付きの履帯など）と、腕による重心移動がそのまま乗り越え能力の差になる。地上版の能動的質量バランスである。
\item
  重心を後ろに置くと、必要な摩擦は減るが、後ろに転倒しやすくなる。トレードオフがある。
\end{itemize}

\section{Part B: 二本脚（Sim09 の機体、股関節の高さ 0.9 m、大腿・下腿 各 0.45 m）}

\begin{enumerate}
\def\labelenumi{\arabic{enumi}.}
\tightlist
\item
  \textbf{可動域の限界}: 振り出す脚の足を段の上（高さ h ＋ すき間 5 cm、前方 0.30 m）に置けるか （股関節の屈曲 ≤ 120°、膝の屈曲 ≤ 140°）。
\item
  \textbf{トルクの限界}: 段の上の脚で体を持ち上げ始める瞬間に、膝トルクと股関節トルクがともに上限（各 250 N·m）以下か。
\end{enumerate}

\begin{longtable}[]{@{}
  >{\raggedright\arraybackslash}p{(\columnwidth - 2\tabcolsep) * \real{0.5000}}
  >{\raggedright\arraybackslash}p{(\columnwidth - 2\tabcolsep) * \real{0.5000}}@{}}
\toprule\noalign{}
\begin{minipage}[b]{\linewidth}\raggedright
項目
\end{minipage} & \begin{minipage}[b]{\linewidth}\raggedright
越えられる段差
\end{minipage} \\
\midrule\noalign{}
\endhead
\bottomrule\noalign{}
\endlastfoot
可動域の限界 & 0.63 m \\
トルクの限界（上体を立てたまま） & 0.18 m \\
\textbf{トルクの限界（上体を前に傾けて重心を前へ移す）} & \textbf{0.60 m} \\
\end{longtable}

\begin{itemize}
\tightlist
\item
  上体を立てたままだと、膝だけで体を持ち上げることになり、18 cm 程度が限界になる。
\item
  \textbf{上体を前に傾けて重心を膝の上に寄せると、膝の負担を股関節に分けられ、60 cm まで上れる。} 人が階段や高い段差を上るときに前かがみになるのと同じで、ここでも重い上体を動かすこと（腰の戦略）が効いている。
\item
  人間大では、\textbf{二本脚（0.60 m）は戦車型（0.25 m）の 2 倍以上の段差を越えられる。} Sim09 の注釈どおり、段差では脚が有利。
\end{itemize}

\section{Part C: 大きさを変えると逆転する}

長さ s 倍・質量 s³ 倍に相似拡大し、アクチュエータのトルク上限は s³ 倍（Sim04 の相似則）とした。

\begin{longtable}[]{@{}
  >{\raggedright\arraybackslash}p{(\columnwidth - 6\tabcolsep) * \real{0.2500}}
  >{\raggedright\arraybackslash}p{(\columnwidth - 6\tabcolsep) * \real{0.2500}}
  >{\raggedright\arraybackslash}p{(\columnwidth - 6\tabcolsep) * \real{0.2500}}
  >{\raggedright\arraybackslash}p{(\columnwidth - 6\tabcolsep) * \real{0.2500}}@{}}
\toprule\noalign{}
\begin{minipage}[b]{\linewidth}\raggedright
s（全高）
\end{minipage} & \begin{minipage}[b]{\linewidth}\raggedright
戦車型
\end{minipage} & \begin{minipage}[b]{\linewidth}\raggedright
二本脚
\end{minipage} & \begin{minipage}[b]{\linewidth}\raggedright
有利な方
\end{minipage} \\
\midrule\noalign{}
\endhead
\bottomrule\noalign{}
\endlastfoot
0.5（約 0.9 m） & 0.12 m（大きさ比 0.25） & 0.32 m（0.63） & 脚 \\
\textbf{1（約 1.8 m、人間大）} & 0.25 m（0.25） & \textbf{0.60 m（0.60）} & \textbf{脚} \\
2（約 3.6 m） & \textbf{0.50 m（0.25）} & 0.37 m（0.18） & 戦車 \\
5（約 9 m） & \textbf{1.25 m（0.25）} & 0.12 m（0.02） & 戦車 \\
\textbf{10（約 18 m 級）} & \textbf{2.50 m（0.25）} & \textbf{0.05 m（0.003）} & \textbf{戦車} \\
\end{longtable}

\FIG{sim10_steps.png}

\begin{itemize}
\tightlist
\item
  \textbf{戦車型は大きさによらず、機体の大きさの 1/4 の段差を越えられる}（形と摩擦で決まり、どちらもスケール不変）。
\item
  \textbf{二本脚は大型化すると急速に段差を越えられなくなる。} 体を持ち上げるのに必要なトルクは s\textsuperscript{4} で増えるが、 使えるトルクは s³ でしか増えないため（二乗三乗則）。人間の 2 倍（約 3.6 m）で戦車型に逆転される。
\item
  18 m 級の二本脚は、相似則のアクチュエータでは 5 cm の段差しか越えられない。 巨大ロボットが脚で段差を越えるには、Sim04 の旋回と同じく、\textbf{相似則をはるかに超えるトルク密度のアクチュエータ} が必要になる。
\end{itemize}

\section{「足は飾り」の判定の更新}

\begin{longtable}[]{@{}
  >{\raggedright\arraybackslash}p{(\columnwidth - 4\tabcolsep) * \real{0.3333}}
  >{\raggedright\arraybackslash}p{(\columnwidth - 4\tabcolsep) * \real{0.3333}}
  >{\raggedright\arraybackslash}p{(\columnwidth - 4\tabcolsep) * \real{0.3333}}@{}}
\toprule\noalign{}
\begin{minipage}[b]{\linewidth}\raggedright
場面
\end{minipage} & \begin{minipage}[b]{\linewidth}\raggedright
人間大
\end{minipage} & \begin{minipage}[b]{\linewidth}\raggedright
18 m 級
\end{minipage} \\
\midrule\noalign{}
\endhead
\bottomrule\noalign{}
\endlastfoot
宇宙（AMBAC） & 脚は有用 & 脚は有用（ただし遅い：Sim04） \\
地上・平地の作業 & 戦車型＋腰が有利 & 戦車型＋腰が有利 \\
\textbf{地上・段差} & \textbf{脚が有利（2 倍以上）} & \textbf{戦車型が圧倒的に有利} \\
\end{longtable}

\textbf{脚の利点は「人間大」で最も大きく、大型化すると段差でも失われる。} 人間の体が脚を持っているのは、人間の大きさだからこそ、と言える。

\section{限界}

\begin{itemize}
\tightlist
\item
  準静的なモデル。勢いを使った乗り越え（走り込み、蹴り出し、ジャンプ）は扱っていない。動的には、どちらもより高い段差を越えられる。
\item
  戦車型の前端の高さ（0.25 m）は仮定。フリッパー付きの履帯なら形の限界は上がる。
\item
  二本脚の股関節トルクのモデルは単純化している（重心を前へ移す分の上体の前傾だけを考えた）。手すりや手をつくことは扱っていない。
\item
  段の角での履帯の変形、足の滑り、膝の摩擦は扱っていない。
\end{itemize}

\Detail{ch10}

\part{上半身：肩と肩甲骨}
\chapter{3 次元化：球面座標の肩と SO(3)（Sim07）}\label{ch:3d}

\section{座標の選び方}

3 次元の角運動量と特異点の式は付録~\ref{app:3d} を参照。

\begin{figure}[htbp]\centering
\resizebox{0.6\linewidth}{!}{% 球面座標の肩と円錐運動
\tdplotsetmaincoords{65}{120}
\begin{tikzpicture}[>=Stealth,font=\small,tdplot_main_coords,scale=2.4]
\draw[->] (0,0,0)--(1.3,0,0) node[anchor=north east]{$x$（前）};
\draw[->] (0,0,0)--(0,1.2,0) node[anchor=north west]{$y$（左）};
\draw[->] (0,0,0)--(0,0,1.2) node[anchor=south]{$z$（上）};
\pgfmathsetmacro{\psA}{40}\pgfmathsetmacro{\alA}{30}
\coordinate (U) at ({cos(\alA)*cos(\psA)},{cos(\alA)*sin(\psA)},{sin(\alA)});
\coordinate (P) at ({cos(\alA)*cos(\psA)},{cos(\alA)*sin(\psA)},0);
\draw[line width=2.5pt,blue!60] (0,0,0)--(U);
\draw[dashed] (0,0,0)--(P)--(U);
\tdplotdrawarc[->]{(0,0,0)}{0.45}{0}{\psA}{anchor=north}{$\psi$}
\node at (0.62,0.55,0.25) {$\alpha$};
\draw[red,thick] plot[domain=0:360,samples=72,variable=\t] ({cos(30)},{sin(30)*cos(\t)},{sin(30)*sin(\t)});
\draw[red,->,thick] (0,0,0)--(1.25,0,0);
\node[red,font=\footnotesize,anchor=east] at (0.9,-0.6,0.5) {円錐（軸：前方）};
\fill (0,0,0) circle (0.8pt) node[anchor=north east,font=\footnotesize]{肩};
\node[blue!70!black,font=\footnotesize,anchor=south west] at (U) {腕 $\bm u(\psi,\alpha)$};
\fill[gray] (0,0,1) circle (0.6pt) node[anchor=west,font=\footnotesize]{極（特異点）};
\end{tikzpicture}
}
\caption{肩を原点とする球面座標 $(\psi,\alpha)$ と、前方を軸とする円錐運動。真上（極）は座標の特異点になる。}\label{fig:tikz-sph}
\end{figure}


\begin{longtable}[]{@{}
  >{\raggedright\arraybackslash}p{(\columnwidth - 4\tabcolsep) * \real{0.3333}}
  >{\raggedright\arraybackslash}p{(\columnwidth - 4\tabcolsep) * \real{0.3333}}
  >{\raggedright\arraybackslash}p{(\columnwidth - 4\tabcolsep) * \real{0.3333}}@{}}
\toprule\noalign{}
\begin{minipage}[b]{\linewidth}\raggedright
対象
\end{minipage} & \begin{minipage}[b]{\linewidth}\raggedright
座標
\end{minipage} & \begin{minipage}[b]{\linewidth}\raggedright
理由
\end{minipage} \\
\midrule\noalign{}
\endhead
\bottomrule\noalign{}
\endlastfoot
腕の向き & \textbf{肩を原点とする球面座標} (ψ: 方位角, α: 仰角)、u = (cos α cos ψ, cos α sin ψ, sin α) & 2 自由度の肩（旋回＋俯仰）の関節角そのもの。旋回動作を自然に書ける \\
胴体の姿勢 & \textbf{回転行列（SO(3)）} & 3 次元の回転は順序で結果が変わる（非可換）。オイラー角は特異点を持つため、積分には使わない \\
\end{longtable}

\begin{itemize}
\tightlist
\item
  \textbf{球面座標の弱点}: α = ±90°（真上・真下）で ψ が定まらない（座標の特異点）。 これは座標の都合だけでなく、2 軸の肩（ジンバル）が実際に持つ \textbf{ジンバルロック} と同じ現象である（Part A3）。
\item
  \textbf{モデルの検証}: 腕を x-z 平面内（ψ = 0）だけで動かすと、平面モデル（Sim01）と 5 桁一致した。
\item
  \textbf{腕のモデル化}: 腕（一様な棒）は 40 個の点質量で表し、角運動量と反力を点質量の和として計算した。
\end{itemize}

3 次元での AMBAC の基本式は、平面の式（第~\ref{ch:theory}章）がベクトル・行列になったものである。

\[
\mathbf L = A(q)\,\boldsymbol\omega + B(q)\,\dot q = 0 \quad\Rightarrow\quad
\boldsymbol\omega = -A(q)^{-1}B(q)\,\dot q,\qquad \dot R = R\,[\boldsymbol\omega]_\times
\]

\begin{itemize}
\tightlist
\item
  A(q): 3×3 の慣性テンソル（形状 q で固まった系全体）
\item
  B(q): 3×2 の結合行列
\item
  ω: 胴体座標で表した胴体の角速度
\end{itemize}

\section{Part A: 宇宙での 3 次元 AMBAC（人間大）}

\FIG{sim07_space.png}

\subsection{A1: 腕で円錐を描く（宇宙飛行士の「腕回し」）}

\begin{longtable}[]{@{}
  >{\raggedright\arraybackslash}p{(\columnwidth - 4\tabcolsep) * \real{0.3333}}
  >{\raggedright\arraybackslash}p{(\columnwidth - 4\tabcolsep) * \real{0.3333}}
  >{\raggedright\arraybackslash}p{(\columnwidth - 4\tabcolsep) * \real{0.3333}}@{}}
\toprule\noalign{}
\begin{minipage}[b]{\linewidth}\raggedright
半頂角 β
\end{minipage} & \begin{minipage}[b]{\linewidth}\raggedright
1 周あたりの胴体回転
\end{minipage} & \begin{minipage}[b]{\linewidth}\raggedright
回転軸
\end{minipage} \\
\midrule\noalign{}
\endhead
\bottomrule\noalign{}
\endlastfoot
10° & 0.65° & ほぼ −x（円錐の軸と逆向き） \\
30° & 5.7° & 〃 \\
45° & 12.4° & 〃 \\
70° & 25.4° & 〃 \\
\end{longtable}

\begin{itemize}
\tightlist
\item
  \textbf{2 自由度の肩 1 つ（腕 1 本）だけで、任意の角度まで回れる。} 平面では関節 2 つ（腕＋脚）が必要だったが、 3 次元の肩はそれだけで 2 自由度の形状空間を持つ。
\item
  小さい円錐では、回転角 ∝ β\^{}2.00（理論どおりの 2 次効果）。Sim05 で議論した「幾何学的位相は振幅の 2 次効果」が、3 次元でも厳密に成り立つ。
\item
  円錐を逆回りにすると、回転は正確に符号反転する。
\end{itemize}

\subsection{A2: 円錐の軸の向き（β = 45°）}

\begin{longtable}[]{@{}
  >{\raggedright\arraybackslash}p{(\columnwidth - 4\tabcolsep) * \real{0.3333}}
  >{\raggedright\arraybackslash}p{(\columnwidth - 4\tabcolsep) * \real{0.3333}}
  >{\raggedright\arraybackslash}p{(\columnwidth - 4\tabcolsep) * \real{0.3333}}@{}}
\toprule\noalign{}
\begin{minipage}[b]{\linewidth}\raggedright
円錐の軸
\end{minipage} & \begin{minipage}[b]{\linewidth}\raggedright
1 周あたり
\end{minipage} & \begin{minipage}[b]{\linewidth}\raggedright
胴体の回転軸と円錐の軸のずれ
\end{minipage} \\
\midrule\noalign{}
\endhead
\bottomrule\noalign{}
\endlastfoot
前方 & 12.4° & 17° \\
右横 & 11.1° & 20° \\
真上（方位角を 1 周させる旋回） & 58.0° & 8° \\
前上方 45° & 39.7° & 31° \\
\end{longtable}

\begin{itemize}
\tightlist
\item
  胴体はおおむね円錐の軸の逆向きに回るが、\textbf{ぴったり逆向きではない}（8〜31° ずれる）。 肩が胴体の中心からずれていること、胴体の慣性テンソルが非等方であることによる。
\item
  「真上」軸の円錐は、方位角 ψ を 360° 回し続ける動作で、腕が胴体のまわりを 1 周する。 胴体の主軸（縦軸）まわりの慣性が最も小さいので、回転量が最も大きい。
\end{itemize}

\subsection{A3: 球面座標の特異点（ジンバルロック）}

円錐が極（真上）の近くを通ると、腕の先端は一定の速さで動いているのに、方位角の関節速度 ψ' が発散する。

\begin{longtable}[]{@{}
  >{\raggedright\arraybackslash}p{(\columnwidth - 4\tabcolsep) * \real{0.3333}}
  >{\raggedright\arraybackslash}p{(\columnwidth - 4\tabcolsep) * \real{0.3333}}
  >{\raggedright\arraybackslash}p{(\columnwidth - 4\tabcolsep) * \real{0.3333}}@{}}
\toprule\noalign{}
\begin{minipage}[b]{\linewidth}\raggedright
ループと極との最小角距離
\end{minipage} & \begin{minipage}[b]{\linewidth}\raggedright
最大 ψ'
\end{minipage} & \begin{minipage}[b]{\linewidth}\raggedright
最大 α'
\end{minipage} \\
\midrule\noalign{}
\endhead
\bottomrule\noalign{}
\endlastfoot
30° & 3.1 rad/s & 1.57 rad/s \\
3° & 30 rad/s & 1.57 rad/s \\
0.3° & 300 rad/s & 1.57 rad/s \\
\end{longtable}

ψ' は極からの距離に反比例する。\textbf{球面座標は旋回動作には便利だが、極の近くでは使えない。} 対策は次のとおり。 - 計画は腕の向きのベクトル u（または四元数）で行い、関節角への変換は極から離れた領域だけで行う。 - 肩の軸の取り付け方向を工夫し、極が作業領域の外に来るようにする（実機のジンバル設計と同じ考え方）。 - 冗長な関節（3 自由度の肩など）を持たせる。

\subsection{A4: 回転の非可換性}

ヨー用の円錐 1 周（約 58°）とピッチ用の円錐 3 周（約 33°）を、順序を入れ替えて実行した。

\begin{longtable}[]{@{}
  >{\raggedright\arraybackslash}p{(\columnwidth - 2\tabcolsep) * \real{0.5000}}
  >{\raggedright\arraybackslash}p{(\columnwidth - 2\tabcolsep) * \real{0.5000}}@{}}
\toprule\noalign{}
\begin{minipage}[b]{\linewidth}\raggedright
順序
\end{minipage} & \begin{minipage}[b]{\linewidth}\raggedright
最終姿勢（回転ベクトル）
\end{minipage} \\
\midrule\noalign{}
\endhead
\bottomrule\noalign{}
\endlastfoot
ヨー → ピッチ & (8.2°, 41.4°, −64.1°) \\
ピッチ → ヨー & (−22.4°, 33.7°, −65.3°) \\
\textbf{2 つの最終姿勢の差} & \textbf{29.4°} \\
\end{longtable}

平面ではループの順序を入れ替えても結果は同じだった（回転が 1 軸だけで可換）。 \textbf{3 次元では動作の順序そのものが結果を変えるので、姿勢変更の計画は「どの動作をどの順で行うか」という系列の問題になる。} これは学習・計画（フィジカル AI）にとって、平面より本質的に難しい点である。

\section{Part B: 地上での 3 次元（戦車型台車＋球面座標の腕）}

\FIG{sim07_ground.png}

\subsection{B1: 静的な定格荷重マップ}

腕の向き（ψ, α）ごとに、静止状態で転倒しない最大荷重を求めた（安全余裕 5 cm）。 クレーンや油圧ショベルの \textbf{定格荷重表} に相当する。

\begin{longtable}[]{@{}
  >{\raggedright\arraybackslash}p{(\columnwidth - 2\tabcolsep) * \real{0.5000}}
  >{\raggedright\arraybackslash}p{(\columnwidth - 2\tabcolsep) * \real{0.5000}}@{}}
\toprule\noalign{}
\begin{minipage}[b]{\linewidth}\raggedright
腕の向き（水平）
\end{minipage} & \begin{minipage}[b]{\linewidth}\raggedright
最大荷重
\end{minipage} \\
\midrule\noalign{}
\endhead
\bottomrule\noalign{}
\endlastfoot
前（ψ = 0°） & 41.5 kg \\
斜め前（45°） & 42.2 kg \\
\textbf{横（90°）} & \textbf{24.3 kg} \\
後（180°） & 41.5 kg \\
\end{longtable}

\begin{itemize}
\tightlist
\item
  接地面が前後に長く左右に狭いので、\textbf{横向きが最も弱い}（前後の約 6 割）。
\item
  腕を高く上げる（α ≳ 70°）と荷物が接地面の真上に来て、静的には転倒しなくなる（図上部の空白）。
\end{itemize}

\subsection{B2: 腕を水平に 180° 旋回（前 → 左 → 後、荷物 10 kg）}

すべての方策に同じ制約（転倒・滑り・関節速度 6 rad/s・加速度 40 rad/s²・仰角 ≥ −35°）を課し、最短時間を比べた。

\begin{longtable}[]{@{}
  >{\raggedright\arraybackslash}p{(\columnwidth - 4\tabcolsep) * \real{0.3333}}
  >{\raggedright\arraybackslash}p{(\columnwidth - 4\tabcolsep) * \real{0.3333}}
  >{\raggedright\arraybackslash}p{(\columnwidth - 4\tabcolsep) * \real{0.3333}}@{}}
\toprule\noalign{}
\begin{minipage}[b]{\linewidth}\raggedright
方策
\end{minipage} & \begin{minipage}[b]{\linewidth}\raggedright
地球
\end{minipage} & \begin{minipage}[b]{\linewidth}\raggedright
月
\end{minipage} \\
\midrule\noalign{}
\endhead
\bottomrule\noalign{}
\endlastfoot
水平のまま旋回 & 1.83 s & 4.48 s \\
途中で 60° 持ち上げる & 1.24 s & 3.04 s \\
途中で 35° 下げる & 1.01 s & 2.48 s \\
\textbf{最適化（仰角の時間変化と旋回のタイミング）} & \textbf{0.88 s（52\% 短縮）} & \textbf{2.07 s（54\% 短縮）} \\
\end{longtable}

\textbf{考察} - 水平のまま速く旋回すると、遠心力で ZMP が外側に大きく振られ、横方向（最も狭い方向）で転倒する（中央図の破線）。 - 腕を上げても下げても旋回半径（cos α）が小さくなり、慣性と遠心力が減る。ただし \textbf{下げる方が有利} である。 転倒させるモーメントは「高さ × 水平方向の加速度」なので、荷物を低くするとその分だけ効きが弱まる。 - 最適化は、仰角の下限（−35°、地面とのすき間）まで腕を下げ、旋回の中盤（横向き＝最も弱い方向）で最も低くする解を選んだ。 Sim06 の「荷物を体に引き寄せる」と同じく、\textbf{慣性と転倒モーメントを減らす動き} であり、人間やクレーン操作員が経験的に行う動作と一致する。

\section{フィジカル AI への示唆}

\begin{enumerate}
\def\labelenumi{\arabic{enumi}.}
\tightlist
\item
  \textbf{座標の選択はモデル化の一部である。} 球面座標は人間にとって直感的で、旋回動作の指令に向く。 一方で特異点を持ち、学習方策の出力をそのまま関節角にすると、極の近くで関節速度が発散する。 \textbf{方策が何を出力するか（関節角か、手先の向きか）は、学習の安定性と実機の安全性に直結する設計判断} である。
\item
  \textbf{3 次元の姿勢は非可換なので、計画は順序の問題になる。} 平面の結果（Sim01〜05）をそのまま 3 次元に拡張するだけでは不十分で、 系列としての計画・学習が必要になる。
\item
  \textbf{定格荷重マップのような「安全な作業範囲の地図」は、シミュレーションで安価に作れる。} 方向ごとに異なる限界（前後と左右）を学習方策の制約として与えれば、実機で転倒させずに学習を進められる。
\end{enumerate}

\section{限界と次の課題}

\begin{itemize}
\tightlist
\item
  腕は 1 リンク（肘なし）。肘を加えると手先位置と腕の向きを独立に選べ、冗長性を姿勢制御に使える。
\item
  地上モデルは台車が静止している前提（転倒・滑りの開始条件のみ）。
\item
  宇宙での 3 次元の最適ループ（与えた目標姿勢に最短で到達する腕の動作系列）は、まだ求めていない。 非可換性があるため、これは学習・計画の良い課題になる。
\end{itemize}

\Detail{ch07}
\chapter{肩の 3 次元運動：AMBAC から人体工学へ（Sim08）}\label{ch:shoulder}

\section{なぜ肩なのか}

Sim07 で、球面座標の肩（2 軸ジンバル）を使って 3 次元の AMBAC を扱った。 ところが、人の肩は 2 軸ではない。少なくとも次の 3 つを持つ。

\begin{itemize}
\tightlist
\item
  腕の向き（2 自由度）
\item
  腕の長軸まわりのひねり（内旋・外旋）
\item
  肩甲骨・鎖骨（上肢帯）の動きによる、肩関節の中心そのものの移動
\end{itemize}

本章では、AMBAC で使った幾何学的位相の考え方を人の肩に当てはめる。 そのうえで、\textbf{人間らしい動きを実演するために機械に何が必要か} を、腕組みを例に定量化する。

\section{Part A: コッドマンのパラドックス ＝ 肩の幾何学的位相}

ガウス・ボンネの定理による導出は付録~\ref{app:codman} を参照。

\begin{figure}[htbp]\centering
\resizebox{0.7\linewidth}{!}{% コッドマンのパラドックス：球面三角形
\tdplotsetmaincoords{70}{125}
\begin{tikzpicture}[>=Stealth,font=\small,tdplot_main_coords,scale=1.9]
\shade[ball color=gray!10,opacity=0.4,tdplot_screen_coords] (0,0) circle (1);
\draw[->,gray] (0,0,0)--(1.3,0,0) node[anchor=north east]{前};
\draw[->,gray] (0,0,0)--(0,-1.3,0) node[anchor=north]{横（右）};
\draw[->,gray] (0,0,0)--(0,0,-1.3) node[anchor=north]{下};
\draw[very thick,red,->] plot[domain=0:90,samples=30,variable=\t] ({sin(\t)},0,{-cos(\t)});
\draw[very thick,red,->] plot[domain=0:90,samples=30,variable=\t] ({cos(\t)},{-sin(\t)},0);
\draw[very thick,red,->] plot[domain=0:90,samples=30,variable=\t] (0,{-cos(\t)},{-sin(\t)});
\node[red,font=\footnotesize,align=center] at (0.9,-0.9,-0.8) {囲んだ立体角\\$\Omega=\pi/2$};
\node[font=\footnotesize,align=left,tdplot_screen_coords] at (2.3,0.2) {下 → 前 → 横 → 下\\ひねらずに動かすと\\上腕が 90° ひねられる};
\end{tikzpicture}
}
\caption{コッドマンのパラドックス：腕の向きが単位球面上で「下 → 前 → 横 → 下」の球面三角形（立体角 $\pi/2$）を描くと、上腕は 90° ひねられる。}\label{fig:tikz-codman}
\end{figure}


\textbf{コッドマンのパラドックス} \cite{R34,R35}: 腕を「下 → 前 → 横 → 下」と動かして元の向きに戻すと、 一度もひねっていないのに、上腕が長軸まわりに 90° ひねられている。

「ひねらずに動かす」とは、各瞬間の回転軸が常に腕に垂直であること（ひねり方向の角速度がゼロ）を指す。 このとき、ループ 1 周後のひねりは次のようになる（ガウス・ボンネの定理 \cite{R11}）。

\[
\text{ひねり角} = \text{腕の向きが単位球面上で囲んだ立体角}
\]

\begin{longtable}[]{@{}
  >{\raggedright\arraybackslash}p{(\columnwidth - 4\tabcolsep) * \real{0.3333}}
  >{\raggedright\arraybackslash}p{(\columnwidth - 4\tabcolsep) * \real{0.3333}}
  >{\raggedright\arraybackslash}p{(\columnwidth - 4\tabcolsep) * \real{0.3333}}@{}}
\toprule\noalign{}
\begin{minipage}[b]{\linewidth}\raggedright
腕の動き
\end{minipage} & \begin{minipage}[b]{\linewidth}\raggedright
上腕のひねり
\end{minipage} & \begin{minipage}[b]{\linewidth}\raggedright
囲んだ立体角
\end{minipage} \\
\midrule\noalign{}
\endhead
\bottomrule\noalign{}
\endlastfoot
下 → 前 → 横 → 下（古典的なコッドマンの動作） & +90.00° & +90.00° \\
半頂角 20° の円錐を描いて下ろす & +21.71° & +21.71° \\
半頂角 45° の円錐を描いて下ろす & +105.44° & +105.44° \\
半頂角 90° の円錐を描いて下ろす & +360.00° & +360.00° \\
\end{longtable}

\FIG{sim08_shoulder.png}

\textbf{AMBAC との対応}

\begin{longtable}[]{@{}
  >{\raggedright\arraybackslash}p{(\columnwidth - 4\tabcolsep) * \real{0.3333}}
  >{\raggedright\arraybackslash}p{(\columnwidth - 4\tabcolsep) * \real{0.3333}}
  >{\raggedright\arraybackslash}p{(\columnwidth - 4\tabcolsep) * \real{0.3333}}@{}}
\toprule\noalign{}
\begin{minipage}[b]{\linewidth}\raggedright
\end{minipage} & \begin{minipage}[b]{\linewidth}\raggedright
AMBAC（第~\ref{ch:theory}章）
\end{minipage} & \begin{minipage}[b]{\linewidth}\raggedright
肩（本章）
\end{minipage} \\
\midrule\noalign{}
\endhead
\bottomrule\noalign{}
\endlastfoot
形状空間 & 関節角 (q\textsubscript{1}, q\textsubscript{2}) & 腕の向き（単位球面） \\
ループ 1 周の正味の変化 & 胴体の回転 Δθ & 上腕のひねり \\
決まり方 & 曲率 F の面積分 & 球面の曲率（=1）の面積分 ＝ 立体角 \\
\end{longtable}

\textbf{どちらも「形を元に戻しても、向きだけが進む」幾何学的位相} である。 AMBAC の数学は、宇宙の巨大ロボットと人の肩に共通している。

\textbf{2 軸ジンバル（Sim07 の肩）の場合}: 上腕の姿勢が関節角 (ψ, α) だけで決まるので、関節角が戻れば姿勢も戻り、 パラドックスは起きない。代わりに、ひねりを独立に選ぶことができない。

\textbf{示唆} - 人の肩の「ひねり」は、どの順序・どの基準で測るかによって値が変わる。 バイオメカニクスで肩の角度の測り方の標準（回転の順序の取り決め）が必要とされるのは、このためである。 - 人の動きをモーションキャプチャーでロボットに移すとき、ひねりの扱いを誤ると、同じ腕の向きでも手のひらの向きが変わる。 \textbf{3 自由度目（ひねり）を持ち、それを明示的に計画すること} が、人間らしい動きを再現する機械の要件になる。

\section{Part B: 腕組みは可能か}

\begin{figure}[htbp]\centering
\resizebox{0.6\linewidth}{!}{% 上から見た肩：胸郭、鎖骨、前方突出、腕組み
\begin{tikzpicture}[>=Stealth,font=\small,scale=11]
\draw[thick,fill=gray!12] (0,0) ellipse (0.11 and 0.14);
\node at (0,0) {胸郭};
\coordinate (SC) at (0.06,-0.02);
\coordinate (GH0) at (0.0,-0.19);
\coordinate (GH1) at ($(SC)!1!20:(GH0)$);
\draw[thick] (SC)--(GH0); \draw[thick,blue!70] (SC)--(GH1);
\fill (SC) circle (0.004) node[right=3pt,font=\footnotesize,fill=white,inner sep=1pt]{胸鎖関節};
\fill (GH0) circle (0.006) node[left,font=\footnotesize]{肩関節（基準）};
\fill[blue!70] (GH1) circle (0.006);
\draw[->,blue!70] ($(SC)!0.8!(GH0)$) to[bend right=30] ($(SC)!0.8!(GH1)$);
\node[blue!70!black,font=\footnotesize,align=left] at (0.17,-0.21) {前方突出\\（肩甲骨・鎖骨）};
\draw[line width=4pt,orange!70,opacity=0.8] (GH1)--(0.15,-0.10)--(0.20,0.15);
\fill[orange!90!black] (0.20,0.15) circle (0.007) node[right,font=\footnotesize]{手（左上腕へ）};
\draw[->] (-0.2,-0.33)--(-0.12,-0.33) node[right,font=\footnotesize]{$x$（前）};
\draw[->] (-0.2,-0.33)--(-0.2,-0.25) node[above,font=\footnotesize]{$y$（左）};
\end{tikzpicture}
}
\caption{上から見た肩のモデル。鎖骨が胸鎖関節まわりに回ると（前方突出）、肩関節の中心が前方・内側へ動き、腕を胸の前で交差させやすくなる。}\label{fig:tikz-girdle}
\end{figure}

\textbf{姿勢}: 右手を左上腕（肘の少し上）に置き、前腕を胸の前で交差させる（目標点 (0.20, 0.15, 0.00) m、許容誤差 1.5 cm）。 \textbf{制約}: 腕（上腕・前腕を太さのある円柱とする）が胸郭（楕円柱）に食い込まないこと、関節の可動域。

\begin{longtable}[]{@{}
  >{\raggedright\arraybackslash}p{(\columnwidth - 2\tabcolsep) * \real{0.5000}}
  >{\raggedright\arraybackslash}p{(\columnwidth - 2\tabcolsep) * \real{0.5000}}@{}}
\toprule\noalign{}
\begin{minipage}[b]{\linewidth}\raggedright
肩の機構
\end{minipage} & \begin{minipage}[b]{\linewidth}\raggedright
中身
\end{minipage} \\
\midrule\noalign{}
\endhead
\bottomrule\noalign{}
\endlastfoot
2 軸ジンバル & Sim07 の肩。ひねりを独立に選べない \\
球関節（人の可動域） & 3 自由度。水平内転 ≤ 45°、内旋 ≤ 70°、外旋 ≤ 90°、肘 ≤ 145° \\
球関節＋上肢帯 & 上記に加えて、鎖骨の回転による肩の前方突出（≤ 25°） \\
球関節（可動域制限なし） & 幾何学的に可能かどうかだけを見る \\
\end{longtable}

\subsection{B1: 合否（標準的な胸板、奥行 0.22 m）}

\begin{longtable}[]{@{}
  >{\raggedright\arraybackslash}p{(\columnwidth - 4\tabcolsep) * \real{0.3333}}
  >{\raggedright\arraybackslash}p{(\columnwidth - 4\tabcolsep) * \real{0.3333}}
  >{\raggedright\arraybackslash}p{(\columnwidth - 4\tabcolsep) * \real{0.3333}}@{}}
\toprule\noalign{}
\begin{minipage}[b]{\linewidth}\raggedright
肩の機構
\end{minipage} & \begin{minipage}[b]{\linewidth}\raggedright
結果
\end{minipage} & \begin{minipage}[b]{\linewidth}\raggedright
手の誤差
\end{minipage} \\
\midrule\noalign{}
\endhead
\bottomrule\noalign{}
\endlastfoot
2 軸ジンバル & \textbf{不可} & 25 cm \\
球関節（人の可動域） & 不可 & 6.6 cm \\
球関節＋上肢帯 & 不可（惜しい） & 2.6 cm \\
可動域制限なし & 可能 & 0 cm（ただしひねりが不自然） \\
\end{longtable}

\begin{itemize}
\tightlist
\item
  \textbf{2 軸ジンバルでは、胸板の厚さによらず届かない}（誤差 24〜26 cm）。 腕を前に出すと、前腕の曲がる面（肘の屈曲面）が上腕と同じ鉛直面に固定され、前腕を胸の前で横に渡せない。
\item
  人の可動域を与えた球関節では、内旋 70° の上限に張り付いて届かなかった。 ただし、内旋 70° は腕を横に 90° 上げた姿勢での標準値で、腕を前に下ろした姿勢でどこまでひねれるかは別である。 さらに Part A のとおり、ひねりの数値は測り方で変わる。そこで B2 では、合否ではなく「最低限必要な内旋量」を求めた。
\end{itemize}

\subsection{B2: 腕組みに最低限必要な内旋量}

水平内転 ≤ 45°、前方突出 ≤ 25° の下で、手を目標に届かせるのに必要なひねりの最小値を求めた。

\begin{longtable}[]{@{}
  >{\raggedright\arraybackslash}p{(\columnwidth - 4\tabcolsep) * \real{0.3333}}
  >{\raggedright\arraybackslash}p{(\columnwidth - 4\tabcolsep) * \real{0.3333}}
  >{\raggedright\arraybackslash}p{(\columnwidth - 4\tabcolsep) * \real{0.3333}}@{}}
\toprule\noalign{}
\begin{minipage}[b]{\linewidth}\raggedright
胸板の奥行
\end{minipage} & \begin{minipage}[b]{\linewidth}\raggedright
球関節のみ
\end{minipage} & \begin{minipage}[b]{\linewidth}\raggedright
球関節＋上肢帯（必要な前方突出）
\end{minipage} \\
\midrule\noalign{}
\endhead
\bottomrule\noalign{}
\endlastfoot
0.18 m（細身） & 86° & 71°（12°） \\
0.22 m（標準） & 93° & 73°（16°） \\
0.26 m（厚い胸板・装甲） & \textbf{110°（人の範囲を超える）} & 75°（20°） \\
\end{longtable}

参考: 人の内旋の可動域はおよそ 70〜90°（姿勢と測り方に依存）。

\textbf{考察} - \textbf{肩関節だけでは、腕組みに人の限界ぎりぎりか、それ以上の内旋が必要になる。} 上肢帯（肩甲骨・鎖骨）が肩関節を前に出すことで、必要な内旋が 15〜35° 減り、人の範囲に収まる。 \textbf{人が腕を組めるのは、肩甲骨が動くから} と言える。 - \textbf{胸板が厚いほど難しい。} 厚い胸板・装甲では、肩関節だけでは人の範囲を超える内旋が必要になり、上肢帯の動き（20°）が必須になる。 アクションフィギュアで、肩の付け根を前に動かせる関節を胸に設けると腕組みができるようになるのと同じ理屈である。 装甲で胸の厚いガンダム型の機体が腕を組むには、肩の付け根ごと前に出す機構が要る。

\section{人間らしい動きのために機械が必要とするもの（ここまでの要件）}

\begin{longtable}[]{@{}
  >{\raggedright\arraybackslash}p{(\columnwidth - 2\tabcolsep) * \real{0.5000}}
  >{\raggedright\arraybackslash}p{(\columnwidth - 2\tabcolsep) * \real{0.5000}}@{}}
\toprule\noalign{}
\begin{minipage}[b]{\linewidth}\raggedright
要件
\end{minipage} & \begin{minipage}[b]{\linewidth}\raggedright
根拠
\end{minipage} \\
\midrule\noalign{}
\endhead
\bottomrule\noalign{}
\endlastfoot
肩は 3 自由度（ひねりを含む）にする & 2 軸ジンバルでは腕組みが不可能（B1）。ひねりを独立に選べない（A） \\
ひねりを明示的に計画・計測する & コッドマンのパラドックス：腕の向きの経路によってひねりが勝手に変わる（A） \\
上肢帯（肩の付け根の移動）を持たせる & 腕組みに必要な内旋が人の範囲に収まる。胸板が厚い機体では必須（B2） \\
座標の特異点を作業範囲の外に置く & 球面座標の極で関節速度が発散する（Sim07 A3） \\
胸・装甲の厚さと腕の可動域を同時に設計する & 胸が厚いほど必要な可動域が増える（B2） \\
\end{longtable}

\section{限界と次の課題}

\begin{itemize}
\tightlist
\item
  胸郭は楕円柱、腕は円柱で近似した。筋肉の盛り上がり・衣服・装甲の形状は扱っていない。
\item
  可動域は「水平内転」「スイング・ツイスト分解によるひねり」で単純化した。実際の可動域は姿勢に依存して連続的に変わる。
\item
  左右の前腕の重なり（どちらが上か）と、左右の腕どうしの干渉は扱っていない。
\item
  股関節（同じく球関節だが、骨盤のくぼみが深く安定な代わりに可動域が狭い）への展開は今後の課題。
\end{itemize}

\Detail{ch08}

\part{結論}
\chapter{結論と第 2 段階の構想}\label{ch:conclusion}

\section{問いへの答え}

「二足歩行は本当にバランスが良いのか？　なぜ人は二本脚なのか？」という問いに対し、第 1 段階の結果は次のように答える。

\begin{enumerate}
\item \textbf{静的には悪い。} 二本脚は支持範囲が狭く重心が高いため、荷物を持つだけで転倒しうる（第~\ref{ch:legs}章）。
\item \textbf{能動的には良い。} 動かせる重い上体（腰の関節）が、腕の反作用を打ち消し、段差を上る力を分担する。
  これは手足の代わりに上体を使う地上版の AMBAC である（第~\ref{ch:legs}章、第~\ref{ch:steps}章）。
\item \textbf{人間の大きさで最も活きる。} 段差では人間大の脚が戦車型の 2 倍以上を越えるが、大型化すると逆転する（第~\ref{ch:steps}章）。
\item \textbf{宇宙では、脚は姿勢制御の道具になる。} 脚は AMBAC で最も効く重い手足で、外すと姿勢変更の能力がほぼ半減する（第~\ref{ch:legs}章）。
  ただし推進剤を惜しまなければ、機動力はスラスタが圧倒的に上である（第~\ref{ch:dyn}章、第~\ref{ch:ai}章）。
\end{enumerate}

\section{第 1 段階で分かったこと}

\subsection{1. AMBAC の物理}

\begin{longtable}{p{0.27\linewidth}p{0.56\linewidth}p{0.08\linewidth}}
\toprule
問い & 答え & 章 \\ \midrule \endhead
AMBAC は物理的に成立するか & 成立する。角運動量保存の下で、手足のループ運動（幾何学的位相）により任意角の旋回が可能 & \ref{ch:theory}, \ref{ch:planar} \\
人工衛星の技術との関係は & リアクションホイールは「曲率ゼロ・連続回転できる関節」で、AMBAC の特殊な場合。運動量飽和も同じ現象 & \ref{ch:rw} \\
巨大ロボで使えるか & 回転角はスケール不変だが、速さは $1/s$ で落ちる。18 m 級で機敏に動くには、相似則の約 170 倍のトルクが必要 & \ref{ch:dyn} \\
学習で獲得できるか & 「ループを描く振動子」という構造を与えれば可能で、専門家制御を上回った。構造なしでは失敗する & \ref{ch:ai} \\
最適制御と学習の分担は & 計画は最適制御、実行時の誤差・外乱・揺れへの対処は学習が得意 & \ref{ch:ai} \\
\bottomrule
\end{longtable}

\subsection{2. 下半身の形状}

\begin{longtable}{p{0.27\linewidth}p{0.56\linewidth}p{0.08\linewidth}}
\toprule
問い & 答え & 章 \\ \midrule \endhead
腕を速く動かすと何が起きるか & 腕の反作用で台車が転倒・滑る。安全な最短時間は $\sqrt{\text{大きさ}/\text{重力}}$ に比例（フルード則）。荷物を体に引き寄せる軌道で 36\% 短縮 & \ref{ch:tank} \\
足は飾りか & 宇宙では飾りではない。推進剤を惜しまなければ飾りになる。地上の平地作業ではほぼ飾り（戦車型＋腰の方が速く安定） & \ref{ch:legs} \\
二本脚はバランスが良いか & 静的には悪い。腰の戦略（上体の前後傾）で補って初めて荷物を運べる & \ref{ch:legs} \\
段差ではどうか & 人間大では脚が戦車型の 2 倍以上。大型化すると逆転し、18 m 級の脚は 5 cm しか越えられない & \ref{ch:steps} \\
\bottomrule
\end{longtable}

\subsection{3. 上半身（肩と肩甲骨）}

\begin{longtable}{p{0.27\linewidth}p{0.56\linewidth}p{0.08\linewidth}}
\toprule
問い & 答え & 章 \\ \midrule \endhead
3 次元の肩で何ができるか & 2 自由度の肩 1 つの円錐運動で任意の姿勢へ回れる。ただし回転の順序で結果が変わる（非可換） & \ref{ch:3d} \\
肩の座標の注意点は & 球面座標（2 軸の肩）は極の近くで関節速度が発散する（ジンバルロック） & \ref{ch:3d} \\
地上で腕を旋回させると & 接地面の狭い横方向が最も弱い。腕を下げて旋回半径と高さを減らすと 52\% 短縮 & \ref{ch:3d} \\
人の肩との関係は & コッドマンのパラドックスは肩の幾何学的位相（ひねり＝腕の向きが囲んだ立体角） & \ref{ch:shoulder} \\
腕組みには何が必要か & ひねりを含む 3 自由度の肩と、肩甲骨による肩の付け根の移動。胸が厚いほど肩甲骨の動きが必須 & \ref{ch:shoulder} \\
\bottomrule
\end{longtable}

\section{フィジカル AI に必要なスペックとトレードオフ}

\subsection{必要なスペック}

\begin{longtable}{p{0.40\linewidth}p{0.53\linewidth}}
\toprule
スペック & 根拠 \\ \midrule \endhead
肩は 3 自由度（ひねりを含む） & 2 軸では腕組みができず、ひねりを独立に選べない（第~\ref{ch:shoulder}章） \\
肩甲骨（上肢帯）に当たる、肩の付け根を動かす機構 & 腕組みに必要な内旋が人の範囲に収まる。胸の厚い機体では必須（第~\ref{ch:shoulder}章） \\
動かせる重い上体（腰の関節） & 二本脚の静的な不安定さを補う。段差の上りにも効く（第~\ref{ch:legs}章、第~\ref{ch:steps}章） \\
重い下半身（宇宙で使うなら） & 脚は AMBAC で最も効く重い手足（第~\ref{ch:legs}章） \\
座標の特異点を作業範囲の外に置く設計 & ジンバルロックで関節速度が発散する（第~\ref{ch:3d}章） \\
大きさに見合ったトルク密度 & 必要トルクは大きさの 4〜5 乗、相似則で使えるトルクは 3 乗（第~\ref{ch:dyn}章、第~\ref{ch:steps}章） \\
\bottomrule
\end{longtable}

\subsection{トレードオフしなければならない事象}

\begin{longtable}{p{0.30\linewidth}p{0.63\linewidth}}
\toprule
トレードオフ & 内容 \\ \midrule \endhead
支持範囲 ↔ 段差への強さ & 戦車型は平地で安定だが段差に弱い。二本脚はその逆（第~\ref{ch:legs}章、第~\ref{ch:steps}章） \\
機体の大きさ ↔ 機敏さ・段差 & 大きいほど遅く、脚は段差を越えられなくなる（第~\ref{ch:dyn}章、第~\ref{ch:tank}章、第~\ref{ch:steps}章） \\
胸・装甲の厚さ ↔ 腕の可動域 & 胸が厚いほど、腕組みに必要なひねり・肩甲骨の動きが増える（第~\ref{ch:shoulder}章） \\
肩の自由度 ↔ 制御の複雑さ & 3 自由度にすると、経路によってひねりが勝手に変わる（コッドマンのパラドックス）ので、ひねりの計画・計測が必要（第~\ref{ch:shoulder}章） \\
推進剤 ↔ 時間 & スラスタは速いが推進剤を使う。AMBAC は推進剤ゼロだが遅く、周回内で大きく揺れる（第~\ref{ch:dyn}章、第~\ref{ch:ai}章） \\
重心を前に出す ↔ 後ろへの転倒 & 戦車型の段差越えで、重心の位置は必要な摩擦と転倒のしやすさを入れ替える（第~\ref{ch:steps}章） \\
\bottomrule
\end{longtable}

\section{用途に応じた形状の選択}

\begin{longtable}{p{0.40\linewidth}p{0.25\linewidth}p{0.25\linewidth}}
\toprule
用途 & 有利な下半身 & 根拠 \\ \midrule \endhead
宇宙・長期運用・推進剤に制約あり & 重い脚（AMBAC 用） & 第~\ref{ch:rw}〜\ref{ch:ai}章、第~\ref{ch:legs}章 \\
宇宙・短時間の高機動 & 推進剤（脚は飾り） & 第~\ref{ch:dyn}章、第~\ref{ch:ai}章 \\
地上・平地での作業 & 戦車型＋腰 & 第~\ref{ch:tank}章、第~\ref{ch:legs}章 \\
地上・平地で旋回が多い & 車輪（全方向移動）など & 第~\ref{ch:legs}章 \\
地上・段差（人間大） & 脚 & 第~\ref{ch:steps}章 \\
地上・段差（巨大機） & 戦車型 & 第~\ref{ch:steps}章 \\
\bottomrule
\end{longtable}

\section{人間らしい動きの「基準値」をどこに置くか}

人間らしい動きの実現には、モーションキャプチャー（人体運動の計測データ）を手本にする方法と、制御理論で動作を設計する方法がある。
第 1 段階の結果からは、次のように整理できる。

\begin{itemize}
\item \textbf{人体運動は「形（スタイル）」の手本として強い。} 自然さ・予測しやすさ・人間との協調に向く。
\item \textbf{ただし、物理的に正しい手本とは限らない。} モーションキャプチャーが測るのは動きの形で、力は測らない。
  AMBAC の効き方は手足の質量比で決まるため（第~\ref{ch:planar}章、第~\ref{ch:rw}章）、質量配分の違うロボットに人の動きをそのまま移すと、同じ動きでも胴体の回り方が変わる。
  肩のひねりは測り方で値が変わる（第~\ref{ch:shoulder}章）。
\item \textbf{性能の上限ではない。} 学習方策は、人手で設計したループを上回った（第~\ref{ch:ai}章）。
\item \textbf{一方で、人の動きの合理性は物理から説明できる場合がある。}
  荷物を体に引き寄せる（第~\ref{ch:tank}章）、腕を下げて旋回する（第~\ref{ch:3d}章）、重い物を持つときに体を反らせる（第~\ref{ch:legs}章）、
  段差で前かがみになる（第~\ref{ch:steps}章）は、いずれも最適化が再発見した人間らしい動きである。
\end{itemize}

したがって有力なのは組み合わせである。人体運動を「目指す形」として与え、物理モデル・制御理論で「実現可能性（力・運動量・制約）」を保証し、
学習でその間を埋める（モーションキャプチャーを手本に物理シミュレーション上で強化学習する手法 \cite{R24} がこれに近い）。

\section{第 2 段階：校正値を使った実現性の検証}

第 2 段階では、組み合わせた形状（アセンブル）について最適な動作の校正値を定義し、要求した動きや性能を実現できるかを検証する。

\begin{figure}[htbp]\centering
\begin{tikzpicture}[box/.style={draw,rounded corners,align=left,font=\footnotesize,text width=3.2cm,minimum height=3.6cm,fill=gray!6},
  head/.style={font=\small\bfseries},>=Stealth]
\node[box] (a) at (0,0) {下半身：脚／戦車／車輪\\肩の自由度・肩甲骨\\腰・上体の可動域\\質量配分\\アクチュエータ\\大きさ・重力環境};
\node[box] (b) at (3.9,0) {姿勢変更能力（第~\ref{ch:rw}, \ref{ch:ai}, \ref{ch:3d}章）\\転倒・滑り（第~\ref{ch:tank}, \ref{ch:legs}章）\\段差（第~\ref{ch:steps}章）\\定格荷重（第~\ref{ch:3d}章）\\人らしい姿勢（第~\ref{ch:shoulder}章）};
\node[box] (c) at (7.8,0) {慣性パラメータ（AMBAC ループで自己同定）\\摩擦・接触\\可動域・アクチュエータ特性\\人の動きのデータ};
\node[box] (d) at (11.7,0) {要求した動き・性能を満たすか\\制御ゲイン\\軌道最適化\\学習方策の制約\\形状の選び直し};
\node[head,above=2pt of a] {形状の組み合わせ};
\node[head,above=2pt of b] {シミュレーション評価};
\node[head,above=2pt of c] {校正値};
\node[head,above=2pt of d] {実現性の検証};
\draw[->,thick] (a)--(b); \draw[->,thick] (b)--(c); \draw[->,thick] (c)--(d);
\draw[->,thick,dashed] (d.south) -- ++(0,-0.5) -| node[pos=0.25,below,font=\footnotesize]{満たさなければ形状・制御を見直す} (a.south);
\end{tikzpicture}
\caption{第 2 段階の枠組み：アセンブル型の設計・評価・校正・検証。}\label{fig:phase2}
\end{figure}

\begin{enumerate}
\item \textbf{形状の組み合わせ}：下半身の形式、肩の自由度、肩甲骨、腰、質量配分、アクチュエータ、大きさ、重力環境を部品として組み合わせる。
\item \textbf{シミュレーションによる評価}：第 1 段階で作った評価を、組み合わせた機体ごとにまとめて実行し、性能表を作る。
\item \textbf{校正値の定義}：シミュレーションと実機の差を埋める値を定義する。慣性パラメータ（AMBAC ループによる自己同定）、
  摩擦・接触、可動域、アクチュエータ特性（油圧・電動）、人の動きのデータなど。
\item \textbf{実現性の検証とチューニング}：校正値を使って、要求した動き・性能（例：腕組み、段差、旋回時間、荷物運び）を実現できるかを検証する。
  満たさなければ、制御・軌道・学習方策の制約を調整し、必要なら形状を選び直す。
\end{enumerate}

\section{残された課題}

\begin{itemize}
\item アクチュエータの動特性（油圧・電動）：油の圧縮性、バルブの非線形性、応答遅れ・帯域制限。油圧はトルク密度の係数を上げるが、相似拡大の指数（$\propto s^3$）は変えない。
\item 動的な動作：踏み出し（ステップ戦略）、走り込み・蹴り出しによる段差の乗り越え、不整地。
\item 周回内の揺れを抑える AMBAC ループの最適化、弾性要素によるエネルギー回生。
\item 最適制御の計画と学習した残差方策のハイブリッド制御、3 次元での目標姿勢への最短到達。
\item 股関節を含む全身モデル、柔軟体・流体などの非剛体効果。
\end{itemize}


\appendix
\chapter{式の導出}\label{app:deriv}

本文で結果だけを示した式を、ここで導出する。記号は本文と同じものを使う。
平面の外積は $\bm u\times\bm v = u_xv_y-u_yv_x$（スカラー）とし、$J=\begin{psmallmatrix}0&-1\\1&0\end{psmallmatrix}$ を 90° 回転、$R(\theta)$ を角度 $\theta$ の回転行列とする。
平面ベクトルについて $\bm u\times(J\bm v)=\bm u\cdot\bm v$ が成り立つことを繰り返し使う。

\section{胴体＋手足 1 本の角運動量（閉じた形）}\label{app:twobody}

胴体（質量 $m_b$、重心まわりの慣性 $I_b$）に、回転関節で手足（質量 $m_l$、重心まわりの慣性 $I_l$）が付いた平面の系を考える。
胴体座標で、関節の位置を $\bm a$（胴体重心から）、関節から手足の重心までの距離を $\ell_c$ とする。
胴体の姿勢角を $\theta$、関節角を $q$ とすると、胴体重心から手足重心へのベクトルは
\begin{equation}
\bm r = R(\theta)\,\bm r_b,\qquad \bm r_b(q) = \bm a + \ell_c\,\bm e(q),\qquad \bm e(q)=(\cos q,\ \sin q)^\top .
\end{equation}

\paragraph{2 体問題への帰着}
外力がなければ系重心は等速運動し、重心系で見れば静止とみなせる。2 体の重心まわりの角運動量は、各剛体の自転の角運動量と、
相対運動の角運動量（換算質量 $\mu = m_bm_l/(m_b+m_l)$ を使う）の和になる：
\begin{equation}
L = I_b\dot\theta + I_l(\dot\theta+\dot q) + \mu\,\bm r\times\dot{\bm r}.
\end{equation}

\paragraph{相対運動の項}
$\dot{\bm r} = R(\theta)\left(\dot\theta J\bm r_b + \dot{\bm r}_b\right)$、$\dot{\bm r}_b = \dot q\,\ell_c J\bm e(q)$ である。
外積は回転で不変なので、
\begin{align}
\bm r\times\dot{\bm r} &= \bm r_b\times\left(\dot\theta J\bm r_b + \dot q\,\ell_c J\bm e\right)
 = \dot\theta\,|\bm r_b|^2 + \dot q\,\ell_c\,\bm r_b\cdot\bm e \\
&= \dot\theta\left(|\bm a|^2+\ell_c^2+2\ell_c\,\bm a\cdot\bm e\right) + \dot q\left(\ell_c^2+\ell_c\,\bm a\cdot\bm e\right).
\end{align}
$\bm a = |\bm a|(\cos\varphi_a,\sin\varphi_a)$ と書けば $\bm a\cdot\bm e = |\bm a|\cos(q-\varphi_a)$ である。

\paragraph{結果}
\begin{equation}\boxed{
\begin{aligned}
L &= A(q)\dot\theta + B(q)\dot q,\\
A(q) &= I_b + I_l + \mu\left(|\bm a|^2+\ell_c^2+2\ell_c|\bm a|\cos(q-\varphi_a)\right),\\
B(q) &= I_l + \mu\left(\ell_c^2+\ell_c|\bm a|\cos(q-\varphi_a)\right).
\end{aligned}}\label{eq:AB-closed}
\end{equation}
$L$ は $\theta$ に依存しない（回転対称性）。$A$ は「形状 $q$ で固まった系全体の慣性モーメント」、$B$ は関節運動と姿勢の結合である。

\paragraph{数値の確認}
Sim01 の人間大モデル（$m_b=60$ kg、$I_b=60(0.35^2+1.0^2)/12=5.613$ kg\,m$^2$、腕 $m_l=4$ kg、長さ 0.7 m、$\bm a=(0,0.35)$ m、$\ell_c=0.35$ m）で、
腕を前に水平に向けた $q=0$ では $\cos(q-\varphi_a)=\cos(-90^\circ)=0$ なので、
$\mu=3.75$ kg、$A = 5.613+0.163+3.75\times0.245 = 6.695$、$B = 0.163+3.75\times0.1225=0.623$、
$\Gamma=B/A=0.0930$ となる。これは数値計算（平面モデル・3 次元モデル）の値 0.09302 と一致する。

\section{AMBAC の基本式と幾何学的位相}\label{app:connection}

\paragraph{基本式}
手足が $n$ 本のときも、同じ議論で $L = A(q)\dot\theta + B(q)\cdot\dot q$（$B$ は $n$ 成分）となる。
静止から始めれば $L\equiv 0$ なので
\begin{equation}
\dot\theta = -\Gamma(q)\cdot\dot q,\qquad \Gamma(q) = B(q)/A(q). \label{eq:basic}
\end{equation}
$\Gamma$ は幾何力学で機械的接続（mechanical connection）と呼ばれる \cite{R5,R4}。
式~\eqref{eq:basic} には時間が陽に現れないので、$\theta$ の変化は関節空間の経路だけで決まり、動かす速さに依存しない。

\paragraph{手足 1 本の場合}
$n=1$ なら $\theta(q) = \theta_0 - \int_{q_0}^{q}\Gamma(q')\,dq'$ は $q$ だけの関数である。関節を元に戻せば姿勢も必ず元に戻る（ホロノミック）。

\paragraph{手足 2 本の場合：ストークスの定理}
関節空間の閉曲線 $C$ を 1 周すると、
\begin{equation}
\Delta\theta = -\oint_C \left(\Gamma_1\,dq_1+\Gamma_2\,dq_2\right)
 = -\iint_S\left(\frac{\partial\Gamma_2}{\partial q_1}-\frac{\partial\Gamma_1}{\partial q_2}\right)dq_1\,dq_2
 = -\iint_S F\,dq_1\,dq_2. \label{eq:stokes}
\end{equation}
ここで $S$ は $C$ が反時計回りに囲む領域、$F$ は接続の曲率である。$F\neq 0$ なら、形状が戻っても姿勢は $\Delta\theta$ だけ進む（幾何学的位相 \cite{R6}）。

\paragraph{小さなループは 2 次の効果}
中心 $\bar q$、振幅 $\varepsilon$ の小さなループでは、$\Delta\theta \approx -F(\bar q)\times(\text{面積}) = O(\varepsilon^2)$ である。
1 次の効果（往復で打ち消す可逆な首振り）は $O(\varepsilon)$ なので、静止に近い方策から学習を始めると、ループを大きくする方向の勾配は小さい（第~\ref{ch:ai}章）。
制御理論の言葉では、正味の回転は入力ベクトル場のリー括弧から生じる \cite{R10}。

\paragraph{1 周あたりの回転の上界}
関節の可動域 $Q$ の中で、$\Delta\theta$ を最大にする閉曲線は、$F$ が同じ符号の領域をちょうど囲むものである。したがって
\begin{equation}
|\Delta\theta| \le \max\left(\iint_{Q\cap\{F<0\}}|F|\,dq,\ \iint_{Q\cap\{F>0\}}F\,dq\right).
\end{equation}

\section{リアクションホイールは曲率ゼロの関節}\label{app:rw}

ホイールは重心が回転軸上にある手足とみなせる。式~\eqref{eq:AB-closed} で $\ell_c=0$ とおくと
\begin{equation}
A = I_b + I_w + \mu|\bm a|^2,\qquad B = I_w,\qquad \Gamma = \frac{I_w}{I_b+I_w+\mu|\bm a|^2} = \text{一定}.
\end{equation}
$\Gamma$ が一定なので $F=0$ で幾何学的位相は生じない。その代わり、ホイールは何回転でも回せるので、$\Delta\theta=-\Gamma\Delta\varphi_w$ は無制限に大きくできる。

\section{角運動量がある場合と「飽和」}\label{app:saturation}

$L\neq0$ のとき、$\dot\theta = L/A(q) - \Gamma(q)\cdot\dot q$ である。
周波数 $f$ で一定の速さでループを回すと、1 周（時間 $1/f$）あたりの変化は
\begin{equation}
\Delta\theta_{\text{周}} = \Delta\theta_c + L\int_0^{1/f}\frac{dt}{A(q(t))} = \Delta\theta_c + \frac{L}{f}\left\langle\frac1A\right\rangle,
\end{equation}
したがって周回平均の姿勢 $\bar\theta$ の変化率は
\begin{equation}
\dot{\bar\theta} \approx f\,\Delta\theta_c + L\left\langle\frac1A\right\rangle .
\end{equation}
外乱で $L$ が増え続けても姿勢を保つには $f = -L\langle 1/A\rangle/\Delta\theta_c$ が必要で、$|f|\le f_{\max}$ から
\begin{equation}
|L| \le h_{\mathrm{AMBAC}} = \frac{f_{\max}|\Delta\theta_c|}{\langle 1/A\rangle}.
\end{equation}
リアクションホイールでは $\dot\theta = (L-I_w\omega_w)/A$ なので、姿勢を保つには $\omega_w = L/I_w$、$|\omega_w|\le\omega_{\max}$ から $|L|\le h_{\mathrm{RW}} = I_w\omega_{\max}$ となる。
どちらも「角運動量を内部に溜める容量」であり、超えると飽和する。

\section{運動方程式と時間スケーリング}\label{app:eom}

一般化速度 $\bm u = (\dot\theta,\dot q)$ とすると、重心系の運動エネルギーは $T = \tfrac12\bm u^\top M(q)\bm u$ と書ける（$M$ は $\theta$ に依存しない）。
$M$ の第 1 行は $(A, B)$ そのものである。ラグランジュ方程式は
\begin{equation}
\frac{d}{dt}\left(M\bm u\right) - \frac12\begin{pmatrix}0\\ \bm u^\top\frac{\partial M}{\partial q_k}\bm u\end{pmatrix} = \begin{pmatrix}0\\ \bm\tau\end{pmatrix},
\end{equation}
すなわち $M\dot{\bm u} + \dot M\bm u - \tfrac12\bigl(0;\ \bm u^\top\partial_{q_k}M\,\bm u\bigr) = (0;\bm\tau)$ である。
第 1 行は $dL/dt=0$（$\theta$ が循環座標であることに対応する角運動量保存）そのものである。

\paragraph{時間スケーリング}
同じ経路を $f$ 倍速で回すと、速度は $f$ 倍、加速度は $f^2$ 倍になる。上式の左辺は加速度の 1 次と速度の 2 次の項からなるので、
\begin{equation}
\bm\tau \propto f^2,\qquad \text{パワー } \bm\tau\cdot\dot q \propto f^3 .
\end{equation}

\paragraph{長さのスケーリング}
長さを $s$ 倍、密度を同じに保つと、質量は $s^3$ 倍、慣性モーメントは $s^5$ 倍になる。したがって同じ周波数で回すのに必要なトルクは $s^5f^2$ に比例する。
アクチュエータを相似拡大すると、出せるトルクは（断面積 $s^2$）×（レバー $s$）$=s^3$ に比例する。両者を等しくおくと
\begin{equation}
f_{\max}\propto s^{-1},\qquad T_{90^\circ}\propto s,\qquad h_{\mathrm{AMBAC}} \propto f_{\max}\bar A \propto s^{-1}s^5 = s^4 .
\end{equation}
回転角 $\Delta\theta_c$ は $\Gamma$ が質量比・長さ比だけの関数なので、スケールに依存しない。

\section{ZMP と転倒・滑りの条件}\label{app:zmp}

地上で静止している台車に、点質量または剛体 $i$（質量 $m_i$、位置 $(x_i,z_i)$、加速度 $(\ddot x_i,\ddot z_i)$、慣性 $I_i$、角加速度 $\ddot\phi_i$）が載っているとする。
地面からの反力（垂直力の合力 $N$、水平力 $F_x$）は系全体のニュートン・オイラー方程式から決まる：
\begin{equation}
N = \sum_i m_i(\ddot z_i+g),\qquad F_x = \sum_i m_i\ddot x_i .
\end{equation}
地面上の点 $p$ まわりのモーメントのつり合いは
\begin{equation}
\sum_i\left[m_i\left((x_i-p)(\ddot z_i+g) - z_i\ddot x_i\right) + I_i\ddot\phi_i\right] = M_{\text{接触}}(p).
\end{equation}
接触力の合力が通る点（ZMP）では $M_{\text{接触}}=0$ なので
\begin{equation}
p = \frac{\sum_i m_i\left(x_i(\ddot z_i+g) - z_i\ddot x_i\right) + \sum_i I_i\ddot\phi_i}{\sum_i m_i(\ddot z_i+g)} . \label{eq:zmp}
\end{equation}
接地は押すことしかできないので、ZMP は接地範囲 $[-a,a]$ の中にしか来られない。外に出ようとすれば転倒が始まる \cite{R28}。
水平力は摩擦で出すので $|F_x|\le\mu N$ が必要で、超えると滑る。3 次元で点質量だけからなる場合は
\begin{equation}
p_x = \frac{\sum m_i\left(x_i(\ddot z_i+g)-z_i\ddot x_i\right)}{\sum m_i(\ddot z_i+g)},\qquad
p_y = \frac{\sum m_i\left(y_i(\ddot z_i+g)-z_i\ddot y_i\right)}{\sum m_i(\ddot z_i+g)} .
\end{equation}

\paragraph{フルード則}
長さを $s$ 倍、動作時間を $T$ とすると、位置は $s$ 倍、加速度は $s/T^2$ 倍になる。式~\eqref{eq:zmp} の分子・分母を $g$ で割ると、
$p/s$ は無次元量 $s/(gT^2)$ だけの関数になる。したがって転倒限界 $p/s=a/s$（一定）を与える時間は
\begin{equation}
T_{\min}\propto\sqrt{s/g}
\end{equation}
となる（フルード数 $v^2/(g\ell)$ を一定に保つ相似則 \cite{R29}）。Sim06 では比 $T_{\min}/\sqrt{s\,g_{\text{地球}}/g}$ が 12 通りの条件で 1.616 に一致した。

\section{履帯のヨーとスティックスリップ}\label{app:stickslip}

台車中心のターレットで腕を水平に振る。台車の並進は止まっているとし、台車中心まわりの角運動量を $L_o = A\dot\theta + B\dot q$（$A=I_b+I_a$、$B=I_a$ は一定）とする。
摩擦トルクの上限は、接地面（面積 $S$）に一様に圧力がかかるとして
\begin{equation}
\tau_f = \mu N\bar r,\qquad \bar r = \frac{1}{S}\iint_S|\bm r|\,dS .
\end{equation}
クーロン摩擦（静止・動摩擦）のモデル \cite{R30} では、
\begin{itemize}
\item 台車が止まっていて（$\dot\theta=0$）、必要な摩擦トルク $|B\ddot q|\le\tau_f$ なら止まったまま。
\item そうでなければ $A\ddot\theta = -B\ddot q - \tau_f\,\mathrm{sgn}(\dot\theta)$（止まっていた場合は滑り出す向きの符号）。
\end{itemize}
最小躍度軌道で角度 $\Delta q$ を時間 $T$ で振ると最大角加速度は $5.77\Delta q/T^2$ なので、滑り始める条件は $T<\sqrt{5.77B\Delta q/\tau_f}$ である。
「速く振ってゆっくり戻す」往復では、速い振りの間だけ滑るので、正味の回転が残る。

\section{3 次元の AMBAC}\label{app:3d}

胴体の姿勢を回転行列 $R$、胴体座標での角速度を $\bm\omega$ とする。重心系の角運動量（胴体座標）は
\begin{equation}
\bm L = A(q)\bm\omega + B(q)\dot q,\qquad
A = I_b + \sum_i m_i\left(|\bm r_i|^2E - \bm r_i\bm r_i^\top\right),\qquad
B = \sum_i m_i\,\bm r_i\times\frac{\partial(\bm r_i-\bm c)}{\partial q},
\end{equation}
（$\bm r_i$ は系重心から各点質量への位置、$E$ は単位行列）。$\bm L=0$ から
\begin{equation}
\bm\omega = -A^{-1}B\,\dot q,\qquad \dot R = R\,[\bm\omega]_\times .
\end{equation}
平面と違い、$\bm\omega$ は 3 成分のベクトルで、回転の合成は順序に依存する。小さなループ（振幅 $\beta$）での正味の回転ベクトルは、
接続の曲率（非可換なので $d\Gamma$ に加えて交換子の項を含む）とループの面積の積で与えられ、$O(\beta^2)$ である。Sim07 では指数 2.00 を得た。

\paragraph{球面座標の特異点}
肩の向きを $\bm u(\psi,\alpha) = (\cos\alpha\cos\psi,\ \cos\alpha\sin\psi,\ \sin\alpha)$ で表すと、
$|\partial\bm u/\partial\psi| = \cos\alpha$ なので、腕先の角速度 $\omega_{\text{tip}}$ を保つには $\dot\psi\approx\omega_{\text{tip}}/\cos\alpha$ が必要である。
極（$\alpha\to\pm90^\circ$）で、極からの角距離を $\delta$ とすると $\cos\alpha\approx\delta$ なので $\dot\psi\propto1/\delta$ で発散する（ジンバルロック \cite{R18}）。

\section{コッドマンのパラドックスとガウス・ボンネの定理}\label{app:codman}

上腕を「ひねらずに」動かすとは、上腕の角速度 $\bm\omega$ が常に腕の軸 $\bm u$ に垂直であること（$\bm\omega\cdot\bm u=0$）を言う。
このとき上腕に固定した、腕に垂直な方向ベクトルは、単位球面上で $\bm u$ が描く曲線に沿って平行移動される。
閉曲線 $C$ を 1 周したとき、平行移動されたベクトルの回転角は、$C$ が囲む領域 $S$ のガウス曲率の積分に等しい（ガウス・ボンネの定理 \cite{R11}）。
単位球面のガウス曲率は 1 なので
\begin{equation}
\text{上腕のひねり} = \iint_S K\,dA = \Omega\quad(\text{立体角},\ \bmod 2\pi).
\end{equation}

\paragraph{例}
\begin{itemize}
\item 下 → 前 → 横 → 下：3 つの直角を持つ球面三角形で、全球（$4\pi$）の 1/8。$\Omega=\pi/2$、すなわち 90°（古典的なコッドマンのパラドックス \cite{R34,R35}）。
\item 半頂角 $\beta$ の円錐：$\Omega = 2\pi(1-\cos\beta)$。$\beta=20^\circ$ で 21.71°、$45^\circ$ で 105.44°。
\end{itemize}
これは AMBAC の式~\eqref{eq:stokes} で、形状空間を単位球面、曲率を $K=1$ とした場合にあたる。

\section{スラスタ・RW・AMBAC の配分の線形計画}\label{app:lp}

旋回時間 $T$ を $N$ 区間（幅 $\Delta t$）に分け、各区間で一定のスラスタトルク $\tau_k=\tau_k^+-\tau_k^-$、AMBAC の周回周波数 $f_k$、RW のトルク $w_k$ を選ぶ。
平均化モデル $\dot{\bar\theta} = (L-h)/\bar A + f\Delta\theta_c$、$\dot L=\tau$、$\dot h = w$ を厳密に離散化すると、
終端の姿勢 $\bar\theta_N$、角運動量 $L_N$、ホイールの角運動量 $h_N$ はすべて決定変数の 1 次式になる。問題は
\begin{align}
\min\ &\sum_k \frac{\Delta t}{r\,I_{sp}\,g_0}\left(\tau_k^++\tau_k^-\right) \quad(\text{推進剤の質量})\\
\text{s.t. }& \bar\theta_N = 90^\circ,\quad L_N=0,\quad h_N=0,\quad |h_k|\le h_{\max},\\
& 0\le\tau_k^\pm\le\tau_{\max},\quad |f_k|\le f_{\max},\quad |w_k|\le w_{\max}
\end{align}
の線形計画で、HiGHS \cite{R25} で厳密に解ける（$r$：スラスタのモーメントアーム、$I_{sp}$：比推力、$g_0$：標準重力）。
MPC では、現在の状態から残り時間について同じ問題を 2 秒ごとに解き直す。

\section{段差の乗り越えの静力学}\label{app:steps}

\paragraph{戦車型}
車体が後ろの接地点 R と段の角 E の 2 点で支えられ、傾き $\vartheta$（$\sin\vartheta = h/s_e$、$s_e$ は R から E までの車体底面に沿った距離）にあるとする。
未知数は R の垂直力 $N_1$・水平力 $T_1$、E の車体底面に垂直な力 $N_2$・沿う力 $T_2$ である。
水平・垂直方向の力と R まわりのモーメントの 3 つのつり合い式に、摩擦の条件 $|T_i|\le\mu N_i$、$N_i\ge0$ を課し、
実行可能な最小の $\mu$（必要摩擦係数）を線形計画と二分法で求める。重心が R より後ろに来ると $N_2<0$ が必要になり、後ろに転倒する。

\paragraph{二本脚}
大腿と下腿の長さが等しく $\ell$、段の上の足首の真上に股関節があり、両者の距離を $H-h$（$H$ は直立時の股関節の高さ、$h$ は段差）とすると、
膝は前方に $d = \sqrt{\ell^2-(H-h)^2/4}$ はみ出す。足首から見た重心の前方位置を $x_c$ とすると、膝より上の質量 $m$ を支える膝トルクは
\begin{equation}
\tau_{\text{膝}} = m g\,(d - x_c),
\end{equation}
上体を前に傾けて $x_c$ を足裏の範囲内で前に移すと膝の負担が減り、その分を股関節が担う（$\tau_{\text{股}}\approx m g\,x_c$）。
相似拡大では $\tau\propto s^3\cdot s = s^4$、使えるトルクは $s^3$ なので、大きいほど上れる段差（大きさとの比）が小さくなる。

\chapter{数値計算の方法と検証}\label{app:num}

\section{数値計算の方法}

\begin{longtable}{p{0.16\linewidth}p{0.77\linewidth}}
\toprule
対象 & 方法 \\ \midrule \endhead
運動学（$L=0$） & 式~\eqref{eq:basic} を常微分方程式として積分。Sim01・02 は適応刻みの Runge--Kutta 法（相対誤差 $10^{-10}$）、Sim03 は 4 次 Runge--Kutta 法（刻み 2 ms）、Sim05A は中点法（刻み 40 ms、接続を格子上に事前計算して双一次補間）。 \\
接続 $\Gamma$ と曲率 $F$ & 角運動量が速度に線形であることを使い、単位速度を入れて $A$, $B$ を求める。$F$ は格子上の中心差分、または中心差分で各点を計算。 \\
順動力学（Sim04） & 重心系の質量行列 $M(q)$ を各剛体のヤコビアンから組み立て、$\partial M/\partial q$ は中心差分（刻み $10^{-6}$）。積分は適応刻みの Runge--Kutta 法。 \\
3 次元（Sim07・09） & 腕・脚を 40 個の点質量で表し、$\bm\omega=-A^{-1}B\dot q$ を中点則で積分、姿勢は回転ベクトルの指数写像で更新。 \\
最適化 & 軌道最適化は SLSQP \cite{R26}（多点から開始し、制約を満たす最良解を採用）。推進剤配分と段差の摩擦は線形計画（HiGHS \cite{R25}）。 \\
学習（Sim05） & OpenAI-ES \cite{R19}：対称サンプリング、順位変換した適応度、Adam \cite{R20} による更新。方策は tanh を活性化関数とする全結合ネットワーク。母集団全体を行列演算で同時に評価。 \\
地上の反力（Sim06・07・09） & 関節経路を時間 $\tau\in[0,1]$ の等間隔点で与え、位置を数値的に 2 回微分して加速度を求め、式~\eqref{eq:zmp} で ZMP を計算。動作時間 $T$ に対しては加速度が $1/T^2$ で変わることを使う。 \\
スティックスリップ（Sim06） & 静止摩擦・動摩擦を区別する Karnopp 型のモデル \cite{R30}、刻み 0.2 ms の陽解法。 \\
\bottomrule
\end{longtable}

\section{検証}

モデルと数値計算の正しさを、独立した方法どうしの一致で確認した。

\begin{longtable}{p{0.40\linewidth}p{0.53\linewidth}}
\toprule
検証の内容 & 結果 \\ \midrule \endhead
腕を往復させた後の胴体の姿勢（Sim01） & 残留 $5\times10^{-9}$°（可逆性） \\
人間大と 18 m 級（長さ 10 倍・質量 1000 倍）の回転角（Sim01） & 両方とも $-16.0432$°（スケール不変性） \\
ループ 1 周の回転：数値積分とストークスの定理（Sim02） & $+2.9778$° で一致 \\
トルクだけを入力した順動力学と運動学（Sim04） & 2 周後の回転 $+13.4086$° で一致、関節軌道の誤差 $2.6\times10^{-8}$° \\
順動力学の角運動量・エネルギー収支（Sim04） & $|L|\le 2\times10^{-9}$ kg\,m$^2$/s、$|\Delta T - \int P\,dt| \le 9\times10^{-4}$ J \\
スケーリング則の指数（Sim04） & トルク $s^{5.00}$、$f_{\max}$ $s^{-1.00}$、容量 $s^{4.00}$（理論どおり） \\
平均化モデルと位相モデル（Sim05B） & 60 s 後の周回平均の姿勢 89.63° 対 90.00° \\
格子補間と厳密な接続（Sim05A） & 最大差 0.32°（学習方策の軌道を再生して確認） \\
フルード則（Sim06） & $T_{\min}/\sqrt{s\,g_{\text{地球}}/g}$ が 12 条件で 1.616 \\
3 次元モデルと平面モデル（Sim07） & 腕を平面内で動かした場合に 5 桁一致 \\
円錐運動の回転の指数（Sim07） & 半頂角の 2.00 乗（2 次効果） \\
コッドマンのパラドックス（Sim08） & ひねりと立体角が全例で小数第 2 位まで一致 \\
\bottomrule
\end{longtable}

\section{修正の記録}

\begin{itemize}
\item \textbf{Sim06・Sim09 の関節制約}：腕の軌道最適化で関節速度・加速度の制約を間引いて課す際、配列の取り方の誤りにより、肘の関節に制約がかかっていなかった。
  修正後、Sim06 の地球での最短時間は 1.00 s（38\% 短縮）から 1.03 s（36\% 短縮）に変わった。月面の結果と Sim09 の結論は変わらない。
\item \textbf{Sim05B の揺れ補償 MPC}：当初は目標を「現在の位相での揺れ」で補正していたが、必要なのは「旋回終了時の位相での揺れ」であった。計画から終了時の位相を予測するよう改めた。
\item \textbf{Sim08 の腕組みの判定}：人の可動域の内旋 70° は腕を 90° 外転した姿勢での標準値であり、ひねりの値は測り方で変わる（第~\ref{ch:shoulder}章）。
  そのため、合否の判定に加えて「最低限必要な内旋量」を求める分析を行った。
\item \textbf{Sim09 の比較の公平性}：二本脚にだけ上体の傾きを許していたため、戦車型にも腰の関節を与え、上体の速度・加速度を腕より厳しく制限した。
\item \textbf{Sim10 の二本脚}：当初は膝だけで体を持ち上げる悲観的なモデルだったため、上体の前傾による重心移動と股関節トルクを加えた。
\end{itemize}

\section{結果の限界}

\begin{itemize}
\item 平面モデルは人の体を矢状面に投影した粗いモデルで、左右の手足の区別や体のねじれは扱っていない。
\item 関節トルクの上限・可動域は目安の値であり、個人差・機体差は扱っていない。
\item 地上の評価は準静的（台車が静止した状態での転倒・滑りの開始条件）で、転倒が始まった後の運動は扱っていない。
\item 学習の結果は乱数の種と超パラメータに依存しうる。主な結論（構造の有無による差）は再実行でも同じであった。
\end{itemize}

\chapter{パラメータと再現手順}\label{app:param}

\section{機体モデルのパラメータ}

\begin{longtable}{p{0.20\linewidth}p{0.73\linewidth}}
\toprule
モデル & パラメータ \\ \midrule \endhead
人間大・平面（Sim01〜05） & 胴体 60 kg（0.35×1.0 m の一様な長方形）。腕 4 kg・0.7 m、肩 (0, 0.35) m。脚 12 kg・0.9 m、股関節 (0, $-0.15$) m。可動域：腕 $-150$°〜150°、脚 $-120$°〜30°（$-90$° が真下、0° が前方）。 \\
リアクションホイール（Sim03） & 2 kg、半径 0.12 m の一様な円盤、最大回転数 600 rad/s。外乱 0.01 N\,m、姿勢ゲイン 0.5 s$^{-1}$、手足の周回周波数の上限 1 Hz。 \\
関節トルク（Sim04） & 持続トルクの目安：肩 40 N\,m、股関節 120 N\,m。相似拡大ではトルク上限 $\propto s^3$。 \\
スラスタ（Sim04・05） & モーメントアーム $0.5s$ m（18 m 級で 5 m）、比推力 300 s、最大トルク $2.5\times10^5$ N\,m（Sim05B）。 \\
18 m 級の配分（Sim05B） & $s=10$、$f_{\max}=0.102$ Hz、RW：$h_{\max}=8.64\times10^4$ kg\,m$^2$/s、トルク上限 200 N\,m。旋回時間 60 s、再計画 2 s ごと、終盤 10 s は AMBAC 停止（揺れ補償 MPC）。プラントの誤差：$\kappa\in[0.7,1.3]$、$\alpha\in[0.9,1.1]$、$\tau_d\in[-500,500]$ N\,m。 \\
学習（Sim05A） & 刻み 40 ms、25 s、目標 8 通り。ES：対称ペア 40、ノイズ 0.08（tg は 0.1）、学習率 0.02（tg は 0.03）、250 回。関節速度の上限 $2\pi\times1.02\times R$。 \\
学習（Sim05B） & ネットワーク 6-32-32-2、対称ペア 32、ノイズ 0.05、学習率 0.02、800 回、プラント 24 例/回。 \\
戦車型・平面（Sim06） & 台車＋胴体 80 kg（重心高さ 0.35 m）、履帯接地長 0.8 m。肩の高さ 0.75 m、上腕 4 kg・0.5 m、前腕 3 kg・0.5 m、荷物 10 kg。摩擦係数 0.6、ZMP の安全余裕 5 cm、関節速度 6 rad/s・加速度 60 rad/s$^2$。ヨー：接地面 0.8×0.6 m、腕 7 kg・1.0 m＋荷物 10 kg、振り 90°、戻り 3 s。 \\
人間大・3 次元（Sim07） & 胴体 60 kg（奥行 0.25×幅 0.4×高さ 1.0 m）、右肩 (0, $-0.2$, 0.35) m、腕 4 kg・0.7 m。 \\
戦車型・3 次元（Sim07） & 台車 80 kg（重心 (0, 0, 0.35) m）、接地面 0.8×0.6 m、肩 (0, 0, 0.75) m、腕 7 kg・1.0 m、荷物 10 kg。関節速度 6 rad/s・加速度 40 rad/s$^2$、仰角 $\ge -35$°。 \\
肩（Sim08） & 胸郭（楕円柱）の半奥行 0.09 / 0.11 / 0.13 m・半幅 0.14 m。胸鎖関節 (0.06, $-0.02$, 0.17) m、肩関節 (0, $-0.19$, 0.15) m。上腕 0.30 m（半径 5 cm）、前腕＋手 0.27 m（半径 4 cm）。可動域：水平内転 45°、内旋 70°、外旋 90°、肘 145°、前方突出 25°。目標 (0.20, 0.15, 0.00) m、許容誤差 1.5 cm。 \\
脚の比較（Sim09） & 戦車型：下部 50 kg（重心 0.2 m）、上体 30 kg（腰 0.4 m）、接地 $\pm0.4$ m。二本脚：脚 24 kg（重心 0.5 m）、上体 56 kg（股関節 0.95 m、重心は股関節から 0.25 m）、足裏 $\pm0.125$ m。上体の傾き $\pm40$°、速度 3 rad/s、加速度 15 rad/s$^2$、垂直反力 $\ge$ 自重の 30\%。3 次元：胴体 46 kg、腕 4 kg・0.7 m、脚 12 kg・0.9 m。 \\
段差（Sim10） & 戦車型：履帯接地長 0.8 m、前端の高さ 0.25 m、摩擦係数 0.6。二本脚：大腿・下腿 各 0.45 m、股関節の高さ 0.9 m、総質量 97 kg、膝・股関節トルク上限 各 250 N\,m、重心を前に移せる範囲 0.15 m、足を置く位置 前方 0.30 m・すき間 5 cm、股関節屈曲 120°、膝屈曲 140°。 \\
\bottomrule
\end{longtable}

\section{再現手順}

ソースコード・データ・図はすべて次のリポジトリにある。

\begin{center}\url{\repo}\end{center}

\begin{longtable}{p{0.30\linewidth}p{0.63\linewidth}}
\toprule
場所 & 内容 \\ \midrule \endhead
\texttt{src/ambac/} & モデル（平面・3 次元・地上・肩、動力学、学習器） \\
\texttt{sims/} & シミュレーション Sim01〜Sim10 \\
\texttt{data/}, \texttt{figures/} & 出力（数値・図）。\texttt{data/sim*\_summary.txt} に本文の数値がある \\
\texttt{docs/} & 各章のノート（Markdown）と参考文献一覧 \\
\texttt{report/} & 本報告書の \LaTeX{} ソース \\
\bottomrule
\end{longtable}

\noindent 実行手順（Python 3、NumPy・SciPy・Matplotlib \cite{R39,R40,R41}）：
\begin{Verbatim}[fontsize=\small,frame=single]
pip install -r requirements.txt
cd sims
python sim01_single_arm.py
python sim02_cyclic_gait.py
python sim03_rw_vs_ambac.py        # 約 30 秒（Sim04, 05 が使う最適ループを保存）
python sim04_dynamics_scaling.py
python sim05a_learn_ambac.py       # 学習あり、約 5 分
python sim05b_allocation.py        # 学習済み方策を再利用（再学習は RETRAIN=1、約 12 分）
python sim06_ground_tank.py
python sim07_spatial.py
python sim08_shoulder.py
python sim09_legs.py
python sim10_steps.py              # 約 3 分
\end{Verbatim}

\noindent 報告書の組版（TeX Live の LuaLaTeX、jlreq クラス）：
\begin{Verbatim}[fontsize=\small,frame=single]
cd report
latexmk -lualatex main.tex
\end{Verbatim}

\Detail{glossary}

\begin{thebibliography}{R44}
\small
\item[] \textbf{A. 出発点となった作品・設定・報道}
\bibitem[R1]{R1} 『機動戦士ガンダム』（テレビアニメ, 1979 年, 日本サンライズ）。AMBAC を備えた人型機動兵器（モビルスーツ）の作品世界。※ AMBAC という用語・概念は劇中ではなく、設定考証 \cite{R42} で示されたものとされる。
\bibitem[R42]{R42} 『ガンダムセンチュリー』（ムック, みのり書房, 1981 年）。AMBAC（Active Mass Balance Auto Control）の概念が SF 考証として最初に示された資料とされる。スタジオぬえのメンバーらが考証に関わったとされる。みのり書房は 1990 年代に活動を終えており、2000 年に樹想社から復刻版『GUNDAM CENTURY RENEWAL VERSION 宇宙翔ける戦士達』（ISBN 4-87777-028-3）が出版されている。
\bibitem[R2]{R2} 『機動戦士ガンダム』第 42 話「宇宙要塞ア・バオア・クー」。整備兵がシャアに、脚のない機体ジオングについて言う台詞「あんなの飾りです。偉い人にはそれがわからんのですよ」。これが広まる過程で「足（脚）なんて飾りです」の形で知られるようになった。
\bibitem[R3]{R3} 『アーマード・コア』シリーズ（フロム・ソフトウェア, 1997 年〜）。パーツを組み合わせて機体を構成する「アセンブル」の例として参照。
\bibitem[R43]{R43} オリコンニュースの記事. \url{https://www.oricon.co.jp/news/2385521/full/} 大西卓哉さんの X への投稿 \cite{R44} を紹介し、宇宙での姿勢制御を「ガンダムの AMBAC」と表現したことがガンダムファンの間で話題になったことを伝える記事（角運動量保存則などの物理の解説記事ではない）。※ 記事名は未照合。
\bibitem[R44]{R44} 大西卓哉（@Astro\_Onishi）の X への投稿, 2025 年 5 月 18 日. \url{https://x.com/Astro_Onishi/status/1924053663238046048} （国際宇宙ステーションで脚を使って姿勢を制御し、AMBAC に言及したもの）※ 本リポジトリの作業環境からは閲覧できず、本文は未照合。
\item[] \textbf{B. 幾何力学・幾何学的位相・非ホロノミック系}
\bibitem[R4]{R4} R. Montgomery, "Gauge theory of the falling cat," in M. J. Enos (ed.), \textit{Dynamics and Control of Mechanical Systems: The Falling Cat and Related Problems}, Fields Institute Communications, vol. 1, American Mathematical Society, 1993.
\bibitem[R5]{R5} J. E. Marsden, \textit{Lectures on Mechanics}, London Mathematical Society Lecture Note Series 174, Cambridge University Press, 1992.
\bibitem[R6]{R6} A. Shapere and F. Wilczek (eds.), \textit{Geometric Phases in Physics}, World Scientific, 1989.
\bibitem[R7]{R7} T. R. Kane and M. P. Scher, "A dynamical explanation of the falling cat phenomenon," \textit{International Journal of Solids and Structures}, 5(7), 663–670, 1969.
\bibitem[R8]{R8} T. R. Kane and M. P. Scher, "Human self-rotation by means of limb movements," \textit{Journal of Biomechanics}, 3(1), 39–49, 1970. doi:10.1016/0021-9290(70)90049-7 $\\checkmark$
\bibitem[R9]{R9} R. W. Brockett, "Asymptotic stability and feedback stabilization," in R. W. Brockett, R. S. Millman, H. J. Sussmann (eds.), \textit{Differential Geometric Control Theory}, Birkhäuser, pp. 181–191, 1983.
\bibitem[R10]{R10} R. M. Murray and S. S. Sastry, "Nonholonomic motion planning: steering using sinusoids," \textit{IEEE Transactions on Automatic Control}, 38(5), 700–716, 1993.
\bibitem[R11]{R11} M. P. do Carmo, \textit{Differential Geometry of Curves and Surfaces}, Prentice-Hall, 1976.（ガウス・ボンネの定理）
\item[] \textbf{C. 宇宙ロボット・人工衛星の姿勢制御}
\bibitem[R12]{R12} Y. Nakamura and R. Mukherjee, "Nonholonomic path planning of space robots via a bidirectional approach," \textit{IEEE Transactions on Robotics and Automation}, 7(4), 500–514, 1991. doi:10.1109/70.86080 $\\checkmark$
\bibitem[R13]{R13} E. Papadopoulos and S. Dubowsky, "On the nature of control algorithms for free-floating space manipulators," \textit{IEEE Transactions on Robotics and Automation}, 7(6), 750–758, 1991.
\bibitem[R14]{R14} Y. Umetani and K. Yoshida, "Resolved motion rate control of space manipulators with generalized Jacobian matrix," \textit{IEEE Transactions on Robotics and Automation}, 5(3), 303–314, 1989.
\bibitem[R15]{R15} D. N. Nenchev, K. Yoshida, P. Vichitkulsawat, M. Uchiyama, "Reaction null-space control of flexible structure mounted manipulator systems," \textit{IEEE Transactions on Robotics and Automation}, 15(6), 1011–1023, 1999.
\bibitem[R16]{R16} B. Wie, \textit{Space Vehicle Dynamics and Control}, 2nd ed., AIAA, 2008.（リアクションホイール・CMG・スラスタによる姿勢制御）
\bibitem[R17]{R17} M. J. Sidi, \textit{Spacecraft Dynamics and Control: A Practical Engineering Approach}, Cambridge University Press, 1997.
\bibitem[R18]{R18} M. D. Shuster, "A survey of attitude representations," \textit{Journal of the Astronautical Sciences}, 41(4), 439–517, 1993.（回転の表現と特異点）
\item[] \textbf{D. 学習・最適化}
\bibitem[R19]{R19} T. Salimans, J. Ho, X. Chen, S. Sidor, I. Sutskever, "Evolution strategies as a scalable alternative to reinforcement learning," arXiv:1703.03864, 2017.
\bibitem[R20]{R20} D. P. Kingma and J. Ba, "Adam: A method for stochastic optimization," \textit{ICLR}, 2015. arXiv:1412.6980.
\bibitem[R21]{R21} A. Iscen, K. Caluwaerts, J. Tan, T. Zhang, E. Coumans, V. Sindhwani, V. Vanhoucke, "Policies modulating trajectory generators," \textit{Proceedings of the 2nd Conference on Robot Learning (CoRL)}, PMLR 87, 916–926, 2018. $\\checkmark$
\bibitem[R22]{R22} A. J. Ijspeert, "Central pattern generators for locomotion control in animals and robots: A review," \textit{Neural Networks}, 21(4), 642–653, 2008.
\bibitem[R23]{R23} J. Tobin, R. Fong, A. Ray, J. Schneider, W. Zaremba, P. Abbeel, "Domain randomization for transferring deep neural networks from simulation to the real world," \textit{IEEE/RSJ IROS}, 2017.
\bibitem[R24]{R24} X. B. Peng, P. Abbeel, S. Levine, M. van de Panne, "DeepMimic: Example-guided deep reinforcement learning of physics-based character skills," \textit{ACM Transactions on Graphics}, 37(4), 2018.（モーションキャプチャーを手本とする物理ベースの強化学習）
\bibitem[R25]{R25} Q. Huangfu and J. A. J. Hall, "Parallelizing the dual revised simplex method," \textit{Mathematical Programming Computation}, 10(1), 119–142, 2018.（HiGHS: scipy の linprog で使用）
\bibitem[R26]{R26} D. Kraft, "A software package for sequential quadratic programming," Technical Report DFVLR-FB 88-28, DLR, 1988.（SLSQP: scipy の minimize で使用）
\bibitem[R27]{R27} T. Flash and N. Hogan, "The coordination of arm movements: an experimentally confirmed mathematical model," \textit{Journal of Neuroscience}, 5(7), 1688–1703, 1985.（最小躍度軌道）
\item[] \textbf{E. 地上での安定性・移動}
\bibitem[R28]{R28} M. Vukobratović and B. Borovac, "Zero-moment point — thirty five years of its life," \textit{International Journal of Humanoid Robotics}, 1(1), 157–173, 2004.
\bibitem[R29]{R29} R. McN. Alexander and A. S. Jayes, "A dynamic similarity hypothesis for the gaits of quadrupedal mammals," \textit{Journal of Zoology}, 201(1), 135–152, 1983.（フルード数による相似則）
\bibitem[R30]{R30} D. Karnopp, "Computer simulation of stick-slip friction in mechanical dynamic systems," \textit{Journal of Dynamic Systems, Measurement, and Control}, 107(1), 100–103, 1985.
\bibitem[R31]{R31} F. B. Horak and L. M. Nashner, "Central programming of postural movements: adaptation to altered support-surface configurations," \textit{Journal of Neurophysiology}, 55(6), 1369–1381, 1986.（足首の戦略・腰の戦略）$\\checkmark$
\bibitem[R32]{R32} J. Y. Wong, \textit{Theory of Ground Vehicles}, 4th ed., Wiley, 2008.（履帯車両の走行・段差乗り越え）
\bibitem[R33]{R33} G. Galilei, \textit{Discorsi e dimostrazioni matematiche intorno a due nuove scienze}（新科学対話）, 1638.（二乗三乗則）
\item[] \textbf{F. 肩・人体工学}
\bibitem[R34]{R34} E. A. Codman, \textit{The Shoulder}, Thomas Todd, Boston, 1934.（コッドマンのパラドックス）
\bibitem[R35]{R35} J. C. Politti, G. Goroso, M. E. Valentinuzzi, O. Bravo, "Codman's paradox of the arm rotations is not a paradox: mathematical validation," \textit{Medical Engineering \& Physics}, 20(4), 257–260, 1998. $\\checkmark$（その後も解釈をめぐる議論がある）
\bibitem[R36]{R36} G. Wu et al., "ISB recommendation on definitions of joint coordinate systems of various joints for the reporting of human joint motion — Part II: shoulder, elbow, wrist and hand," \textit{Journal of Biomechanics}, 38(5), 981–992, 2005.（肩の角度の測り方の標準）
\bibitem[R37]{R37} V. T. Inman, J. B. Saunders, L. C. Abbott, "Observations on the function of the shoulder joint," \textit{Journal of Bone and Joint Surgery}, 26(1), 1–30, 1944.（肩甲上腕リズム）
\bibitem[R38]{R38} American Academy of Orthopaedic Surgeons, \textit{Joint Motion: Method of Measuring and Recording}, AAOS, 1965.（関節可動域の標準値）
\item[] \textbf{G. ソフトウェア}
\bibitem[R39]{R39} C. R. Harris et al., "Array programming with NumPy," \textit{Nature}, 585, 357–362, 2020.
\bibitem[R40]{R40} P. Virtanen et al., "SciPy 1.0: fundamental algorithms for scientific computing in Python," \textit{Nature Methods}, 17, 261–272, 2020.
\bibitem[R41]{R41} J. D. Hunter, "Matplotlib: A 2D graphics environment," \textit{Computing in Science \& Engineering}, 9(3), 90–95, 2007.
\item[] \textbf{H. 関連する事例}
\bibitem[R45]{R45} T. Libby, T. Y. Moore, E. Chang-Siu, D. Li, D. J. Cohen, A. Jusufi, R. J. Full, "Tail-assisted pitch control in lizards, robots and dinosaurs," \textit{Nature}, 481, 181–184, 2012. doi:10.1038/nature10710 $\\checkmark$（尻尾を振って空中で姿勢を制御する）
\bibitem[R46]{R46} M. R. Yeadon, "The biomechanics of twisting somersaults. Part III: aerial twist," \textit{Journal of Sports Sciences}, 11, 209–218, 1993. $\\checkmark$（空中での腕などの非対称な動きでひねりを生む）
\bibitem[R47]{R47} Boston Dynamics による油圧式 Atlas の引退と電動式 Atlas の発表（2024 年 4 月）. 報道の例: \url{https://builtin.com/articles/boston-dynamics-electric-atlas-20240417} $\\checkmark$
\end{thebibliography}

\end{document}
